% NetHack 5.0  Guidebook.tex $NHDT-Date: 1782274458 2026/06/23 23:14:18 $  $NHDT-Branch: NetHack-5.0 $:$NHDT-Revision: 1.630 $ */
\documentclass[titlepage]{article}
\usepackage[utf8]{inputenc}
\usepackage[T1]{fontenc}
\usepackage[french]{babel}
\usepackage[hidelinks]{hyperref}
\usepackage{longtable}
\usepackage{enumitem}
\usepackage[a4paper, text={160mm, 220mm}, centering]{geometry}

\setlist[description]{leftmargin=30mm, topsep=2mm, partopsep=0mm, parsep=0mm, itemsep=1mm, labelwidth=28mm, labelsep=2mm}

\hyphenation{CRASHREPORTURL}

\begin{document}
%
% input file: guidebook.mn
%
%.ds h0 "
%.ds h1 %.ds h2 \%
%.ds f0 "

%.mt
\title{\LARGE Un guide des Mazes of Menace :\\
\Large Guide de l'exploration du donjon de \textit{NetHack}}

%.au
\author{Version originale - Eric S. Raymond\\
(Révisée et augmentée pour 5.0.1 par Mike Stephenson et d'autres)\\[2ex]
Traduction française du \textit{Guidebook for NetHack} d'origine}
%DO NOT REMOVE NH_DATESUB \date{Date(%B %-d, %Y)}
\date{4 juillet 2026}

\maketitle

%.pg
%.hn 1
\section{Introduction}

%.pg

Depuis quelque temps, vous vous sentez insatisfait et distant
dans vos occupations quotidiennes.  D'étranges rêves de prospection, de vol,
de croisade et de combat hantent votre sommeil depuis de nombreux mois,
sans que vous sachiez trop pourquoi.  Vous vous demandez si, en fait,
vous n'avez pas fait ces rêves toute votre vie, en parvenant on ne sait
comment à les oublier jusqu'à présent.  Certaines nuits, vous vous réveillez
en sursaut et poussez un cri, terrifié par le souvenir saisissant des
créatures étranges et puissantes qui semblent se tapir à chaque recoin du
donjon de vos rêves.  Ces détails qui hantent vos rêves pourraient-ils être réels ?
Nuit après nuit, vous sentez grandir le désir de pénétrer dans les mystérieuses
cavernes situées près des ruines.  Chaque matin, pourtant, vous chassez
bien vite cette idée de votre esprit en vous rappelant les récits de ceux qui
sont entrés dans les cavernes avant vous et ne sont jamais revenus.  Un jour,
vous ne pouvez plus résister à l'envie d'aller chercher le lieu fantastique
de vos rêves.
Après tout, quand d'autres aventuriers repassaient par ici après un séjour
dans les cavernes, ils semblaient généralement en meilleure posture qu'à
leur premier passage.  Et qui pouvait affirmer que tous ceux qui n'étaient
pas revenus n'avaient pas tout simplement poursuivi leur route ?
%.pg

En vous renseignant, vous entendez parler d'une babiole, appelée par certains
l'Amulette de Yendor, qui vous apportera une grande richesse si vous parvenez
à la trouver.  Une légende qu'on vous a contée affirmait même que celui qui
trouverait l'amulette se verrait accorder l'immortalité par les dieux.
L'amulette se trouverait quelque part au-delà de la Vallée de Gehennom, tout
au fond des Mazes of Menace.  En entendant ces légendes, vous comprenez
aussitôt qu'il existe une raison profonde, encore inconnue, pour laquelle
vous devez descendre dans les cavernes et y chercher cette amulette dont on
parle.  Même si les rumeurs sur les pouvoirs de l'amulette sont fausses, vous
vous dites que vous devriez au moins pouvoir vendre le récit de vos aventures
aux ménestrels du coin pour une coquette somme, surtout si vous croisez en
chemin quelques-unes des créatures terrifiantes et magiques de vos rêves.
Vous passez une dernière nuit à reprendre des forces à l'auberge du coin,
de plus en plus déprimé à mesure que vous voyez baisser, sur les murs de
l'auberge, la cote de vos chances de succès.

%.pg
\noindent Au matin, vous vous réveillez, rassemblez vos affaires et
partez pour le donjon.  Après plusieurs jours d'un voyage
sans histoire, vous apercevez les ruines antiques qui marquent l'entrée des
Mazes of Menace.  Comme il fait déjà nuit noire, vous campez devant l'entrée
et dormez à la belle étoile.  Au matin, vous rassemblez votre équipement,
prenez ce qui sera peut-être votre dernier repas à l'air libre, et entrez
dans le donjon\ldots

%.hn 1
\section{Que se passe-t-il ici ?}

%.pg
Vous venez de commencer une partie de \textit{NetHack}.  Votre but est de ramasser
autant de trésors que possible, de récupérer l'Amulette de Yendor et de
sortir vivant des Mazes of Menace.

%.pg
Vos capacités et vos atouts face aux périls de l'aventure
varieront selon votre origine et votre formation :

%.pg
%
\begin{description}
\item[Archéologues]
connaissent plutôt bien les donjons, ce qui leur
permet de se déplacer vite et de surprendre les sales bêtes du coin.  Ils
commencent équipés des outils nécessaires à une véritable expédition
scientifique et savent lire les langues anciennes.
%.pg
%
\item[Barbares]
sont des guerriers venus de l'arrière-pays, endurcis au combat.
Ils entament leur quête sans rien d'autre qu'une force peu commune, un fidèle
haubert et une grande épée à deux mains.
%.pg
%
\item[Hommes des cavernes \textrm{\textmd{et}} Femmes des cavernes]
commencent avec une force exceptionnelle mais, malheureusement, avec des armes néolithiques.
%.pg
%
\item[Guérisseurs]
sont versés en médecine et en apothicairerie.  Ils
connaissent les herbes et les simples qui rendent la vitalité, apaisent la
douleur, anesthésient et neutralisent
les poisons ; et grâce à leurs instruments, ils savent deviner l'état de santé
ou de maladie d'un être.  Leur pratique médicale leur rapporte des sommes tout
à fait raisonnables, avec lesquelles ils entrent dans le donjon.
%.pg
%
\item[Chevaliers]
se distinguent du vulgaire escarmoucheur par leur
dévouement aux idéaux de la chevalerie et par l'excellence incomparable
de leur armure.
%.pg
%
\item[Moines]
sont des ascètes qui, par la pratique rigoureuse de
disciplines physiques et mentales, sont devenus capables de se battre aussi
efficacement sans armes qu'avec.  Ils ne portent pas d'armure mais compensent
par une mobilité accrue.
%.pg
%
\item[Prêtres \textrm{\textmd{et}} Prêtresses]
sont des clercs militants, des croisés
qui font avancer la cause de la vertu par les armes, l'armure et les arts
thaumaturgiques.  Leur capacité à communier avec les divinités par la prière
les tire parfois d'un mauvais pas, mais peut aussi les y plonger.
%.pg
%
\item[Rôdeurs]
sont surtout à l'aise dans les bois et, selon certains,
un peu déplacés dans un donjon.  Ce sont pourtant des experts du tir à l'arc,
du pistage et du déplacement furtif.
%.pg
%
\item[Voleurs]
sont des larrons agiles et discrets, qui s'y connaissent en
serrures, en pièges et en poisons.  Leur atout est la surprise, dont ils
tirent le meilleur parti.
%.pg
%
\item[Samouraïs]
sont les guerriers d'élite du Nippon féodal.  Légèrement
armurés et rapides, ils portent le %
\textit{dai-sho}, deux sabres au tranchant
redoutable.
%.pg
%
\item[Touristes]
commencent avec beaucoup d'or (de quoi faire des achats),
une carte de crédit, beaucoup de nourriture, quelques cartes et un appareil
photo coûteux.  La plupart des monstres n'aiment pas être photographiés.
%.pg
%
\item[Valkyries]
sont de robustes guerrières.  Leur éducation dans les rudes
Terres du Nord les rend fortes, les endurcit aux froids extrêmes et leur
inculque discrétion et ruse.
%.pg
%
\item[Magiciens]
commencent avec des connaissances en magie, un assortiment
d'objets magiques et une affinité particulière pour les arts occultes.  Bien
qu'il paraisse faible et facile à vaincre au premier abord, un Magicien
expérimenté est un ennemi redoutable.
\end{description}

%.pg
Vous pouvez aussi choisir la race de votre personnage (dans certaines limites ;
la plupart des rôles n'acceptent que certaines races) :

%.pg
%
\begin{description}
\item[Nains]
sont plus petits que les humains ou les elfes, mais ce sont des
individus trapus et solides.  Le trait le plus remarquable des Nains est leur
grande maîtrise de l'exploitation minière et du travail des métaux.  On dit
que les armures naines ne le cèdent en qualité même pas aux armures de mithril
des Elfes.
%.pg
%
\item[Elfes]
sont agiles, rapides et perspicaces ; bien peu de choses
échappent à un Elfe.  La qualité de l'artisanat elfique leur donne souvent
un avantage en matière d'armes et d'armures.
%.pg
%
\item[Gnomes]
sont plus petits que les nains, mais leur ressemblent dans
l'ensemble.  Les Gnomes sont réputés être des mineurs experts, et l'on sait
qu'un complexe minier souterrain secret construit par cette race existe au
sein des Mazes of Menace, plein de richesses comme de dangers.
%.pg
%
\item[Humains]
sont de loin la race la plus répandue du monde de la surface,
et constituent donc la norme à laquelle on compare souvent les autres races.
Bien qu'ils n'aient aucune capacité particulière, ils peuvent réussir dans
n'importe quel rôle.
%.pg
%
\item[Orques]
sont une race cruelle et barbare qui hait tout être vivant
(y compris les autres orques).  Par-dessus tout, les Orques haïssent les Elfes
avec une passion sans égale, et feront tout pour en tuer un à la moindre occasion.
Les armes et armures façonnées par les Orques sont en général de qualité inférieure.
\end{description}

%.hn 1
\section{Que signifie tout ce qui est affiché à l'écran ?}
%.pg
L'écran présente une carte des endroits où vous êtes allé et de ce que vous
avez vu sur le niveau actuel du donjon ; à mesure que vous explorez le niveau,
il apparaît à l'écran devant vous.

%.pg
Lorsque \textit{rogue}, l'ancêtre de \textit{NetHack}, est apparu, son orientation
vers l'affichage à l'écran était presque unique parmi les jeux de fantasy
sur ordinateur.  Depuis, l'affichage à l'écran est devenu la norme plutôt
que l'exception ; \textit{NetHack} perpétue cette belle tradition.  Contrairement aux
jeux d'aventure textuels qui acceptent des commandes sous forme de phrases en
pseudo-anglais et décrivent les résultats avec des mots, les commandes de
\textit{NetHack} tiennent toutes en une ou deux touches et les résultats sont affichés
graphiquement à l'écran.  Une taille d'écran minimale de 24 lignes sur
80 colonnes est recommandée ; si l'écran est plus grand, seule une zone de
$21\times80$ sera utilisée pour la carte.

%.pg
\textit{NetHack} peut même être joué par des joueurs aveugles, à l'aide de
plages braille ou de synthétiseurs vocaux.  Les instructions pour configurer
\textit{NetHack} pour les aveugles figurent plus loin dans ce document.

%.pg
\textit{NetHack} génère un nouveau donjon à chaque partie ; même les
auteurs le trouvent toujours divertissant et passionnant, bien
qu'ils l'aient gagné plusieurs fois.

%.pg
\textit{NetHack} propose diverses options d'affichage.  Les options dont vous
disposez varient d'un portage à l'autre, selon les capacités de votre matériel et de
votre logiciel, et selon les options de compilation qui étaient activées
lors de la création de votre exécutable.  Les trois modes d'affichage
possibles sont : une interface en caractères monochrome, une interface en
caractères couleur, et une interface graphique utilisant de petites images
appelées tuiles (\textit{tiles}).  Les deux interfaces en caractères permettent de
substituer des polices comportant d'autres caractères, mais les affectations
par défaut utilisent les caractères ASCII standard pour tout représenter.
Les différents modes d'affichage ne changent rien au déroulement du jeu.
Comme nous ne pouvons reproduire ni les tuiles ni les couleurs dans ce guide,
et comme c'est commun à tous les portages, nous utiliserons les caractères
ASCII par défaut de l'affichage en caractères monochrome pour désigner ce
que vous pourriez voir à l'écran pendant votre partie.
%.pg
Pour comprendre ce qui se passe dans \textit{NetHack}, vous devez d'abord
comprendre ce que \textit{NetHack} fait de l'écran.  L'écran de \textit{NetHack}
remplace les descriptions ``Vous voyez\ldots'' des jeux d'aventure textuels.
La figure 1 donne un exemple de ce à quoi peut ressembler un écran de \textit{NetHack}.
L'aspect de l'écran chez vous dépend de votre plateforme.
%.BR 2

% (Either generated by hand or else the composite of two different
% situations.  Originally the character had only reached a second room
% (unchanged here) by turn 257 (now changed to 752) and was already
% Weak from hunger (now changed to just Hungry) and also lacked any of
% Tourist's starting gold.  Confusion is added to include a condition.)
%
% Width is forced to match similar figure in Guidebook.mn where it is
% constrained by the margins of plain text output (Guidebook.txt).
\vbox{
\begin{verbatim}
        La chauve-souris vous mord !

                ------
                |....|    ----------
                |.<..|####...@...$.|
                |....-#   |...B....+
                |....|    |.d......|
                ------    -------|--



        Player le Promeneur  Fo:12 Dx:7 Co:18 In:11 Sa:9 Ch:15 Neutre
        Niv:1 $:993 PV:9(12) Pm:3(3) CA:10 Exp:1/19 T:752 Affamé Conf
\end{verbatim}
\begin{center}
Figure 1
\end{center}
}

% 3-line status includes trailing spaces to force the width to match the
% 2-line data above; unlike Guidebook.pm, we can't add a trailing comment
% to make them visible
\vbox{
\begin{verbatim}
        Player le Promeneur  Fo:12 Dx:7 Co:18 In:11 Sa:9 Ch:15        
        Neutre $:993 PV:9(12) Pm:3(3) CA:10 Exp:1/19 Affamé           
        Niv:1 T:752                                   Conf            
\end{verbatim}
\begin{center}
Figure 2
\end{center}
}

%.hn 2
\subsection*{Les lignes d'état (en bas)}

Les deux (ou trois) dernières lignes de l'écran contiennent plusieurs
informations sibyllines décrivant votre état actuel.
La figure 1 montre la zone d'état traditionnelle sur deux lignes, sous la carte.
La figure 2 montre seulement la zone d'état, lorsque l'option \textit{statuslines:3}
a été activée (toutes les interfaces ne la prennent pas en charge).
Si une ligne d'état devient plus large que l'écran, il se peut que vous ne
la voyiez pas en entier parce qu'elle est tronquée.
Quand les nombres grandissent et que plusieurs \textit{conditions} sont présentes,
le format sur deux lignes manque de place sur la deuxième ligne, mais
\textit{statuslines:2}
est la valeur par défaut, car un terminal basique de 24 lignes n'est pas assez
haut pour la troisième ligne.

%.pg
Voici ce que signifient les différents éléments de la ligne d'état :

%.lp
\begin{description}
\item[Titre]
Le nom de votre personnage et son rang professionnel (fondé sur son rôle et
son \textit{niveau d'expérience}, voir plus bas).
%.lp
\item[Force]
Une mesure de la force de votre personnage ; l'une de vos six caractéristiques
de base.  Les caractéristiques d'un personnage humain peuvent aller de 3 à 18
inclus ; les non-humains peuvent dépasser ces limites
(il arrive que vous obteniez une force surhumaine de la forme 18/xx, et la
magie peut aussi faire dépasser les limites normales aux caractéristiques).
Plus votre force est élevée, plus vous êtes fort.  La force influe sur votre
réussite dans les tâches physiques, sur les dégâts que vous infligez au
combat et sur la quantité de butin que vous pouvez porter.
%.lp
\item[Dextérité]
La dextérité influe sur vos chances de toucher au combat, d'éviter les pièges
et d'accomplir d'autres tâches demandant de l'agilité ou de la manipulation
d'objets.
%.lp
\item[Constitution]
La constitution influe sur votre capacité à vous remettre de vos blessures et
d'autres atteintes à votre endurance.
Quand votre force est faible ou modeste, la constitution influe aussi sur la
quantité que vous pouvez porter.  Avec une force suffisamment élevée, la part
de votre constitution dans votre capacité de charge ne compte plus.
%.lp
\item[Intelligence]
L'intelligence influe sur votre capacité à lancer des sorts et à lire des
grimoires.
%.lp
\item[Sagesse]
La sagesse provient de votre expérience pratique (en particulier en matière de
magie).  Elle influe sur votre énergie magique.
%.lp
\item[Charisme]
Le charisme influe sur la façon dont certaines créatures réagissent envers vous.
En particulier, il peut influer sur les prix que vous proposent les
commerçants.
%.lp
\item[Alignement]
%
\textit{Loyal}, \textit{Neutre} ou \textit{Chaotique}.  On assimile souvent Loyal au
bien et Chaotique au mal, mais légal et moral ne coïncident pas toujours.
Votre alignement influe sur la façon dont les autres
monstres réagissent envers vous.  Les monstres de même alignement sont plus
enclins à ne pas être agressifs, tandis que ceux d'alignement opposé sont plus
enclins à être gravement offensés par votre présence.
%.lp
\item[Niveau du donjon]
La profondeur à laquelle vous vous trouvez dans le donjon.  Vous commencez au
niveau un et le nombre augmente à mesure que vous descendez dans le donjon.
Certains niveaux sont spéciaux et sont identifiés par un nom plutôt que par un
numéro.  L'Amulette de Yendor se trouverait quelque part sous le vingtième
niveau.
%.lp
\item[Or]
Le nombre de pièces d'or que vous portez ouvertement.  L'or que vous avez
caché dans des conteneurs n'est pas compté.
%.lp
\item[Points de vie]
Vos points de vie actuels et maximaux.  Les points de vie indiquent combien
de dégâts vous pouvez encaisser avant de mourir.  Plus vous êtes touché au
combat, plus ils diminuent.  Vous pouvez regagner des points de vie en vous
reposant, ou en utilisant certains objets magiques ou certains sorts.  Le
nombre entre parenthèses est le maximum que vos points de vie peuvent atteindre.
%.lp
\item[Pouvoir]
Points de magie.  Ils indiquent combien d'énergie mystique (\textit{mana})
vous avez à disposition pour lancer des sorts.  Là encore, le repos régénère
la quantité disponible.
%.lp
\item[Classe d'armure]
Une mesure de l'efficacité avec laquelle votre armure arrête les coups des
créatures hostiles.  Plus ce nombre est bas, plus l'armure est efficace ; il est
tout à fait possible d'avoir une classe d'armure négative.
Consultez la sous-section \textit{Armures} de la section \textit{Objets} pour plus de détails.
%.lp
\item[Expérience]
Votre niveau d'expérience actuel.
Si l'option \textit{showexp}
est activée, il est suivi d'une barre oblique et des points d'expérience.
Au fil de vos aventures, vous gagnez des points d'expérience.
À certains totaux de points d'expérience, vous gagnez un niveau d'expérience.
Plus vous êtes expérimenté, mieux vous combattez et résistez aux attaques
magiques.
(Lorsque votre niveau atteint deux chiffres, l'intérêt d'afficher les points
avec lui a nettement diminué.
Vous pouvez utiliser la commande `\texttt{O}' pour désactiver \textit{showexp}
et éviter d'occuper la place limitée de la ligne d'état.)
%.lp
\item[Temps]
Le nombre de tours écoulés jusqu'ici, affiché si l'option
\textit{time} est activée.
%.lp
\item[État]
Faim :
votre état de faim actuel.
Les valeurs sont \textit{Rassasié}, \textit{Pas~affamé} (ou \textit{Normal}),
\textit{Affamé}, \textit{Faible} et \textit{Défaillant}.
%.\" not mentioned: Fainted
Rien n'est affiché à l'état \textit{Normal}.

%.lp ""
Charge :
une indication de la façon dont ce que vous portez gêne vos déplacements.
Les valeurs sont \textit{Non~chargé}, \textit{Chargé}, \textit{Stressé},
\textit{Tendu}, \textit{Exténué} et \textit{Surchargé}.
Rien n'est affiché à l'état \textit{Non~chargé}.

%.lp ""
Conditions~fatales :
\textit{Pierre} (pétrification en cours),
\textit{Glu} (transformation en glu verte),
\textit{Étrang} (strangulation),
\textit{Intox} (intoxication alimentaire aiguë),
\textit{Malade} (maladie mortelle).

%.lp ""
Conditions~non~fatales :
\textit{Aveugle} (vous ne voyez rien), \textit{Sourd} (vous n'entendez rien),
\textit{Étourdi} (étourdi), \textit{Conf} (confus), \textit{Hallu} (hallucinations).

%.lp ""
Modificateurs~de~déplacement :
\textit{Lév} (lévitation), \textit{Vol} (vol), \textit{Selle} (en selle).

%.lp ""
Il existe d'autres conditions et modificateurs, mais il n'y a pas assez de
place pour les afficher avec les autres champs d'état.
\\
% unindented paragraph
La commande \texttt{\#attributes} (touche par défaut \texttt{\textasciicircum X}) affiche
toutes les informations d'état actuelles sous forme non abrégée.
Elle montre aussi d'autres informations qui pourraient figurer sur les
lignes d'état si celles-ci avaient plus de place.

\end{description}

%.hn 2
\subsection*{La ligne de messages (en haut)}

%.pg
La première ligne de l'écran est réservée aux messages décrivant
ce qu'il est impossible de représenter visuellement.  Si vous voyez
``\texttt{--Plus--}'' sur la première ligne, cela signifie que \textit{NetHack} a
un autre message à afficher à l'écran, mais veut s'assurer
que vous avez bien lu celui qui s'y trouve déjà.  Pour lire le message suivant,
appuyez simplement sur la barre d'espace.

%.pg
Pour modifier la manière dont les messages s'affichent sur la ligne de messages,
et lesquels s'affichent, consultez ``\textit{Configurer les types de messages}'' et l'option \textit{verbose}.

%.hn 2
\subsection*{La carte (le reste de l'écran)}

%.pg
Le reste de l'écran est la carte du niveau, telle que vous l'avez explorée
jusqu'ici.  Chaque symbole à l'écran représente quelque chose.  Vous pouvez régler
diverses options graphiques
pour changer certains des symboles qu'utilise le jeu ; sinon, le
jeu utilise les symboles par défaut.  Voici la liste de ce que signifient
les symboles par défaut :

\begin{description}[font=\mdseries\ttfamily]
\item[-]
Les murs horizontaux ou les coins d'une pièce, ou une porte ouverte est/ouest.
\item[|]
Les murs verticaux d'une pièce, ou une porte ouverte nord/sud, ou une tombe.
\item[.]
Le sol d'une pièce, ou de la glace, ou une embrasure sans porte, ou le tablier
d'un pont-levis abaissé.
\item[\#]
Un couloir, ou des barreaux de fer, ou un arbre, ou la herse d'un pont-levis
relevé.\\
%.lp ""
Remarque : les gravures dans les couloirs apparaissent aussi sous la forme \#, mais
dans une couleur différente de celle des cases de couloir normales.
\item[>]
Escalier descendant : un passage vers le niveau suivant.
\item[<]
Escalier montant : un passage vers le niveau précédent.
\item[+]
Une porte fermée, ou un grimoire contenant un sort que vous pourrez peut-être apprendre.
\item[@]
Votre personnage, ou un humain, ou un elfe.
\item[\textdollar]
Un tas d'or.
\item[\textasciicircum]
Un piège (une fois que vous l'avez détecté).
\item[)]
Une arme.
\item[{[}]
Une armure ou une pièce d'armure.
\item[\%]
Quelque chose de comestible (pas forcément sain).
\item[?]
Un parchemin.
\item[/]
Une baguette.
\item[=]
Un anneau.
\item[!]
Une potion.
\item[(]
Un objet utile (pioche, clé, lampe\ldots).
\item["]
Une amulette ou une toile d'araignée.
\item[*]
Une gemme ou une pierre (peut-être précieuse, peut-être sans valeur).
\item[\textasciigrave]
Un rocher, une statue ou une gravure sur le sol d'une pièce.\\
%.lp ""
Remarque : les statues sont affichées comme les monstres qu'elles représentent
et n'apparaissent donc pas sous la forme d'un \textit{accent grave} (alias \textit{back-tick}).
\item[0]
Un boulet de fer.
\item[\textunderscore]
Un autel, ou une chaîne de fer.
\item[\textbraceleft]
Une fontaine ou un évier.
\item[\textbraceright]
Un bassin d'eau, des douves ou un mur d'eau,
ou un bassin de lave ou un mur de lave.
\item[\textbackslash]
Un trône opulent.
\item[a-z \textrm{\textmd{et}}]
\item[A-HJ-Z \textrm{\textmd{et}}]
\item[@\&\textquotesingle :;]
Les lettres et certains autres symboles représentent les divers habitants
des Mazes of Menace.
Méfiez-vous, ils peuvent être méchants et vicieux.
Parfois, cependant, ils peuvent se montrer utiles.
\item[I]
Plutôt qu'un type de monstre précis, ce symbole marque la dernière position
connue d'un monstre invisible ou que vous ne voyez pas pour une autre raison.
Notez que le monstre a pu se déplacer.
Les commandes `\texttt{s}', `\texttt{F}' et `\texttt{m}' peuvent être utiles ici.
\item[1-5]
Les chiffres 1 à 5 peuvent être affichés pour signaler des monstres invisibles
détectés grâce à la propriété d'\textit{Avertissement}.
Les monstres les moins dangereux sont indiqués par les valeurs les plus basses,
les plus dangereux par les valeurs les plus hautes.

\end{description}
%.pg
Inutile de mémoriser tous ces symboles ; vous pouvez demander au jeu ce que
représente n'importe quel symbole avec la commande `\texttt{/}' (voir la section
suivante pour plus de détails).

%.hn 1
\section{Commandes}

%.pg
Une commande se lance en tapant le ou les deux caractères auxquels
elle est associée, ou en tapant son nom dans la saisie des commandes
étendues.  Certaines commandes,
comme ``\texttt{search}'' (chercher), ne nécessitent aucune information supplémentaire
de la part de \textit{NetHack}.  D'autres peuvent demander des précisions, par
exemple une direction ou un objet à utiliser.  Pour les commandes qui
exigent des informations supplémentaires, \textit{NetHack} vous présentera soit
un menu de choix, soit une invite en ligne de commande demandant ces informations.
Ce qui vous sera présenté dépend principalement de la façon dont vous avez
réglé l'option `\textit{menustyle}'.

%.pg
Par exemple, une question courante, de la forme ``\texttt{Que voulez-vous
utiliser ? [a-zA-Z\ ?*]}'', vous demande de choisir un objet que vous portez.
Ici, ``\texttt{a-zA-Z}'' sont les lettres d'inventaire des choix possibles.
Taper `\texttt{?}' affiche la liste de ces objets, ce qui vous permet de voir
à quoi correspond chaque lettre.  Dans cet exemple, il y a aussi un `\texttt{*}'
qui indique que vous pouvez choisir un objet absent de la liste, au cas où
vous voudriez utiliser quelque chose d'inattendu.  Taper `\texttt{*}' affiche tout votre
inventaire, ce qui vous permet de voir les lettres de tous les objets que
vous portez.  Enfin, si vous changez d'avis et décidez de ne plus exécuter
cette commande, vous pouvez appuyer sur la touche `ESC' (Échap) pour
l'annuler.

%.pg
Vous pouvez faire précéder certaines commandes d'un nombre pour les répéter
autant de fois ; par exemple, ``\texttt{10s}'' cherchera dix fois.  Si l'option
\textit{number\textunderscore pad}
est activée, vous devez taper `\texttt{n}' avant le nombre, si bien que l'exemple
ci-dessus s'écrirait plutôt ``\texttt{n10s}''.  Les commandes pour lesquelles un nombre
n'a pas de sens l'ignorent.  De plus, les commandes de déplacement peuvent être
précédées d'un préfixe pour un meilleur contrôle (voir ci-dessous).  Pour annuler
un nombre ou un préfixe, appuyez sur la touche
`ESC'.

%.pg
La liste des commandes est plutôt longue, mais vous pouvez la consulter à
tout moment pendant la partie grâce à la commande `\texttt{?}', qui donne accès
à un menu de textes d'aide.  Voici, pour référence, les associations de touches par défaut :

\begin{description}[font=\mdseries\ttfamily]
%.lp
\item[?]
Menu d'aide : affiche l'un des différents textes d'aide disponibles.
%.lp
\item[/]
La commande \texttt{whatis}, qui
indique ce que représente un symbole.  Vous pouvez désigner un emplacement
ou taper un symbole (voire un mot entier) à expliquer.
Pour désigner un emplacement, déplacez le curseur jusqu'à un endroit précis
de la carte, puis appuyez sur l'une des touches `\texttt{.}', `\texttt{,}', `\texttt{;}'
ou `\texttt{:}'.  `\texttt{.}' explique le symbole à l'emplacement choisi,
demande éventuellement ``\texttt{Plus d'informations ?}'' selon que l'option
`\textit{help}'
est activée ou non, puis vous demande de choisir un autre emplacement ;
`\texttt{,}' explique le symbole mais omet toute information
supplémentaire, puis vous laisse choisir un autre emplacement ;
`\texttt{;}' omet les informations supplémentaires et ne vous demande pas
non plus de choisir un autre emplacement à examiner ; `\texttt{:}' affiche les
informations supplémentaires éventuelles sans demander de confirmation.  Pendant le choix d'un emplacement,
la touche \texttt{ESC} met fin à cette commande, et la touche `\texttt{?}'
affiche un bref rappel de son fonctionnement.

%.lp ""
Si l'option
\textit{autodescribe}
est activée, une courte description de ce que vous voyez à chaque emplacement
s'affiche à mesure que vous déplacez le curseur.  Taper `\texttt{\#}' pendant le choix d'un emplacement
active ou désactive cette option.
L'option
\textit{whatis\textunderscore coord}
détermine si la courte description inclut les coordonnées sur la carte.

%.lp ""
Indiquer un nom plutôt qu'un emplacement
donne toujours toutes les informations supplémentaires disponibles sur ce nom.

%.lp ""
Vous pouvez aussi demander une description des monstres proches,
de tous les monstres actuellement affichés, des objets proches ou de tous les objets.
L'option
\textit{whatis\textunderscore coord}
détermine le format des coordonnées affichées avec leurs
descriptions.
%.lp
\item[\&]
Indique ce que fait une commande.
%.lp
\item[<]
Monte au niveau précédent (si vous êtes sur un escalier ou une échelle).
%.lp
\item[>]
Descend au niveau suivant (si vous êtes sur un escalier ou une échelle).
%.lp
\item[{[yuhjklbn]}]
Fait un pas dans la direction indiquée (voir la figure 3).  Si vous percevez
ou vous rappelez
un monstre à cet endroit, vous le combattrez à la place.  Seules ces
commandes de déplacement d'un pas vous font combattre les monstres ; les autres
(ci-dessous) sont ``sans danger''.
%.sd
\begin{center}
\begin{tabular}{cc}
\verb+   y  k  u   + & \verb+   7  8  9   +\\
\verb+    \ | /    + & \verb+    \ | /    +\\
\verb+   h- . -l   + & \verb+   4- . -6   +\\
\verb+    / | \    + & \verb+    / | \    +\\
\verb+   b  j  n   + & \verb+   1  2  3   +\\
                     & (si \textit{number\textunderscore pad} est activée)
\end{tabular}
\end{center}
%.ed
\begin{center}
Figure 3
\end{center}
%.lp
\item[{[YUHJKLBN]}]
Avance dans cette direction jusqu'à heurter un mur ou quelque chose.
%.lp
\item[m{[yuhjklbn]}]
Préfixe : se déplacer sans ramasser d'objets ni combattre (même si vous vous
rappelez un monstre à cet endroit).\\
%.lp ""
Quelques commandes autres que de déplacement utilisent le préfixe `\texttt{m}' pour demander
un fonctionnement par menu (afin de passer outre temporairement l'option
\textit{menustyle:Traditional}).
C'est surtout utile pour `\texttt{,}' (ramasser) quand une seule classe
d'objets est présente (auquel cas il n'y aura pas d'invite ``quelles sortes d'objets ?'',
et donc aucune occasion d'y répondre `\texttt{m}').
\\
%.lp ""
Ce préfixe
fait afficher par la commande ``\texttt{\#travel}'' un menu des destinations intéressantes en vue.
Il peut aussi servir avec les commandes `\texttt{\textbackslash}' (known, afficher la
liste de tous les objets découverts) et `\texttt{\`{}}' (knownclass,
afficher la liste des objets découverts d'une classe donnée) pour proposer
un menu de plusieurs ordres de tri (qui donne une nouvelle valeur à l'option
\textit{sortdiscoveries}) ;
de même pour les commandes ``\texttt{\#vanquished}'' et ``\texttt{\#genocided}'',
afin de proposer un menu de tri.
\\
%.lp ""
Quelques autres commandes (manger, offrir un sacrifice, appliquer une
trousse de mise en boîte, boire, tremper, vider un contenant) utilisent
le préfixe `\texttt{m}' pour ne pas chercher d'objets utilisables au
sol et passer directement à l'inventaire,
ou (pour ``\texttt{\#loot}'' afin de retirer une selle)
pour ignorer les contenants et passer directement aux monstres adjacents.
\\
%.lp ""
En mode débogage (alias ``mode magicien''), le préfixe `\texttt{m}' peut aussi
servir avec les commandes ``\texttt{\#teleport}'' et ``\texttt{\#wizlevelport}''.
%.lp
\item[F{[yuhjklbn]}]
Préfixe : combattre un monstre (même si vous ne faites que deviner sa présence).
%.lp
\item[g{[yuhjklbn]}]
Préfixe : se déplacer jusqu'à trouver quelque chose d'intéressant.
%.lp
\item[G{[yuhjklbn]} \textrm{\textmd{ou}} <Control>+{[yuhjklbn]}]
Préfixe : comme `\texttt{g}', mais un embranchement de couloirs n'est pas jugé
intéressant.
\\
Remarque : \texttt{<Control>+<touche>} signifie maintenir la touche \texttt{<Control>} ou
\texttt{<Ctrl>} enfoncée, comme \texttt{<Shift>} (Maj), pendant que vous tapez et relâchez
\texttt{<touche>}, puis relâcher \texttt{<Control>}.  \texttt{\textasciicircum <touche>} est utilisé
ailleurs dans ce Guide comme abréviation de la même chose.  Les caractères de
contrôle ne distinguent pas majuscules et minuscules, donc \texttt{\textasciicircum x} et \texttt{\textasciicircum X} sont identiques.
%.lp
\item[M{[yuhjklbn]}]
Les anciennes versions acceptaient `\texttt{M}' comme préfixe de déplacement
combinant l'effet de `\texttt{m}' avec \texttt{<Control>+<direction>}.
Ce préfixe n'existe plus, mais vous pouvez obtenir un effet similaire
en combinant \texttt{m} et \texttt{G<direction>}.
\texttt{m} peut aussi se combiner avec \texttt{g<direction>},
\texttt{<Control>+<direction>} ou \texttt{<Shift>+<direction>}.
%.lp
\item[\textunderscore]
Voyage vers un emplacement de la carte par un algorithme de plus court chemin.\\
%.lp ""
Le plus court chemin
est calculé sur les emplacements de la carte que le héros connaît (par exemple
vus ou déjà parcourus).
S'il n'existe aucun chemin connu, une estimation est faite à la place.
Le voyage s'arrête dans la plupart
des mêmes conditions que la commande `\texttt{G}', mais sans ramasser
d'objets, ce qui revient à imposer implicitement le préfixe `\texttt{m}'.
Sur les portages qui gèrent
la souris, la commande est aussi déclenchée par un clic sur un
emplacement autre que la position actuelle.
%.lp
\item[.]
Attend ou se repose : ne fait rien pendant un tour.
Faites-la précéder du préfixe `\texttt{m}'
pour attendre un tour même à côté d'un monstre hostile, si l'option \textit{safe\textunderscore wait}
est activée.
%.lp
\item[a]
Applique (utilise) un outil (pioche, clé, lampe \ldots).\\
%.lp ""
Utilisée sur une baguette, elle brise cette baguette, ce qui libère sa
magie.
Une confirmation est demandée.
%.lp
\item[A]
Retire un ou plusieurs objets portés, comme des pièces d'armure.\\
%.lp ""
Utilisez `\texttt{T}' (ôter) pour n'ôter qu'une seule pièce d'armure
ou `\texttt{R}' (retirer) pour ne retirer qu'un seul accessoire.
%.lp
\item[\textasciicircum A]
Répète la commande précédente.
%.lp
\item[c]
Ferme une porte.
%.lp
\item[C]
Nomme un monstre, un objet précis ou un type d'objet.\\
%.lp ""
Identique à la commande étendue ``\texttt{\#name}''.
%.lp
\item[\textasciicircum C]
Bouton panique : quitte la partie.
%.lp
\item[d]
Pose quelque chose.\\
Par exemple, \texttt{d7a} --- poser sept exemplaires de l'objet
\textit{a}.
%.lp
\item[D]
Pose plusieurs choses.\\
%.lp ""
En réponse à la question\\
``\texttt{Quels types d'objets voulez-vous lâcher ? [!\%= BUCXPaium]}''\\
vous devez taper zéro, un ou plusieurs symboles d'objets, éventuellement suivis de
`\texttt{a}' et/ou `\texttt{i}' et/ou `\texttt{u}' et/ou `\texttt{m}'.
Vous pouvez en outre taper un ou plusieurs des
groupes béni/\-non maudit/\-maudit.\\
%.sd
%.si
\texttt{DB}  --- pose tous les objets connus comme bénis.\\
\texttt{DU}  --- pose tous les objets connus comme non maudits.\\
\texttt{DC}  --- pose tous les objets connus comme maudits.\\
\texttt{DX}  --- pose tous les objets dont l'état B/U/C est inconnu.\\
\texttt{DP}  --- pose les derniers objets ramassés.\\
\texttt{Da}  --- pose tous les objets, sans demander de confirmation.\\
\texttt{Di}  --- examine votre inventaire avant de poser quoi que ce soit.\\
\texttt{Du}  --- ne pose que les objets non payés (dans une boutique).\\
\texttt{Dm}  --- utilise un menu pour choisir le ou les objets à poser.\\
\texttt{D\%u} --- ne pose que la nourriture non payée.
%.ei
%.ed
Le dernier exemple montre une combinaison.
Il existe quatre catégories de filtrage des objets : la classe (`\texttt{!}' pour
les potions, `\texttt{?}' pour les parchemins, etc.), l'état en boutique (`\texttt{u}'
pour non payé, autrement dit appartenant à la boutique), l'état béni/maudit
(`\texttt{B}', `\texttt{U}', `\texttt{C}' et `\texttt{X}', comme ci-dessus),
et la nouveauté (`\texttt{P}', objets récemment ramassés ; déterminée par le fait de
ramasser ou de poser des objets plutôt que par un quelconque facteur temps).
%.lp ""
\\
Si vous indiquez plus d'une valeur dans une catégorie (par exemple ``\texttt{!?}'' pour
les potions et les parchemins, ou ``\texttt{BU}'' pour béni et non maudit), un objet
de l'inventaire satisfait le critère s'il correspond à l'une quelconque des
valeurs indiquées (ainsi ``\texttt{!?}'' signifie `\texttt{!}' ou `\texttt{?}').
Si vous indiquez plus d'une catégorie, un objet de l'inventaire doit satisfaire
les critères de chacune des catégories (ainsi ``\texttt{\%u}'' signifie classe `\texttt{\%}'
et non payé `\texttt{u}').
Enfin, vous pouvez indiquer plusieurs valeurs dans plusieurs catégories :
``\texttt{!?BU}'' sélectionne toutes les potions et tous les parchemins connus comme
bénis ou non maudits.
(Dans les versions antérieures à 3.6, les combinaisons de filtres se comportaient différemment.)
%.lp
\item[\textasciicircum D]
Donne un coup de pied dans quelque chose (généralement une porte).
%.lp
\item[e]
Mange de la nourriture.\\
%.lp ""
Cherche normalement d'abord des aliments au sol, puis, s'il n'y en a aucun
ou si aucun n'est choisi, des aliments dans l'inventaire.
Faites précéder `\texttt{e}' du préfixe `\texttt{m}' pour ne pas essayer de manger
quoi que ce soit au sol.\\
%.lp ""
Si vous essayez de manger alors que vous êtes déjà rassasié, vous risquez de
vous étouffer à mort.
Si vous prenez ce risque, on vous demandera
``Continuer à manger ?'' \textit{si vous survivez à la première bouchée}.
Vous pouvez régler l'option
\textit{paranoid\textunderscore confirmation:eating}
pour exiger la réponse ``\texttt{oui}'' au lieu de simplement `\texttt{y}'.
%.lp
% Make sure Elbereth is not hyphenated below, the exact spelling matters.
% (Only specified here to parallel Guidebook.mn; use of \tt font implicitly
% prevents automatic hyphenation in TeX and LaTeX.)
\hyphenation{Elbereth}		%override the deduced syllable breaks
\item[E]
Grave un message sur le sol.\\
%.sd
%.si
\texttt{E-} --- écrire dans la poussière avec les doigts.\\
%.ei
%.ed
%.lp ""
Graver le mot ``\texttt{Elbereth}'' empêche la plupart des monstres de vous attaquer
au corps à corps (mais si vous attaquez, vous l'effacerez) ; c'est
souvent utile pour souffler un peu.
%.lp
\item[f]
Tire (projette ou lance) l'un des objets placés dans votre carquois (ou
sac de munitions, ou que vous tenez prêts).
Vous pouvez choisir les munitions au préalable avec la commande `\texttt{Q}', ou
laisser l'ordinateur choisir quelque chose d'approprié si l'option \textit{autoquiver} est activée.
Si l'arme que vous maniez revient quand on la lance, que votre carquois
est vide et que l'option \textit{autoquiver}
est désactivée, vous lancerez cette arme maniée au lieu de remplir le carquois.
Cette commande utilise aussi automatiquement une arme d'hast si vous en maniez une.
Si l'option \textit{fireassist} est activée, tirer essaiera automatiquement de manier un lanceur
(par exemple un arc ou une fronde) correspondant aux munitions du carquois ; cela peut
prendre plusieurs tours et être interrompu par un monstre.
Pensez à reprendre ensuite votre arme de corps à corps principale.
%.lp ""
\\
Voir aussi `\texttt{t}' (lancer) pour lancer et tirer de façon plus générale.
%.lp
\item[i]
Affiche votre inventaire (tout ce que vous portez).
%.lp
\item[I]
Affiche certaines parties de votre inventaire, généralement en indiquant le
caractère d'un ensemble d'objets, comme `\texttt{[}' pour les armures ou `\texttt{!}'
pour les potions.\\
%.sd
%.si
\texttt{I*} --- affiche toutes les gemmes de l'inventaire ;\\
\texttt{Iu} --- affiche tous les objets non payés ;\\
\texttt{Ix} --- affiche tous les objets consommés qui figurent sur votre facture ;\\
\texttt{IB} --- affiche tous les objets connus comme bénis ;\\
\texttt{IU} --- affiche tous les objets connus comme non maudits ;\\
\texttt{IC} --- affiche tous les objets connus comme maudits ;\\
\texttt{IX} --- affiche tous les objets dont l'état béni/maudit est inconnu ;\\
\texttt{IP} --- affiche les derniers objets ramassés ;\\
\texttt{I\$} --- compte votre argent.
%.ei
%.ed
%.lp
\item[o]
Ouvre une porte.
%.lp
\item[O]
Règle les options.\\
%.lp ""
Un menu indiquant les valeurs actuelles des options
s'affiche.  Vous pouvez modifier la plupart des valeurs simplement en sélectionnant
l'entrée de menu de l'option voulue (c'est-à-dire en tapant sa lettre ou en
cliquant dessus, selon votre interface).  Pour les options non booléennes,
un autre menu ou une invite apparaît une fois ce menu fermé.
Les options disponibles
sont décrites plus loin dans ce Guide.  Les options se règlent généralement avant la
partie plutôt qu'avec la commande `\texttt{O}' ; voir la section sur les options plus bas.
Faites précéder \texttt{O} du préfixe \texttt{m} pour afficher les options avancées.
%.lp
\item[\textasciicircum O]
Affiche le résumé du donjon.\\
%.lp ""
Raccourci de ``\texttt{\#overview}'' :
liste les niveaux intéressants du donjon déjà visités.\\
%.lp ""
(Avant la version 3.6.0, `\texttt{\textasciicircum O}' était une commande du mode débogage qui
indiquait l'emplacement de tous les niveaux spéciaux.
Utilisez ``\texttt{\#wizwhere}'' pour lancer cette commande.)
%.lp
\item[p]
Paie vos achats.
%.lp
\item[P]
Met un accessoire (anneau, amulette ou bandeau).\\
%.lp ""
Cette commande peut aussi servir à revêtir une armure.  L'invite
demandant quel objet de l'inventaire utiliser ne liste que les accessoires, mais choisir
une pièce d'armure non listée tentera de la revêtir.
(Voir la commande `\texttt{W}' plus bas.  Elle propose les armures comme choix
d'inventaire, mais accepte un accessoire et tente alors de le mettre.)
%.lp
\item[\textasciicircum P]
Répète le message précédent.\\
%.lp ""
Des \texttt{\textasciicircum P} successifs répètent des messages plus anciens.
Pour certaines interfaces, ce comportement peut être modifié par l'option
\textit{msg\textunderscore window}.
%.lp
\item[q]
Boit quelque chose (potion, eau, etc.).\\
%.lp ""
En présence d'une fontaine ou d'un évier, demande si vous voulez y
boire.
En cas de refus, propose de choisir une potion dans
l'inventaire.
Faites précéder \texttt{q} du préfixe \texttt{m} pour ne pas être interrogé
sur la fontaine ou l'évier.
%.lp
\item[Q]
Choisit un objet pour votre carquois, votre sac de munitions, ou simplement
à tenir prêt (un seul de ces emplacements est disponible à la fois).  Vous pouvez ensuite
le lancer (ou lancer l'un d'eux) avec la commande `\texttt{f}'.
%.lp
\item[r]
Lit un parchemin ou un grimoire.
%.lp
\item[R]
Retire un accessoire porté (anneau, amulette ou bandeau).\\
%.lp ""
Si vous en portez plusieurs, on vous demande lequel
retirer.  Si vous n'en portez qu'un, il est par défaut retiré
sans question, mais vous pouvez régler l'option
\textit{paranoid\textunderscore confirmation:Remove}
pour exiger une confirmation.\\
%.lp ""
Cette commande peut aussi servir à ôter une armure.  L'invite demandant quel
objet retirer ne liste que les accessoires portés, mais une pièce
d'armure portée peut être choisie.
(Voir la commande `\texttt{T}' plus bas.  Elle propose les armures comme choix
d'inventaire, mais accepte un accessoire et tente alors de le retirer.)
%.lp
\item[\textasciicircum R]
Redessine l'écran.
%.lp
\item[s]
Cherche les portes secrètes et les pièges autour de vous.
Il faut généralement plusieurs essais pour trouver quelque chose.
Faites-la précéder du préfixe `\texttt{m}' pour chercher pendant un tour
même à côté d'un monstre hostile, si l'option \textit{safe\textunderscore wait}
est activée.\\
%.lp ""
Peut aussi servir à déterminer s'il y a toujours un monstre sous un
marqueur adjacent de ``monstre mémorisé mais invisible''.
%.lp
\item[S]
Sauvegarde la partie (ce qui interrompt le jeu et quitte le programme).
La partie sauvegardée sera restaurée automatiquement la prochaine fois que
vous jouerez avec le même nom de personnage.\\
%.lp ""
En jeu normal, une fois la partie sauvegardée restaurée, le fichier
qui contenait la sauvegarde est supprimé.
En mode exploration, une fois la restauration effectuée, on vous demande si
vous voulez conserver ou supprimer le fichier.
Le conserver permet de jouer un moment, puis de quitter
sans sauvegarder et de restaurer à nouveau plus tard.\\
%.lp ""
Il n'existe pas de commande ``sauvegarder l'état actuel et continuer à jouer'',
pas même en mode exploration où les fichiers de sauvegarde peuvent être conservés et réutilisés.
%.lp
\item[t]
Lance un objet ou tire un projectile.\\
%.lp ""
Il n'existe pas de commande ``tirer'' distincte.
Si vous ``lancez'' une flèche en maniant un arc, vous tirez
cette flèche, et les bonus ou malus de compétence à l'arc s'appliquent.
Si vous ``lancez'' une flèche sans manier d'arc, vous la lancez
à la main, et elle sera généralement moins efficace que si elle était tirée.\\
%.lp ""
Voir aussi `\texttt{f}' (tirer) pour lancer ou tirer un objet présélectionné
avec la commande `\texttt{Q}' (carquois), avec une aide supplémentaire.
%.lp
\item[T]
Ôte une pièce d'armure.\\
%.lp ""
Si vous en portez plusieurs, on vous demande laquelle
ôter.  (Notez qu'une cape couvrant une armure de corps
et/ou une chemise, ou une armure de corps couvrant une chemise, est traitée
comme si les objets en dessous n'existaient pas.)
Si vous n'en portez qu'une, elle est par défaut
ôtée sans question, mais vous pouvez régler l'option
\textit{paranoid\textunderscore confirmation:Remove}
pour exiger une confirmation.\\
%.lp ""
Cette commande peut aussi servir à retirer des accessoires.  L'invite
demandant quel objet ôter ne liste que les armures portées, mais un accessoire
porté peut être choisi.
(Voir la commande `\texttt{R}' plus haut.  Elle propose les accessoires comme choix
d'inventaire, mais accepte une pièce d'armure et tente alors de l'ôter.)
%.lp
\item[\textasciicircum T]
Vous téléporte, si vous en avez la capacité.
%.lp
\item[v]
Affiche la liste des événements marquants de la partie en cours, également affichée par
\texttt{\#chronicle}.
\\
%.lp ""
`\texttt{v}' servait autrefois à afficher le numéro de version du jeu.
Utilisez désormais `\texttt{V}' pour cela.

%.lp
\item[V]
Affiche le numéro de version.
\\
%.lp ""
`\texttt{V}' servait autrefois à afficher un résumé de l'histoire du jeu.
C'est toujours possible avec \texttt{\#history}.
%.lp
\item[w]
Manie une arme.\\
%.sd
%.si
\texttt{w-} --- ne rien manier, utiliser vos mains nues (ou gantées).\\
%.ei
%.ed
Certains personnages peuvent manier deux armes à la fois ; utilisez pour cela la commande `\texttt{X}'
(ou la commande étendue ``\texttt{\#twoweapon}'').
%.lp
\item[W]
Revêt une pièce d'armure.\\
%.lp ""
Cette commande peut aussi servir à mettre un accessoire (anneau, amulette ou
bandeau).  L'invite demandant quel objet de l'inventaire utiliser ne liste que
les armures, mais choisir un accessoire non listé tentera de le mettre.
(Voir la commande `\texttt{P}' plus haut.  Elle propose les accessoires comme choix
d'inventaire, mais accepte une pièce d'armure et tente alors de la revêtir.)
%.lp
\item[x]
Échange l'arme que vous maniez avec l'objet de votre emplacement d'arme secondaire.\\
%.lp ""
Ce dernier sert d'arme secondaire lorsque vous combattez
à deux armes.  Notez que l'échange a lieu même si l'un de ces
emplacements est vide.
%.lp
\item[X]
Active ou désactive le combat à deux armes, si votre personnage en est capable.  Également disponible
avec la commande étendue ``\texttt{\#twoweapon}''.\\
%.lp ""
(Dans les versions antérieures à 3.6, cette touche lançait la commande qui fait passer du jeu
normal au ``mode exploration'', aussi appelé ``mode découverte'', désormais
déplacée vers ``\texttt{\#exploremode}'' et \texttt{M-X}.)
%.lp
\item[\textasciicircum X]
Affiche les informations de base sur votre personnage.\\
%.lp ""
Affiche le nom, le rôle, la race, le sexe (sauf si le nom du rôle le rend
redondant, comme \texttt{Femme des cavernes} ou \texttt{Prêtresse}) et l'alignement,
ainsi que votre divinité tutélaire et celle qui s'y oppose.  Affiche aussi
la plupart des informations de la ou des lignes d'état,
sous une forme moins concise, y compris plusieurs éléments supplémentaires qui
n'apparaissent pas dans l'affichage d'état normal faute de place.\\
%.lp ""
En jeu normal, c'est tout ce qu'affiche `\texttt{\textasciicircum X}'.
En mode exploration, les informations sur le rôle et l'état sont complétées par
celles que fournit la magie d'\textit{illumination}.
%.lp
\item[z]
Utilise une baguette.\\
%.sd
%.si
\texttt{z.} --- pour vous viser vous-même, utilisez `\texttt{.}' comme direction.
%.ei
%.ed
%.lp
\item[Z]
Lance un sort.\\
%.sd
%.si
\texttt{Z.} --- pour lancer un sort sur vous-même, utilisez `\texttt{.}' comme direction.
%.ei
%.ed
%.lp
\item[\textasciicircum Z]
Suspend la partie (versions UNIX avec contrôle des tâches uniquement).
Voir ``\#suspend'' plus bas pour plus de détails.
%.lp
\item[:]
Regarde ce qui se trouve ici.
%.lp
\item[;]
Indique à quel type de chose correspond un symbole visible.
%.lp
\item[,]
Ramasse des objets sur le sol, sous vous.\\
%.lp ""
Peut être précédée de `\texttt{m}' pour imposer un menu de sélection.
%.lp
\item[@]
Active ou désactive l'option \textit{autopickup}.
%.lp
\item[\textasciicircum]
Demande le type d'un piège adjacent que vous avez déjà découvert.
%.lp
\item[)]
Indique quelle arme vous maniez.
%.lp
\item[{]}]
Indique quelle armure vous portez.
%.lp
\item[=]
Indique quels anneaux vous portez.
%.lp
\item["]
Indique quelle amulette vous portez.
%.lp
\item[(]
Indique quels outils vous utilisez.
%.lp
\item[*]
Indique quel équipement vous utilisez.\\
%.lp ""
Réunit en une seule les cinq commandes précédentes
propres à chaque type.
%.lp
\item[\$]
Indique l'or que vous portez, et éventuellement aussi votre crédit et/ou votre dette en boutique.
%.lp
\item[+]
Affiche la liste des sorts que vous connaissez.\\
%.lp ""
Cette commande permet aussi de réorganiser
l'ordre dans lequel vos sorts sont listés, soit en triant toute la
liste, soit en choisissant un sort dans le menu puis un autre avec lequel
l'échanger.  Échanger des paires de sorts modifie leurs lettres de lancement,
si bien que le changement persiste après la fin de la commande `\texttt{+}' en cours.  Le tri
de la liste entière est temporaire.  Pour que le dernier ordre de tri persiste
au-delà de la commande `\texttt{+}' en cours, choisissez à nouveau l'option de tri, puis
``réattribuer les lettres pour garder l'ordre actuel''.  (Les sorts appris ensuite seront
ajoutés en fin de liste plutôt qu'insérés dans l'ordre
trié.)
%.lp
\item[\textbackslash]
Affiche les types d'objets découverts.
\\
%.lp ""
Peut être précédée de `\texttt{m}' pour choisir l'ordre d'affichage préféré.
%.lp
\item[\textasciigrave]
Affiche les types découverts d'une classe d'objets.
\\
%.lp ""
Peut être précédée de `\texttt{m}' pour choisir l'ordre d'affichage préféré.

%.lp
\item[|]
Si l'affichage permanent de l'inventaire est pris en charge et activé (avec l'option
\textit{perm\textunderscore invent}), agit sur cet affichage au lieu de la carte.
\\
%.lp ""
Permet de faire défiler avec les touches
menu\textunderscore first\textunderscore page, menu\textunderscore previous\textunderscore page,
menu\textunderscore next\textunderscore page et menu\textunderscore last\textunderscore page
(`\texttt{\textasciicircum}', `\texttt{<}', `\texttt{>}', `\texttt{|}' par défaut).
Certaines interfaces gèrent aussi les touches menu\textunderscore shift\textunderscore left et menu\textunderscore shift\textunderscore right
(`\texttt{\textbraceleft}' et `\texttt{\textbraceright}' par défaut).
Utilisez la touche \textit{Entrée} (\textit{Return} ou \textit{Enter}) ou \textit{Échap} pour
reprendre la partie.

%.lp
\item[!]
Passe dans un shell.
Voir ``\#shell'' plus bas pour plus de détails.
%.lp
\item[Del]
Affiche la carte sans obstacles.
Vous pouvez voir la partie explorée de la carte du niveau actuel sans
les monstres ; sans les monstres ni les objets ; ou sans les monstres, les objets
ni les pièges.\\
%.lp ""
La touche \texttt{<del>} est aussi indiquée \texttt{<delete>} (Suppr) sur certains claviers, ou
\texttt{<rubout>} sur d'autres.
Elle est parfois affichée \texttt{\textasciicircum ?} bien que ce ne soit pas un vrai
caractère de contrôle.\\
%.lp ""
Beaucoup de terminaux ont une option pour échanger les touches \texttt{<delete>} et \texttt{<backspace>},
si bien que taper \texttt{<del>} pourrait ne pas exécuter cette commande.
Dans ce cas, vous pouvez utiliser à la place la commande étendue ``\texttt{\#terrain}''.
%.lp
\item[\#]
Exécute une commande étendue.\\
%.lp ""
Comme vous le voyez, les auteurs de \textit{NetHack}
ont épuisé toutes les lettres ; c'est donc un moyen d'introduire les commandes
moins fréquemment utilisées.
Les commandes étendues disponibles dépendent des fonctionnalités avec lesquelles
le jeu a été compilé.
%.lp
\item[\#adjust]
Réorganise les lettres de l'inventaire (surtout utile quand l'option
\textit{fixinv}
est ``activée''). Complétion automatique. Touche par défaut : `\texttt{M-a}'.\\
%.lp ""
Cette commande permet de déplacer un objet d'un emplacement de l'inventaire
vers un autre, afin qu'il ait une lettre plus parlante pour vous
ou qu'il apparaisse à un endroit précis dans les listes
d'inventaire.
Vous pouvez le déplacer vers un emplacement actuellement vide ; si la destination est
occupée---et que les objets ne se fusionnent pas---l'objet
qui s'y trouve échange sa place avec celui que vous déplacez.
``\texttt{\#adjust}'' peut aussi servir à diviser une pile d'objets : au moment
de choisir l'objet à réorganiser, tapez un nombre avant sa lettre.\\
%.lp ""
Autrefois, réorganiser sans nombre rassemblait toutes les piles compatibles lors
du déplacement vers la destination.  Ce comportement a changé ; pour rassembler
des piles compatibles, utilisez ``\texttt{\#adjust}'' pour déplacer une pile vers son propre emplacement.
Si elle porte un nom, les autres piles de même nom ou sans
nom fusionneront, à condition que tous leurs autres attributs correspondent.
Si elle n'a pas de nom, seules les autres piles sans nom sont concernées.
Dans les deux cas, les piles par ailleurs compatibles mais portant un nom différent
ne fusionneront pas.  Ce comportement diffère de celui de ``\texttt{\#adjust}'' pour déplacer
un objet d'un emplacement vers un autre.  Dans ce cas, déplacer (sans
nombre) une pile compatible la fusionne si l'une des piles porte un
nom et pas l'autre, et donne ce nom au résultat, tandis que
diviser (avec un nombre) ignore le nom de la pile d'origine pour
décider de la fusion avec la pile de destination.
%.lp
\item[\#annotate]
Permet d'associer une ligne de texte au niveau actuel
du donjon.  Tous les niveaux annotés sont affichés par la
commande ``\texttt{\#overview}''. Complétion automatique.
Touche par défaut : `\texttt{M-A}',
et aussi `\texttt{\textasciicircum N}' si l'option \textit{number\textunderscore pad} est activée.
\\
%.lp ""
Faire précéder \#annotate du préfixe `\texttt{m}' revient à utiliser
\#overview avec ce préfixe.
%.lp
\item[\#apply]
Applique (utilise) un outil, comme une pioche, une clé ou une lampe.
Touche par défaut : `\texttt{a}'.\\
%.lp ""
Si l'outil utilisé agit sur des objets au sol, le préfixe `\texttt{m}'
ignore ces objets.\\
%.lp ""
Utilisée sur une baguette, elle brise cette baguette, ce qui libère sa
magie.  Une confirmation est demandée.
%.lp
\item[\#attributes]
Affiche vos attributs. Touche par défaut : `\texttt{\textasciicircum X}'.
%.lp
\item[\#autopickup]
Active ou désactive l'option \textit{autopickup}. Touche par défaut : `\texttt{@}'.
%.lp
\item[\#bugreport]
Ouvre une fenêtre de navigateur pour envoyer un rapport à l'\textit{équipe de développement
de NetHack}.
Peut être désactivée lors de la compilation du programme ; si elle est activée,
CRASHREPORTURL doit être défini dans le fichier de configuration système.
%.lp
\item[\#call]
Nomme un monstre, ou un objet de l'inventaire, au sol
ou dans la liste des découvertes, ou ajoute une annotation au
niveau actuel (comme ``\texttt{\#annotate}''). Touche par défaut : `\texttt{C}'.
%.lp
\item[\#cast]
Lance un sort. Touche par défaut : `\texttt{Z}'.
%.lp
\item[\#chat]
Parle à quelqu'un. Touche par défaut : `\texttt{M-c}'.
%.lp
\item[\#chronicle]
Affiche la liste des événements importants de la partie. Touche par défaut : `\texttt{v}'.
%.lp
\item[\#close]
Ferme une porte. Touche par défaut : `\texttt{c}'.
%.lp
\item[\#conduct]
Liste les défis volontaires que vous avez respectés. Complétion automatique.
Touche par défaut : `\texttt{M-C}'.\\
%.lp ""
Voir la section ``Conduite'' plus bas pour plus de détails.
%.lp
\item[\#debugfuzzer]
Lance le testeur aléatoire (fuzzer).
Mode débogage uniquement.
%.lp
\item[\#dip]
Trempe un objet dans quelque chose. Complétion automatique. Touche par défaut : `\texttt{M-d}'.\\
%.lp ""
Le préfixe \texttt{m} évite de tremper l'objet dans une fontaine ou un bassin
s'il y en a un à votre emplacement.
%.lp
\item[\#down]
Descend un escalier. Touche par défaut : `\texttt{>}'.
%.lp
\item[\#drop]
Pose un objet. Touche par défaut : `\texttt{d}'.
%.lp
\item[\#droptype]
Pose des types d'objets précis. Touche par défaut : `\texttt{D}'.
%.lp
\item[\#eat]
Mange quelque chose. Touche par défaut : `\texttt{e}'.
Le préfixe `\texttt{m}' ignore les aliments au sol.
%.lp
\item[\#engrave]
Grave une inscription sur le sol. Touche par défaut : `\texttt{E}'.
%.lp
\item[\#enhance]
Améliore ou consulte vos compétences d'armes et de sorts. Complétion automatique.
Touche par défaut : `\texttt{M-e}'.
%.lp
\item[\#exploremode]
Fait passer du jeu normal au mode exploration, sans score.
Touche par défaut : `\texttt{M-X}'.\\
%.lp ""
Demande une confirmation ; la réponse par défaut est `\texttt{n}' (non).
Pour passer réellement en mode exploration, répondez `\texttt{y}'.
Vous pouvez régler l'option
\textit{paranoid\textunderscore confirmation:quit}
pour exiger plutôt la réponse ``\texttt{oui}''.
%.lp
\item[\#fight]
Touche préfixe pour forcer le combat dans une direction, même si vous n'y voyez
rien à combattre.
Touche par défaut : `\texttt{F}', ou `\texttt{-}' avec l'option
\textit{number\textunderscore pad}
%.lp
\item[\#fire]
Tire les munitions du carquois, en maniant éventuellement un lanceur automatiquement,
ou frappe avec une arme d'hast maniée.
Touche par défaut : `\texttt{f}'.
%.lp
\item[\#force]
Force une serrure. Complétion automatique. Touche par défaut : `\texttt{M-f}'.
%.lp
\item[\#genocided]
Liste les types de monstres qui ont été génocidés.
En mode exploration et en mode débogage, affiche aussi les types qui ont
disparu.
\\
%.lp ""
L'ordre d'affichage est le même que celui de \texttt{\#vanquished}.
Le préfixe `\texttt{m}' ouvre un menu des ordres de tri disponibles, et
l'utiliser pour \texttt{\#genocided} ou pour \texttt{\#vanquished} modifie l'ordre des deux.
\\
%.lp ""
Si l'ordre de tri est ``par nombre, décroissant'' ou ``par nombre, croissant''
(applicables à \texttt{\#vanquished}), il est ignoré
pour \texttt{\#genocided} et l'ordre alphabétique est utilisé à la place.
Le menu omet ces deux choix lorsqu'il est utilisé pour \texttt{\#genocide}.
\\
%.lp ""
Complétion automatique.
Touche par défaut : `\texttt{M-g}'.
%.lp
\item[\#glance]
Indique à quel type de chose correspond un symbole de la carte. Touche par défaut : `\texttt{;}'.
%.lp
\item[\#help]
Affiche le menu d'aide.
Touche par défaut : `\texttt{?}',
et aussi `\texttt{h}' si l'option \textit{number\textunderscore pad} est activée.
%.lp
\item[\#herecmdmenu]
Affiche un menu des actions possibles visant votre emplacement actuel.
Le menu se limite à un sous-ensemble des actions les plus probables, et non à
une liste exhaustive de toutes les possibilités.
Complétion automatique.\\
%.lp ""
Si la souris est prise en charge et que l'option \textit{herecmd\textunderscore menu}
est activée, cliquer sur le héros (ou sur la monture, si vous êtes en selle)
exécute cette commande.
%.lp
\item[\#history]
Affiche un résumé de l'histoire du développement du jeu.
%.lp
\item[\#inventory]
Affiche votre inventaire. Touche par défaut : `\texttt{i}'.
%.lp
\item[\#inventtype]
Affiche l'inventaire de types d'objets précis. Touche par défaut : `\texttt{I}'.
%.lp
\item[\#invoke]
Invoque les pouvoirs spéciaux d'un objet. Complétion automatique. Touche par défaut : `\texttt{M-i}'.
%.lp
\item[\#jump]
Saute vers un autre endroit. Complétion automatique.
Touche par défaut : `\texttt{M-j}',
et aussi `\texttt{j}' si l'option \textit{number\textunderscore pad} est activée.
%.lp
\item[\#kick]
Donne un coup de pied dans quelque chose.
Touche par défaut : `\texttt{\textasciicircum D}',
et aussi `\texttt{k}' si l'option \textit{number\textunderscore pad} est activée.
%.lp
\item[\#known]
Affiche les types d'objets découverts.
Touche par défaut : `\texttt{\textbackslash}'.
\\
%.lp ""
Le préfixe `\texttt{m}' permet de donner une nouvelle valeur à l'option
\textit{sortdiscoveries}
pour choisir l'ordre d'affichage des découvertes.
%.lp
\item[\#knownclass]
Affiche les types découverts d'une classe d'objets.
Touche par défaut : `\texttt{`}'.
\\
%.lp ""
Le préfixe `\texttt{m}' fonctionne comme pour \texttt{\#known}.
%.lp
\item[\#levelchange]
Change votre niveau d'expérience.
Complétion automatique.
Mode débogage uniquement.
%.lp
\item[\#lightsources]
Affiche les sources de lumière mobiles.
Complétion automatique.
Mode débogage uniquement.
%.lp
\item[\#look]
Regarde ce qui se trouve ici, sous vous. Touche par défaut : `\texttt{:}'.
%.lp
\item[\#lookaround]
Décrit ce que vous voyez, ou ce dont vous vous souvenez, de vos alentours.
%.lp
\item[\#loot]
Pille une boîte ou un sac posé au sol sous vous, ou retire la selle
d'une monture qui se trouve à côté de vous. Complétion automatique.
Faites-la précéder du préfixe `\texttt{m}' pour ignorer les contenants à votre emplacement
et passer directement au retrait d'une selle.
Touche par défaut : `\texttt{M-l}',
et aussi `\texttt{l}' si l'option \textit{number\textunderscore pad} est activée.
%.lp
\item[\#monster]
Utilise la capacité spéciale d'un monstre (quand vous êtes métamorphosé en monstre).
Complétion automatique. Touche par défaut : `\texttt{M-m}'.
%.lp
\item[\#name]
Nomme un monstre, un objet précis ou un type d'objet.
Identique à ``\texttt{\#call}''.
Complétion automatique.
Touches par défaut : `\texttt{N}', `\texttt{M-n}' et `\texttt{M-N}'.
%.lp
\item[\#offer]
Offre un sacrifice aux dieux. Complétion automatique. Touche par défaut : `\texttt{M-o}'.\\
%.lp ""
Il vous faudra trouver un autel pour avoir la moindre chance de succès.
Les cadavres de monstres fraîchement tués sont les offrandes de choix.
%.lp ""
Le préfixe `\texttt{m}' évite d'offrir les objets posés sur l'autel.\\
%.lp
\item[\#open]
Ouvre une porte. Touche par défaut : `\texttt{o}'.
%.lp
\item[\#options]
Affiche et modifie le réglage des options. Touche par défaut : `\texttt{O}'.
Faites-la précéder du préfixe \texttt{m} pour afficher les options avancées.
%.lp
\item[\#optionsfull]
Affiche le réglage des options avancées du jeu.
Aucune touche par défaut.
Faites-la précéder du préfixe `\texttt{m}' pour exécuter la commande d'options simplifiée.
(Surtout utile si vous utilisez \texttt{BINDING=O:optionsfull} pour que
`\texttt{O}' revienne des options simplifiées aux options avancées traditionnelles.)
%.lp
\item[\#overview]
Affiche les informations que vous avez découvertes sur le donjon.
Tout niveau visité
% [note: amnesia no longer causes levels to be forgotten so exclude this]
% (unless forgotten due to amnesia)
portant une annotation est inclus,
et de nombreux éléments (autels, trônes, fontaines, etc. ; escaliers supplémentaires
menant à une autre branche du donjon) déclenchent une annotation automatique.
Si le résumé du donjon est choisi lors des révélations de fin de partie, tous les
niveaux visités sont inclus, qu'ils soient annotés ou non.
\\
%.lp ""
Faites précéder \#overview du préfixe `\texttt{m}' pour afficher le résumé
du donjon sous forme de menu, où vous pouvez choisir n'importe quel niveau visité
pour ajouter ou supprimer une annotation sans devoir y retourner.
Cela force aussi l'affichage de tous les niveaux visités, et pas seulement
du sous-ensemble des niveaux ``intéressants''.
\\
%.lp ""
Complétion automatique.
Touches par défaut : `\texttt{\textasciicircum O}' et `\texttt{M-O}'.
% DON'T PANIC!
%.lp
\item[\#panic]
Teste la routine de panique.
Met fin à la partie en cours.
Complétion automatique.
Mode débogage uniquement.\\
%.lp ""
Demande une confirmation ; la réponse par défaut est `\texttt{n}' (non) : la partie continue.
Pour paniquer réellement, répondez `\texttt{y}'.
Vous pouvez régler l'option
\textit{paranoid\textunderscore confirmation:quit}
pour exiger plutôt la réponse ``\texttt{oui}''.
%.lp
\item[\#pay]
Paie vos achats. Touche par défaut : `\texttt{p}'.
%.lp
\item[\#perminv]
Si l'affichage permanent de l'inventaire est pris en charge et activé (avec
l'option \textit{perm\textunderscore invent}), agit sur cet affichage au lieu de la carte.
On vous demandera des touches de défilement du menu telles
que `\texttt{>}' et `\texttt{<}'.
Appuyez sur \texttt{Entrée} ou \texttt{Échap} pour reprendre le jeu normal.
Touche par défaut : \texttt{|}.
%.lp
\item[\#pickup]
Ramasse des objets à l'endroit actuel. Touche par défaut : `\texttt{,}'.
Le préfixe `\texttt{m}' impose l'utilisation d'un menu.
%.lp
\item[\#polyself]
Vous métamorphose.
Complétion automatique.
Mode débogage uniquement.
%.lp
\item[\#pray]
Prie les dieux de vous venir en aide. Complétion automatique. Touche par défaut : `\texttt{M-p}'.\\
%.lp ""
Prier trop tôt après avoir déjà reçu de l'aide est une mauvaise idée.
(Indice : entrer vivant dans le donjon compte comme avoir reçu de l'aide.
Mieux vaut sans doute ne pas commencer une nouvelle partie en priant tout de suite.)
Comme utiliser cette commande par accident peut causer des ennuis, une
option vous demande de confirmer votre intention avant de prier.  Elle est activée
par défaut, et vous pouvez modifier l'option
\textit{paranoid\textunderscore confirmation}
pour la désactiver.
%.lp
\item[\#prevmsg]
Affiche les messages du jeu précédemment affichés. Touche par défaut : `\texttt{\textasciicircum P}'.
%.lp
\item[\#puton]
Met un accessoire (anneau, amulette, etc.). Touche par défaut : `\texttt{P}'.
%.lp
\item[\#quaff]
Boit quelque chose. Touche par défaut : `\texttt{q}'.\\
%.lp ""
Le préfixe \texttt{m} évite de boire à une fontaine ou à un évier
s'il y en a un à votre emplacement.
%.lp
\item[\#quit]
Quitte le programme sans sauvegarder la partie. Complétion automatique.\\
%.lp ""
Comme utiliser cette commande par accident ferait perdre la partie en cours,
on vous demande de confirmer votre intention avant de quitter.
La réponse par défaut est `\texttt{n}' (non) : la partie continue.
Pour quitter réellement, répondez `\texttt{y}'.
Vous pouvez régler l'option
\textit{paranoid\textunderscore confirmation:quit}
pour exiger plutôt la réponse ``\texttt{oui}''.
%.lp
\item[\#quiver]
Choisit les munitions du carquois. Touche par défaut : `\texttt{Q}'.
%.lp
\item[\#read]
Lit un parchemin, un grimoire ou autre chose. Touche par défaut : `\texttt{r}'.
%.lp
\item[\#redraw]
Redessine l'écran.
Touche par défaut : `\texttt{\textasciicircum R}',
et aussi `\texttt{\textasciicircum L}' si l'option \textit{number\textunderscore pad} est activée.
%.lp
\item[\#remove]
Retire un accessoire (anneau, amulette, etc.). Touche par défaut : `\texttt{R}'.
%.lp
\item[\#repeat]
Répète la commande précédente.
Touche par défaut :~`\texttt{\textasciicircum A}'.
%.lp
\item[\#reqmenu]
Touche préfixe pour modifier le comportement de certaines commandes ou leur demander un menu.
Empêche le ramassage automatique quand elle est utilisée avec les commandes de déplacement.
Touche par défaut : `\texttt{m}'.
%.lp
\item[\#retravel]
Voyage vers la destination de voyage choisie précédemment.
Touche par défaut : `\texttt{C-\textunderscore }'.
Voir aussi \texttt{\#travel}.
%.lp
\item[\#ride]
Monte (ou descend de) une créature sellée. Complétion automatique.
Touche par défaut : `\texttt{M-R}'.
%.lp
\item[\#rub]
Frotte une lampe ou une pierre. Complétion automatique. Touche par défaut : `\texttt{M-r}'.
%.lp
\item[\#run]
Touche préfixe pour courir dans une direction.
Touche par défaut : `\texttt{G}' quand l'option
\textit{number\textunderscore pad}
est désactivée,
`\texttt{5}' quand
\textit{number\textunderscore pad}
vaut 1~ou~3,
et sinon `\texttt{M-5}' quand elle vaut 2~ou~4.
%.lp
\item[\#rush]
Touche préfixe pour foncer dans une direction.
Touche par défaut : `\texttt{g}' quand l'option
\textit{number\textunderscore pad}
est désactivée,
`\texttt{M-5}' quand
\textit{number\textunderscore pad}
vaut 1~ou~3,
et sinon `\texttt{5}' quand elle vaut 2~ou~4.
%.lp
\item[\#save]
Sauvegarde la partie et quitte le programme.
Touche par défaut : `\texttt{S}'.
%.lp
\item[\#saveoptions]
Sauvegarde les options de configuration dans le fichier de configuration.
Cela écrase le fichier et en supprime tous les commentaires ; si vous avez
modifié le fichier de configuration à la main, n'utilisez donc pas cette commande.
%.lp
\item[\#search]
Cherche les pièges et les portes secrètes autour de vous. Touche par défaut : `\texttt{s}'.
%.lp
\item[\#seeall]
Affiche tout l'équipement utilisé. Touche par défaut : `\texttt{*}'.
%.lp ""
\\
Affiche les objets utilisés dans un menu, même s'il n'y en a qu'un.
%.lp
\item[\#seeamulet]
Affiche l'amulette actuellement portée. Touche par défaut : `\texttt{"}'.
%.lp ""
\\
Le préfixe `\texttt{m}' impose l'affichage de l'amulette portée
dans un menu plutôt que par un simple message.
%.lp
\item[\#seearmor]
Affiche l'armure actuellement portée. Touche par défaut : `\texttt{[}'.
%.lp ""
\\
Affiche l'armure portée dans un menu, même si une seule pièce est portée.
%.lp
\item[\#seerings]
Affiche le ou les anneaux actuellement portés. Touche par défaut : `\texttt{=}'.
%.lp ""
Affiche les anneaux portés dans un menu s'il y en a deux (ou s'il n'y en a
qu'un et que c'est un anneau de viande plutôt qu'un ``vrai'' anneau).
Utilisez le préfixe `\texttt{m}' pour imposer un menu pour un seul anneau.
%.lp
\item[\#seetools]
Affiche les outils actuellement utilisés. Touche par défaut : `\texttt{(}'.
%.lp ""
Affiche le résultat dans un message si un seul outil est utilisé (bandeau,
serviette ou lentilles portés, lampe(s) et/ou bougie(s) allumée(s), laisses
attachées à des familiers).
Affiche un menu s'il y en a plusieurs ou si la commande est
précédée du préfixe `\texttt{m}'.
%.lp
\item[\#seeweapon]
Affiche l'arme actuellement maniée. Touche par défaut : `\texttt{)}'.
%.lp ""
Si vous maniez deux armes, un message distinct est donné pour l'arme
secondaire.
Le préfixe `\texttt{m}' impose un menu qui inclut
l'arme principale, l'arme de rechange même hors combat à deux armes, ainsi que
ce qui est actuellement placé dans l'emplacement du carquois.
%.lp
\item[\#shell]
Passe dans un shell, en quittant NetHack pour un sous-processus.
Peut être désactivée lors de la compilation du programme.
Si elle est activée, l'accès peut être réglé utilisateur par utilisateur dans
le fichier de configuration système.
Utilisez la commande shell `\texttt{exit}' pour revenir au jeu.
Touche par défaut : `\texttt{!}'.
%.lp
\item[\#showgold]
Indique l'or de votre inventaire, y compris l'or que vous savez présent dans
les contenants que vous portez.  Si vous êtes dans une boutique, indique aussi
votre crédit ou votre dette dans cette boutique.
Touche par défaut : `\texttt{\$}'.
%.lp
\item[\#showspells]
Liste et réordonne les sorts connus.
Touche par défaut : `\texttt{+}'.
%.lp
\item[\#showtrap]
Décrit un piège adjacent, éventuellement recouvert d'objets ou d'un monstre.
Pour cela, le piège doit déjà avoir été découvert.
(La commande ``\texttt{\#terrain}'' peut afficher votre carte en retirant
temporairement tous les objets et monstres, ce qui permet de voir tous les
pièges découverts.)
Touche par défaut : `\texttt{\textasciicircum }'.
%.lp
\item[\#sit]
Vous assoit. Complétion automatique. Touche par défaut : `\texttt{M-s}'.
%.lp
\item[\#stats]
Affiche les statistiques d'utilisation de la mémoire.
Complétion automatique.
Mode débogage uniquement.
%.lp
\item[\#suspend]
Suspend la partie, en quittant NetHack pour revenir au terminal depuis lequel
il a été lancé, sans sauvegarder ni quitter.
Peut être désactivée lors de la compilation du programme.
Si elle est activée, surtout utile pour les interfaces \textit{tty} et \textit{curses} sous
%.UX \. \" yields "UNIX."
UNIX.
Utilisez la commande shell `\texttt{fg}' pour revenir au jeu.
Touche par défaut : `\texttt{\textasciicircum Z}'.
%.lp
\item[\#swap]
Échange l'arme maniée et l'arme secondaire. Touche par défaut : `\texttt{x}'.
%.lp
\item[\#takeoff]
Ôte une pièce d'armure. Touche par défaut : `\texttt{T}'.
%.lp
\item[\#takeoffall]
Retire toute l'armure. Touche par défaut : `\texttt{A}'.
%.lp
\item[\#teleport]
Vous téléporte sur le niveau. Touche par défaut : `\texttt{\textasciicircum T}'.
%.lp
\item[\#terrain]
Affiche la carte sans obstacles.
En jeu normal, vous pouvez voir la partie explorée de la carte du niveau
actuel sans les monstres ; sans les monstres ni les objets ; ou sans les monstres,
les objets ni les pièges.\\
%.lp ""
Les nuages de gaz visibles sont traités comme des pièges
pour décider s'il faut les afficher ou montrer le sol en dessous.\\
%.lp ""
En mode exploration, vous pouvez choisir d'afficher la carte complète plutôt
que sa seule partie explorée.
En mode débogage, d'autres choix sont proposés.\\
%.lp ""
Complétion automatique.
Touche par défaut : `\texttt{<del>}' ou `\texttt{<delete>}' (voir \textit{Del} plus haut).
%.lp
\item[\#therecmdmenu]
Affiche un menu des actions possibles visant un emplacement voisin du vôtre.
Le menu se limite à un sous-ensemble des actions les plus probables, et non à
une liste exhaustive de toutes les possibilités.
Complétion automatique.
%%--invoking it by mouse seems to be broken
%% \\
%% .lp ""
%% If mouse support is enabled and the \textit{herecmd\textunderscore menu}
%% option is On, clicking on an adjacent location will execute this command.
%.lp
\item[\#throw]
Lance quelque chose. Touche par défaut : `\texttt{t}'.
%.lp
\item[\#timeout]
Examine la file des minuteries.
Complétion automatique.
Mode débogage uniquement.
%.lp
\item[\#tip]
Renverse un contenant (sac ou boîte) pour en vider le contenu.
S'il y a des contenants au sol, le jeu vous demande d'en choisir un
ou de ``vider un objet que vous portez''.
\\
%.lp ""
Si vous choisissez cette dernière possibilité, une autre invite vous demande
quel objet de l'inventaire vider.
Le préfixe `\texttt{m}' fait ignorer à la commande les contenants au
sol pour en choisir un dans l'inventaire, sauf dans le cas particulier de
\textit{menustyle:Traditional}
avec deux contenants ou plus présents ; dans ce cas, le jeu commence par le
menu des contenants au sol.
\\
%.lp ""
Complétion automatique. Touche par défaut : `\texttt{M-T}'.
%.lp
\item[\#toggle]
Active ou désactive une option booléenne.
Exige un paramètre entre parenthèses : le nom de l'option à basculer.
L'option doit être modifiable en cours de partie.

%.lp ""
Par exemple :
%.sd
\begin{verbatim}
    BIND=':toggle(price_quotes)
    BIND=@:toggle(autopickup)
\end{verbatim}
%.ed

%.lp
\item[\#travel]
Voyage vers un endroit précis de la carte.
Touche par défaut : `\texttt{\textunderscore }'.
Le préfixe ``demande de menu'' affiche un menu des destinations intéressantes en vue,
sans demander de déplacer le curseur.
Si vous choisissez une destination avec le curseur et que l'option \textit{autodescribe}
est activée, la ligne du haut affichera ``(pas de chemin)'' si
votre personnage ne connaît aucun chemin vers cet endroit.
Voir aussi \texttt{\#retravel}.
%.lp
\item[\#turn]
Repousse les morts-vivants. Complétion automatique. Touche par défaut : `\texttt{M-t}'.
%.lp
\item[\#twoweapon]
Active ou désactive le combat à deux armes. Complétion automatique.
Touche par défaut : `\texttt{X}',
et aussi `\texttt{M-2}' si l'option \textit{number\textunderscore pad} est désactivée.\\
%.lp ""
Notez que vous devez
utiliser des armes adaptées à ce type de combat, faute de quoi il sera
automatiquement désactivé.
%.lp
\item[\#untrap]
Désamorce quelque chose (piège, porte ou coffre).
Touche par défaut : `\texttt{M-u}', et `\texttt{u}' si l'option \textit{number\textunderscore pad} est activée.\\
%.lp ""
Dans certaines circonstances, sert aussi à délivrer des monstres pris au piège.
%.lp
\item[\#up]
Monte un escalier. Touche par défaut : `\texttt{<}'.
%.lp
\item[\#vanquished]
Liste les monstres vaincus, par type et par nombre.
\\
%.lp ""
Notez que la liste des monstres vaincus inclut tous les monstres tués par
des pièges ou par d'autres monstres, en plus de ceux que vous avez tués, et omet
ceux qui ont été retirés du jeu sans être tués (par exemple par génocide, ou par
un boutiquier apaisé renvoyant les flics de Keystone appelés) ou qui étaient déjà
des cadavres lorsqu'ils ont été placés sur la carte.
\\
%.lp ""
Le préfixe ``demande de menu'' avant \#vanquished ouvre
un menu des ordres de tri disponibles (à condition que la liste des monstres vaincus
contienne au moins deux types de monstres).
L'ordre choisi est affecté à l'option \textit{sortvanquished}
et donc mémorisé pour les utilisations suivantes de \#vanquished.
La commande \texttt{\#genocided} partage cet ordre de tri.
\\
%.lp ""
Lors des révélations de fin de partie, quand on vous demande s'il faut afficher
les monstres vaincus, répondre `\texttt{a}' vous laisse choisir dans le menu de tri.
\\
%.lp ""
Complétion automatique.
Touche par défaut : `\texttt{M-V}'.
%.lp
\item[\#version]
Affiche les options de compilation de cette version de \textit{NetHack}.

%.lp
Le deuxième paragraphe liste la ou les interfaces utilisateur incluses.
S'il y en a plusieurs, vous pouvez utiliser l'option \textit{windowtype}
dans votre fichier de configuration pour choisir celle que vous voulez.

%.lp
Complétion automatique. Touche par défaut : `\texttt{M-v}'.
%.lp
\item[\#versionshort]
Affiche le numéro de version du programme, ainsi que la date et l'heure auxquelles
l'exécutable a été compilé à partir des sources (et non la date de sortie de la version).
Touche par défaut : `\texttt{V}'.
%.lp
\item[\#vision]
Affiche la matrice de vision.
Complétion automatique.
Mode débogage uniquement.
%.lp
\item[\#wait]
Vous repose un tour sans rien faire.
Touche par défaut : `\texttt{.}', et aussi `{\tt{ }}' si l'option
\textit{rest\textunderscore on\textunderscore space} est activée.
%.lp
\item[\#wear]
Revêt une pièce d'armure. Touche par défaut : `\texttt{W}'.
%.lp
\item[\#whatdoes]
Indique ce que fait une touche. Touche par défaut : `\texttt{\&}'.
%.lp
\item[\#whatis]
Indique à quel type de chose correspond un symbole. Touche par défaut : `\texttt{/}'.
%.lp
\item[\#wield]
Manie une arme. Touche par défaut : `\texttt{w}'.
%.lp
\item[\#wipe]
Vous essuie le visage. Complétion automatique. Touche par défaut : `\texttt{M-w}'.
%.lp
\item[\#wizborn]
Affiche les statistiques de naissance, de mort, de génocide et de disparition des monstres.
Mode débogage uniquement.
%.lp
\item[\#wizbury]
Enterre les objets sous vous et autour de vous.
Complétion automatique.
Mode débogage uniquement.
%.lp
\item[\#wizcast]
Lance n'importe quel sort.
Mode débogage uniquement.
%.lp
\item[\#wizdetect]
Révèle les choses cachées (portes secrètes, pièges ou monstres invisibles)
dans un rayon modeste.
Aucun temps ne s'écoule.
Complétion automatique.
Mode débogage uniquement.
Touche par défaut : `\texttt{\textasciicircum E}'.
%.lp
\item[\#wizgenesis]
Crée un monstre.
Peut être précédée d'un nombre pour en créer plusieurs.
Complétion automatique.
Mode débogage uniquement.
Touche par défaut : `\texttt{\textasciicircum G}'.
%.lp
\item[\#wizidentify]
Identifie tous les objets de l'inventaire.
Complétion automatique.
Mode débogage uniquement.
Touche par défaut : `\texttt{\textasciicircum I}'.
%.lp
\item[\#wizintrinsic]
Donne une ou plusieurs intrinsèques.
Complétion automatique.
Mode débogage uniquement.
%.lp
\item[\#wizkill]
Retire des monstres du jeu en les désignant simplement.
Par défaut, le héros reçoit le mérite ou le blâme de leur mort.
Faites précéder cette commande du préfixe `\texttt{m}' pour l'éviter.
Complétion automatique.
Mode débogage uniquement.
%.lp
\item[\#wizlevelport]
Vous téléporte vers un autre niveau.
Complétion automatique.
Mode débogage uniquement.
Touche par défaut : `\texttt{\textasciicircum V}'.
%.lp
\item[\#wizmap]
Cartographie le niveau.
Complétion automatique.
Mode débogage uniquement.
Touche par défaut : `\texttt{\textasciicircum F}'.
%.lp
\item[\#wizrumorcheck]
Vérifie les limites des rumeurs en affichant la première et la dernière rumeur
vraie, ainsi que la première et la dernière rumeur fausse.\\
%.lp ""
Affiche aussi la première, la deuxième et la dernière gravure aléatoire, épitaphe
et monstre hallucinatoire.\\
%.lp ""
Complétion automatique.
Mode débogage uniquement.
%.lp
\item[\#wizseenv]
Affiche les vecteurs de vision des cases de la carte.
Complétion automatique.
Mode débogage uniquement.
%.lp
\item[\#wizsmell]
Renifle un monstre.
Complétion automatique.
Mode débogage uniquement.
%.lp
\item[\#wizwhere]
Affiche l'emplacement des niveaux spéciaux.
Complétion automatique.
Mode débogage uniquement.
%.lp
\item[\#wizwish]
Fait un vœu.
Complétion automatique.
Mode débogage uniquement.
Touche par défaut : `\texttt{\textasciicircum W}'.
Faites précéder cette commande du préfixe `\texttt{m}' pour afficher un menu de l'historique des vœux.
%.lp
\item[\#wmode]
Affiche les modes des murs.
Complétion automatique.
Mode débogage uniquement.
%.lp
\item[\#zap]
Utilise une baguette. Touche par défaut : `\texttt{z}'.
%.lp
\item[\#?]
Menu d'aide : affiche la liste des commandes étendues disponibles.
\end{description}

%.pg
\noindent Si votre clavier possède une touche méta (qui, pressée en combinaison
avec une autre touche, la modifie en activant le bit `méta' [8e bit, ou bit
`de poids fort']), vous pouvez lancer de nombreuses commandes étendues en tapant
la première lettre de la commande avec méta.

Sous \textit{Windows} et \textit{MS-DOS},
la touche `Alt' peut être utilisée de cette façon.
Sur les autres systèmes, si taper `Alt' plus une autre touche transmet une
séquence de deux caractères composée d'un \texttt{Escape} (Échap)
suivi de l'autre touche, vous pouvez activer l'option \textit{altmeta}
pour que \textit{NetHack} les combine en \texttt{méta+<touche>}.
(Cette combinaison n'a lieu que lorsque NetHack attend une
commande à exécuter, et non lors de la saisie d'un nom ou
d'un vœu.)

%.pg
Contrairement aux caractères de contrôle, pour lesquels \texttt{\textasciicircum x} et \texttt{\textasciicircum X} désignent la même
chose, les caractères méta distinguent majuscules et minuscules : \texttt{M-x} et \texttt{M-X}
représentent des choses différentes.  Certaines commandes accessibles par un caractère
méta exigent une lettre majuscule parce que la minuscule
correspondante sert à une autre commande ; il faut alors la combinaison de trois touches
\texttt{méta+Shift+lettre}.

%.BR 1
\begin{description}[font=\mdseries\ttfamily]
%.lp
\item[M-?]
{\tt\#?} (non pris en charge sur toutes les plateformes)
%.lp
\item[M-2]
{\tt\#twoweapon} (sauf si l'option \textit{number\textunderscore pad} est activée)
%.lp
\item[M-a]
{\tt\#adjust}
%.lp
\item[M-A]
{\tt\#annotate}
%.lp
\item[M-c]
{\tt\#chat}
%.lp
\item[M-C]
{\tt\#conduct}
%.lp
\item[M-d]
{\tt\#dip}
%.lp
\item[M-e]
{\tt\#enhance}
%.lp
\item[M-f]
{\tt\#force}
%.lp
\item[M-g]
{\tt\#genocided}
%.lp
\item[M-i]
{\tt\#invoke}
%.lp
\item[M-j]
{\tt\#jump}
%.lp
\item[M-l]
{\tt\#loot}
%.lp
\item[M-m]
{\tt\#monster}
%.lp
\item[M-n]
{\tt\#name}
%.lp
\item[M-o]
{\tt\#offer}
%.lp
\item[M-O]
{\tt\#overview}
%.lp
\item[M-p]
{\tt\#pray}
%.Ip
\item[M-r]
{\tt\#rub}
%.lp
\item[M-R]
{\tt\#ride}
%.lp
\item[M-s]
{\tt\#sit}
%.lp
\item[M-t]
{\tt\#turn}
%.lp
\item[M-T]
{\tt\#tip}
%.lp
\item[M-u]
{\tt\#untrap}
%.lp
\item[M-v]
{\tt\#version}
%.lp
\item[M-V]
{\tt\#vanquished}
%.lp
\item[M-w]
{\tt\#wipe}
%.lp
\item[M-X]
{\tt\#exploremode}
\end{description}

%.pg
\noindent Si l'option \textit{number\textunderscore pad} est activée, quelques commandes par lettre
supplémentaires sont disponibles :
\begin{description}[font=\mdseries\ttfamily]
%.lp
\item[h]
{\tt\#help}
%.lp
\item[j]
{\tt\#jump}
%.lp
\item[k]
{\tt\#kick}
%.lp
\item[l]
{\tt\#loot}
%.lp
\item[N]
{\tt\#name}
%.lp
\item[u]
{\tt\#untrap}
\end{description}

%.BR 1 \"blank line for extra separation; plain text output looks better

%.hn 1
\section{Pièces et couloirs}

%.pg
Les pièces et les couloirs du donjon sont soit éclairés, soit sombres.
Toutes les zones éclairées situées dans votre champ de vision sont affichées ;
les zones sombres ne sont affichées que si elles sont à une case de vous.
Les murs et les couloirs restent sur la carte à mesure que vous les explorez.

%.pg
Les couloirs secrets sont cachés et ressemblent à de la roche pleine.
Vous pouvez les trouver avec la commande `\texttt{s}' (fouiller) lorsque vous
êtes à côté d'eux.
Plusieurs tentatives de fouille peuvent être nécessaires.
Quand la fouille réussit, les couloirs secrets deviennent des cases de
couloir ordinaires et ouvertes.
La magie de cartographie révèle les couloirs secrets ; elle les convertit donc
en couloirs ordinaires et les affiche comme tels.

%.hn 2
\subsection*{Embrasures de porte}

%.pg
Les embrasures de porte relient les pièces et les couloirs.
Certaines embrasures n'ont pas de porte ; vous pouvez passer directement.
D'autres sont munies d'une porte, qui peut être ouverte, fermée ou verrouillée.
Pour ouvrir une porte fermée, utilisez la commande `\texttt{o}' (ouvrir) ;
pour la refermer, utilisez la commande `\texttt{c}' (fermer).
Par défaut, l'option
\textit{autoopen}
est activée : il suffit donc d'essayer de marcher sur la case d'une porte
fermée pour tenter de l'ouvrir, sans avoir besoin de `\texttt{o}'.
L'ouverture via
\textit{autoopen}
ne fonctionne pas si vous êtes \textit{confus} ou \textit{étourdi} ou si vous souffrez
de \textit{maladresse}.

%.pg
On ne peut pas franchir une porte ouverte en diagonale ; vous devez l'aborder
de face, horizontalement ou verticalement.
Les embrasures sans porte ne sont
pas soumises à cette restriction, sauf sur un niveau particulier
%.\" the rogue level
(décrit par ``\texttt{\#overview}'' comme ``une zone primitive'').

%.pg
Il existe une magie de déverrouillage, mais elle n'est généralement pas
disponible en début de partie.
Vous pouvez franchir une porte verrouillée sans magie en utilisant d'abord
un outil de déverrouillage avec la commande `\texttt{a}' (appliquer), puis en l'ouvrant.
Par défaut, l'option
\textit{autounlock}
est également activée : si vous tentez d'ouvrir (via `\texttt{o}' ou
\textit{autoopen})
une porte verrouillée alors que vous portez un outil de déverrouillage, on vous
demandera si vous voulez l'utiliser sur la serrure de la porte.
Vous pouvez aussi enfoncer une porte fermée (verrouillée ou non)
en lui donnant des coups de pied via la commande ``\texttt{\textasciicircum D}'' (coup de pied).
Enfoncer une porte la détruit et fait beaucoup de bruit, ce qui peut
réveiller les monstres endormis.

%.pg
Certaines portes fermées sont piégées et explosent si l'on tente
de les ouvrir (non verrouillées), de les déverrouiller (verrouillées) ou de
les enfoncer.
Comme un coup de pied, une explosion détruit la porte et fait beaucoup de bruit.
La commande ``\texttt{\#untrap}'' permet de chercher des pièges sur une porte, mais
plusieurs tentatives peuvent être nécessaires pour en trouver un.
Lorsqu'un piège est trouvé, on vous demande si vous voulez essayer de le désamorcer.
Si vous acceptez, une réussite élimine le piège, mais
un échec déclenche son explosion.
(Si vous refusez, vous oubliez en pratique qu'un piège a été trouvé à cet endroit.)

%.pg
Les portes fermées peuvent servir à tenir les monstres à distance.
La plupart des monstres ne savent pas ouvrir les portes fermées, mais quelques-uns
n'en ont pas besoin (par exemple, les fantômes traversent les portes et les
nuages de brouillard s'infiltrent dessous).
Certains monstres capables d'ouvrir les portes savent aussi utiliser les outils
de déverrouillage.
Et certains (les géants) peuvent fracasser les portes.

%.pg
Les portes secrètes sont cachées et ressemblent à un mur ordinaire (depuis
l'intérieur d'une pièce) ou à de la roche pleine (depuis l'extérieur).
Vous pouvez les trouver avec la commande `\texttt{s}' (fouiller), mais cela peut
demander plusieurs essais (voire de nombreux essais si votre chance est faible).
Une fois trouvées, elles sont en tout point équivalentes à des portes normales.
La magie de cartographie ne révèle pas les portes secrètes.

%.hn 2
\subsection*{Pièges (`\texttt{\textasciicircum }')}

%.pg
Le donjon est parsemé de pièges destinés à attraper l'intrus imprudent.
Par exemple, vous pouvez soudain tomber dans une fosse et rester coincé
quelques tours à essayer d'en sortir (voir plus bas).
En général, un piège n'apparaît pas sur votre carte tant que vous ne l'avez pas
déclenché en marchant dessus, que vous n'avez pas vu quelqu'un d'autre le
déclencher ou que vous ne l'avez pas découvert avec
la commande `\texttt{s}' (fouiller) (plusieurs tentatives sont souvent nécessaires ;
si votre chance est faible, il peut en falloir beaucoup).
Les \textit{baguettes de détection des portes secrètes} et le sort de \textit{détection de l'invisible}
révèlent aussi les pièges dans un rayon modeste, mais seulement si le piège se trouve
aussi dans votre ligne de vue (que vous puissiez voir à ce moment-là ou non).
D'autres magies peuvent également révéler les pièges.

%.pg
Les monstres peuvent eux aussi être victimes des
pièges, ce qui peut éventuellement servir de stratégie défensive.
Malheureusement, les pièges peuvent aussi blesser votre ou vos animaux de compagnie.
Les monstres, animaux de compagnie compris, évitent généralement de marcher sur
un piège affiché sur votre carte s'ils ont déjà rencontré ce type de piège.

%.pg
Certains pièges, comme les fosses, les pièges à ours et les toiles d'araignée,
vous immobilisent.
Vous pouvez vous échapper en essayant simplement de vous déplacer vers une case
adjacente, autant de fois que nécessaire ; vous finirez par vous libérer.

%.pg
D'autres pièges peuvent vous envoyer ailleurs.
Les pièges de téléportation vous envoient ailleurs sur le même niveau du donjon.
Les téléporteurs de niveau vous envoient sur un niveau aléatoire du donjon, la
destination étant choisie entre quelques niveaux plus bas et tout en haut.
Ces pièges choisissent une nouvelle destination à chaque déclenchement.
Les trappes et les trous vous envoient aussi sur un autre niveau, mais toujours
situé sous le niveau actuel.
Ce sera généralement le niveau juste en dessous, mais cela peut être plus loin.
Contrairement aux téléporteurs (de niveau), le niveau de destination d'une trappe
ou d'un trou donné est fixe : y tomber vous amène chaque fois au même niveau---mais
pas forcément au même endroit de ce niveau.
Les portails magiques se comportent de manière similaire, avec quelques variantes.
Certains portails fonctionnent dans les deux sens et leur destination est
toujours la même : un autre portail qui peut vous ramener.
D'autres sont à sens unique et vous envoient sur un niveau précis, mais pas
forcément à un endroit précis de ce niveau.

%.pg
Il existe une branche spéciale du donjon, sur plusieurs niveaux, dont les plans
sont prédéfinis d'après le jeu vidéo classique ``\textit{Sokoban}''.
Dans ce jeu, vous incarnez un magasinier qui pousse des caisses autour
d'obstacles pour les placer à des emplacements désignés.
Dans NetHack, le but est de pousser des rochers dans des fosses ou des trous
jusqu'à ce que tous ces pièges soient neutralisés, ce qui donne accès à ce qui
se trouve au-delà.
Dans le jeu Sokoban, vous ne pouvez vous déplacer que dans les quatre directions
cardinales, et une caisse arrivée à destination bloque ensuite l'accès
à cette case.
Dans les niveaux Sokoban de NetHack, vous pouvez vous déplacer en diagonale (sauf
si cela vous fait passer entre deux rochers voisins), mais vous ne pouvez
pousser les rochers que dans les quatre directions cardinales, et un rocher qui
comble une fosse ou un trou fait disparaître à la fois le rocher et le piège,
ce qui rend la case normalement accessible.
Avec un peu de prévoyance, il est possible de terminer tous les niveaux
en respectant les règles traditionnelles de Sokoban.
(Astuce : pour résoudre les énigmes de Sokoban, il faut souvent éloigner les objets
de leur destination finale afin de se ménager plus de place pour manœuvrer.)
Comme NetHack ne propose pas de fonction \textit{annuler}, quelques entorses
sont permises au cas où vous seriez coincé.
Par exemple, chaque niveau comporte au moins un rocher supplémentaire.
Il est aussi possible de tout lâcher pour pouvoir vous faufiler
sur la même case qu'un rocher (et ensuite, on l'espère, le dépasser),
de détruire un rocher par la magie ou avec des outils, ou de créer de nouveaux
rochers avec un \textit{parchemin de terre}.
Cependant, de telles actions font baisser votre chance sans qu'aucun message
particulier ne vous le signale.
Consultez la section \textit{Conduite} pour savoir comment obtenir un retour sur
vos actions dans Sokoban.

%.hn 2
\subsection*{Escaliers et échelles (`\texttt{<}', `\texttt{>}')}

%.pg
En général, chaque niveau du donjon possède un escalier montant
(`\texttt{<}') vers le niveau précédent et un escalier descendant (`\texttt{>}')
vers le niveau
suivant.  Il y a toutefois quelques exceptions.  Par exemple, assez tôt
dans le donjon, vous trouverez un niveau comportant deux escaliers descendants,
l'un continuant dans le donjon et l'autre menant à une zone
appelée les Mines des Gnomes.  Ces mines finissent en cul-de-sac ;
après les avoir explorées (si vous choisissez de le faire), vous devrez donc
remonter jusqu'au donjon principal.

%.pg
Lorsque vous empruntez un escalier ou déclenchez un piège qui vous envoie
sur un autre niveau, le niveau que vous quittez est désactivé et
enregistré dans un fichier sur le disque.  Si vous vous rendez sur un niveau
déjà visité, il est chargé depuis son fichier sur le disque et réactivé.  Si
vous vous rendez sur un niveau encore jamais visité, il est
créé (de toutes pièces pour la plupart des niveaux aléatoires, à partir d'un
modèle pour certains niveaux ``spéciaux'', ou à partir des restes d'une
partie précédente pour un niveau ``d'ossements'', comme décrit brièvement plus bas).
Les monstres ne sont actifs que sur le niveau actuel ; ceux des autres niveaux
sont en quelque sorte mis en sommeil.

%.pg
D'ordinaire, lorsque vous empruntez un escalier, vous arrivez sur
l'escalier correspondant à destination.  Cependant, les animaux de compagnie
(voir plus bas) et certains autres monstres vous suivent s'ils sont assez près
quand vous montez ou descendez, et il arrive que l'une de ces créatures
vous déplace pendant le trajet.  Dans ce cas, l'animal de compagnie ou l'autre
monstre arrive sur l'escalier et vous vous retrouvez à proximité.

%.pg
Les échelles ont la même fonction que les escaliers, et ces deux types de
liaisons entre niveaux sont presque indiscernables en cours de partie.

%.hn 2
\subsection*{Boutiques et achats}

%.pg
Il vous arrivera de tomber sur une pièce avec un marchand près de la porte
et de nombreux objets posés par terre.  Vous pouvez acheter des objets en les
ramassant puis en utilisant la commande `\texttt{p}'.  Vous pouvez vous renseigner
sur le prix d'un objet avant de le ramasser en utilisant la commande ``\texttt{\#chat}''
lorsque vous vous tenez dessus.  Utiliser un objet avant de l'avoir payé entraîne des
frais, et le marchand ne vous laissera pas quitter la boutique tant que vous
n'aurez pas réglé toutes vos dettes.

%.pg
Vous pouvez vendre des objets à un marchand en les lâchant par terre à
l'intérieur de sa boutique.  Soit on vous propose une somme en or en vous demandant
si vous acceptez de vendre, soit on vous dit que le marchand
n'est pas intéressé (en général, votre objet doit correspondre au
type de marchandise vendu par la boutique).

%.pg
Si vous lâchez quelque chose dans une boutique par accident, le marchand en
revendique généralement la propriété sans offrir de compensation.  Vous devrez
le racheter si vous voulez le récupérer.

%.pg
Il arrive que les marchands manquent d'argent.
Dans ce cas, on vous
propose un avoir au lieu de l'or lorsque vous essayez de vendre quelque chose.
L'avoir peut servir à payer des achats, mais il n'est valable que dans la boutique
où il a été obtenu ; les autres marchands ne l'honoreront pas.
(S'il vous arrive de
trouver une ``carte de crédit'' dans le donjon, ne vous donnez pas la peine
d'essayer de l'utiliser dans les boutiques ; les marchands ne l'acceptent pas.)

%.pg
La commande \texttt{\$}, qui indique la quantité d'or que vous portez,
affiche aussi votre dette ou votre avoir actuel en boutique, le cas échéant.
La commande \texttt{Iu} liste les objets impayés (ceux qui appartiennent encore à la
boutique) si vous en portez.
La commande \texttt{Ix} affiche, sous une forme proche de l'inventaire, les objets impayés
qui ont été consommés, ainsi que les autres frais de boutique éventuels.

%.hn 3
\subsubsection*{Particularités des boutiques}

%.pg
Plusieurs aspects du comportement des boutiques peuvent surprendre.

\begin{itemize}
% note: a bullet is the default item label so we could omit [$\bullet$] here
%.lp \(bu 2
\item
Le prix d'un même objet peut varier en fonction de nombreux facteurs.
%.lp \(bu 2
\item
Un marchand considère la case située juste à l'intérieur de la porte comme si
elle était hors de la boutique.
%.lp \(bu 2
\item
Pendant que le marchand vous surveille comme un faucon, il ou elle ignore
généralement tous les autres clients.
%.lp \(bu 2
\item
Si une boutique est ``fermée pour inventaire'', elle ne rouvrira pas d'elle-même.
%.lp \(bu 2
\item
Les boutiques ne sont jamais réapprovisionnées, même si leur stock s'épuise.
\end{itemize}

%.hn 2
\subsection*{Informations sur les déplacements}

%.pg
Se déplacer sur la carte ne fournit généralement aucune
information---hormis le dessin du héros à sa nouvelle position---sauf
si vous marchez sur un objet ou un tas d'objets,
ou sur un piège, ou si vous tentez d'aller sur une case occupée par un monstre.
Plusieurs options permettent d'enrichir ces informations.

%.pg
L'option
\textit{pile\textunderscore limit}
détermine combien d'objets doivent former un
tas---partageant la même case de la carte---pour
que le jeu annonce ``Il y a plusieurs objets ici.'' au lieu de les énumérer.
La valeur par défaut est \texttt{5}.
La mettre à \texttt{1} donnerait toujours ce message au lieu d'énumérer
les objets.
La mettre à \texttt{0} est un cas particulier qui énumère toujours tous
les objets, quelle que soit la taille du tas.
Notez que ce nombre correspond au nombre de piles d'objets distinctes
présentes et non à la somme des quantités de ces piles (ainsi,
\texttt{7 flèches} ou \texttt{25 pièces d'or} comptent chacune pour 1 et non
pour 7 et 25 respectivement, et totalisent 2 lorsqu'elles se trouvent sur la
même case).

%.pg
Le préfixe de commande \texttt{nopickup} (par défaut `\texttt{m}') peut
être utilisé avant une direction de déplacement pour marcher sur des objets sans
tenter de les ramasser automatiquement et sans recevoir d'information à leur sujet.

%.pg
L'option
\textit{mention\textunderscore walls}
détermine si vous recevez une information lorsque vous tentez de foncer dans un mur
ou dans la roche pleine, ou de sortir du bord de la carte.
Normalement, rien ne se passe (sauf si le héros est aveugle et qu'aucun mur n'est
affiché ; le mur heurté est alors dessiné sur la carte).
Cette option donne aussi une information lorsqu'une course ou un déplacement
rapide s'interrompt pour une raison peu évidente.

%.pg
L'option
\textit{mention\textunderscore decor}
détermine si vous recevez une information lorsque vous marchez sur du ``mobilier''.
Normalement, marcher sur un escalier, une fontaine, un autel ou divers autres
éléments ne provoque rien, sauf s'ils sont recouverts par un ou plusieurs objets
et donc masqués sur la carte.
Activer cette option décrit ces éléments même lorsqu'ils
ne sont pas masqués.
Les embrasures sans porte et les portes ouvertes ne sont pas jugées dignes d'être
mentionnées ; les portes fermées (si vous pouvez aller sur leur case) et les
portes cassées le sont.
À supposer que vous puissiez le faire, aller sur de l'eau, de la lave ou de la glace
donne une information si vous n'étiez pas déjà sur ce type de terrain, mais ne la
répète pas (sauf si un autre message s'est intercalé) lorsque vous passez d'une case
d'eau à une autre, de lave à lave ou de glace à glace.
Quitter l'un de ces terrains pour revenir sur un terrain ``normal'' donne aussi
un message, sauf s'il y a une information sur un ou plusieurs objets, auquel
cas le retour sur la terre ferme est implicite.

%.pg
Les options
\textit{confirm}
et
\textit{safe\textunderscore pet}
déterminent ce qui se passe lorsque vous tentez d'aller sur la case d'un monstre
pacifique ou d'un monstre apprivoisé.

%.\" getting away from "Movement feedback" here; oh well...
%.pg
Le préfixe de commande \texttt{nopickup} (par défaut `\texttt{m}') sert
aussi de préfixe de déplacement sans attaque et permet d'essayer d'aller
sur la case d'un monstre visible sans que le déplacement soit considéré comme
une attaque (voir la sous-section \textit{Combat} de \textit{Monstres} plus bas).
Le préfixe de commande `\texttt{fight}' (par défaut `\texttt{F}' ;
aussi `\texttt{-}' si
\textit{number\textunderscore pad}
est activée) permet de forcer une attaque, lorsque vous devinez où se trouve
un monstre invisible ou que vous attaquez délibérément une créature pacifique
ou apprivoisée.

%.pg
L'option
\textit{run\textunderscore mode}
détermine la fréquence à laquelle la carte est redessinée lorsque vous parcourez
plus d'une case en une seule commande (donc en courant, en vous précipitant ou en voyageant).

%.hn 2
\subsection*{Niveau Rogue}

%.pg
Un niveau du donjon (situé vers le milieu ou la fin de la quinzaine du donjon principal)
rend hommage à \textit{rogue}, l'inspiration du jeu ancêtre \textit{hack}.

%.pg
Il est généralement affiché différemment des autres niveaux : éventuellement en
caractères au lieu de tuiles, ou sans symboles de tracé de lignes si l'affichage
est déjà en caractères ; de plus, l'or est représenté par \texttt{*} au lieu de \texttt{\textdollar}
et les escaliers par \texttt{\%} au lieu de \texttt{<} et \texttt{>}.
Le jeu lui-même présente quelques différences mineures : les embrasures n'ont pas
de porte ; un parchemin, une baguette ou un sort de lumière utilisé dans une pièce
éclaire toute la pièce au lieu d'un rayon autour de votre personnage.
Et les monstres représentés par des lettres minuscules ne sont pas générés
aléatoirement sur le niveau Rogue.

%.pg
La légère étrangeté de ce niveau est une fonctionnalité, pas un bogue...

%.hn 1
\section{Monstres}

%.pg
Les monstres que vous ne pouvez pas voir ne sont pas affichés à l'écran.  Prudence !
Vous pouvez soudain tomber sur l'un d'eux dans un endroit sombre.  Certains objets
magiques peuvent vous aider à les localiser avant qu'ils ne vous localisent (ce que
certains monstres font très bien).

%.pg
Les commandes `\texttt{/}' et `\texttt{;}' permettent d'obtenir des informations
sur les
monstres affichés à l'écran.  La commande ``\texttt{\#name}''
(associée par défaut à `\texttt{C}') vous permet
de donner un nom à un monstre, ce qui peut aider à les distinguer
les uns des autres lorsque plusieurs monstres sont présents.  Donner un nom
constitué d'une simple espace supprime le nom précédent.

%.pg
La commande étendue ``\texttt{\#chat}'' permet d'interagir avec un monstre
adjacent.  Il n'y a pas de véritable dialogue (autrement dit, vous ne choisissez pas
ce que vous dites), mais discuter avec certains monstres, comme un
marchand ou l'Oracle de Delphes, peut donner des résultats utiles.

%.hn 2
\subsection*{Combat}

%.pg
Si vous voyez un monstre et souhaitez le combattre, essayez simplement de
marcher sur lui.  Beaucoup de monstres que vous rencontrerez s'occuperont de leurs
affaires, sauf si vous les attaquez.  Certains sont très dangereux lorsqu'ils sont en colère.
Souvenez-vous : la prudence est mère de la bravoure.

%.pg
Dans la plupart des cas, si vous tentez d'attaquer un monstre pacifique en
allant sur sa case, on vous demande de confirmer votre intention.  Par
défaut, répondre `\texttt{y}' confirme cette intention,
ce qui peut prêter à erreur si vous utilisez `\texttt{y}' pour vous déplacer.  Vous pouvez régler
l'option \textit{paranoid\textunderscore confirmation:attack}
pour exiger plutôt la réponse ``\texttt{yes}''.
%.pg

Si vous ne pouvez pas voir un monstre (s'il est invisible ou si vous êtes aveugle),
le symbole `\texttt{I}' est affiché lorsque vous apprenez sa présence.
Si vous tentez de marcher sur lui, vous essaierez de le combattre comme
un monstre visible ; bien sûr,
si le monstre s'est déplacé, vous frapperez dans le vide.  Si vous devinez
que le monstre s'est déplacé et ne souhaitez pas combattre, vous pouvez utiliser la
commande `\texttt{m}' pour vous déplacer sans combattre ; de même, si vous ne vous
souvenez d'aucun monstre mais voulez quand même essayer de combattre, vous pouvez
utiliser la commande `\texttt{F}'.

%.hn 2
\subsection*{Votre animal de compagnie}

%.pg
Vous commencez la partie avec un petit chien (`\texttt{d}'), un chaton (`\texttt{f}')
ou un poney (`\texttt{u}'), qui vous suit
dans le donjon et combat les monstres à vos côtés.
Comme vous, votre animal de compagnie a besoin de nourriture pour survivre.
Les chiens et les chats se nourrissent généralement de charognes fraîches et
d'autres viandes ; les chevaux ont besoin d'une nourriture végétarienne, plus
difficile à trouver.
Si vous vous inquiétez pour votre animal ou voulez le dresser, vous
pouvez aussi le nourrir en lui lançant de la nourriture.
Un animal bien dressé peut être très utile dans certaines circonstances.

%.pg
Votre animal de compagnie gagne aussi de l'expérience en tuant des monstres et
peut grandir avec le temps, gagnant des points de vie et infligeant plus de dégâts.
Au début, votre animal peut même être meilleur que vous pour tuer, ce qui rend les
animaux de compagnie utiles aux personnages de bas niveau.

%.pg
Votre animal de compagnie vous suit dans les escaliers, en montée comme en
descente, s'il est à côté de vous au moment où vous partez.  Sinon, il est
abandonné et peut redevenir sauvage.  De même, lorsque vous déclenchez certains
types de pièges qui changent votre position (par exemple une trappe qui vous fait
tomber à un niveau inférieur du donjon), tout animal adjacent vous accompagne et
tout animal non adjacent reste en arrière.  Votre animal peut lui-même déclencher
de tels pièges ; vous ne serez pas emporté avec lui, même si vous êtes adjacent
à ce moment-là.

%.hn 2
\subsection*{Montures}

%.pg
Certains types de créatures du donjon peuvent effectivement être montés si vous
avez l'équipement et la compétence nécessaires.  Convaincre une bête sauvage de se
laisser seller est pour le moins difficile.  Plus d'un explorateur de donjon
a dû recourir à la magie et à la sorcellerie pour sceller l'alliance.
Une fois la bête sous votre contrôle, en revanche, vous pouvez
facilement monter en selle et en descendre avec la commande ``\texttt{\#ride}''.  Une
fois en selle, guidez la bête dans le donjon de la même manière
que vous vous déplaceriez vous-même.  C'est la bête que vous verrez affichée
sur la carte.

%.pg
La compétence d'équitation se gère avec la commande ``\texttt{\#enhance}''.  Consultez la section
sur la maîtrise des armes pour plus d'informations à ce sujet.

%.pg
Utilisez la commande `\texttt{a}' (appliquer) et choisissez une selle dans votre inventaire
pour tenter de la poser sur une créature adjacente.  En cas de réussite,
elle est transférée dans l'inventaire de cette créature.

%.pg
Utilisez la commande ``\texttt{\#loot}'' à côté d'une créature sellée pour
tenter de lui retirer sa selle.  En cas de réussite, elle est
transférée dans votre inventaire.

%.hn 2
\subsection*{Niveaux d'ossements}

%.pg
Vous pouvez rencontrer les spectres et les cadavres d'autres aventuriers (ou même
d'anciennes incarnations de vous-même !) ainsi que leurs effets personnels.  Les
fantômes sont difficiles à tuer, mais faciles à éviter, car ils sont lents et
infligent peu de dégâts.  Vous pouvez piller les possessions de l'aventurier défunt ;
cependant, elles sont probablement maudites.  Méfiez-vous de ce qui a tué
l'ancien joueur ; il rôde sans doute encore dans les parages, savourant sa
dernière victoire.

%.hn 2
\subsection*{Persistance des monstres}

%.pg
Les monstres (terme générique qui inclut aussi les humains et les animaux de
compagnie) ne sont affichés que lorsqu'on peut les voir ou les percevoir d'une
autre manière.
Aller à un endroit d'où vous ne pouvez plus voir ou percevoir un monstre
le fait disparaître de votre carte ; de même si c'est lui
qui s'est déplacé plutôt que vous.

%.pg
Cependant, si vous rencontrez un monstre que vous ne pouvez ni voir ni
percevoir---peut-être est-il invisible et vient-il de vous taper sur la caboche---un
marqueur spécial de ``monstre invisible mémorisé'' est affiché à
l'endroit où vous pensez qu'il se trouve.
Ce marqueur persiste jusqu'à ce que vous ayez
prouvé qu'il n'y a pas de monstre à cet endroit, même si le monstre invisible
se déplace ailleurs ou si vous allez à un endroit d'où l'emplacement du
marqueur ne serait normalement plus visible.

%.hn 1
\section{Objets}

%.pg
Lorsque vous trouvez quelque chose dans le donjon, il est courant de vouloir le
ramasser.  Dans \textit{NetHack}, cela se fait avec la commande `\texttt{,}'.
Si l'option \textit{autopickup} est activée, vous ramasserez automatiquement l'objet
en marchant dessus, sauf si vous vous déplacez avec le préfixe `\texttt{m}'.
%.pg
Si vous portez trop d'objets, \textit{NetHack} vous le dira et vous ne
pourrez plus rien ramasser.  Sinon, il ajoutera le ou les objets
à votre sac et vous indiquera ce que vous venez de ramasser.
%.pg
Chaque objet ajouté à votre inventaire ajoute aussi son poids
à votre charge.  La quantité que vous pouvez porter dépend de votre force et de
votre constitution.  Plus
vous êtes fort et robuste,
moins la charge supplémentaire vous affecte.  Il arrive cependant
un moment où le poids de tout ce bazar que vous transportez
avec vous dans le donjon vous encombre.  Vos réactions
ralentissent et vous brûlez des calories plus vite, ce qui vous oblige à manger
plus souvent pour compenser.  Finalement, vous serez si surchargé qu'il vous faudra
soit vous débarrasser d'une partie de ce que vous portez, soit vous effondrer sous son poids.
%.pg
\textit{NetHack} vous indique à quel point vous êtes chargé.
Si vous êtes encombré, l'un des états
\texttt{Chargé}, \texttt{Stressé}, \texttt{Tendu},
\texttt{Exténué} ou \texttt{Surchargé} est
affiché sur la ligne d'état en bas de l'écran.

%.pg
Quand vous ramassez un objet, une lettre d'inventaire lui est attribuée.  Beaucoup
de commandes qui agissent sur des objets doivent vous demander quel objet
vous voulez utiliser.  Lorsque \textit{NetHack} vous demande de choisir un objet précis
que vous portez, on vous présente généralement une liste de lettres
d'inventaire parmi lesquelles choisir (voir Commandes, plus haut).

%.pg
Certains objets, comme les armes, se distinguent facilement.  D'autres, comme
les parchemins et les potions, reçoivent des descriptions qui varient selon leur
type.  Au cours d'une partie, deux objets ayant la même description sont
du même type.  En revanche, les descriptions changent d'une partie à l'autre.

%.pg
Lorsque vous utilisez l'un de ces objets et que son effet est évident, \textit{NetHack}
se souviendra pour vous de ce qu'il est.  Si son effet n'est pas tout à fait
évident, on vous demandera comment vous voulez appeler ce type d'objet
afin de le reconnaître plus tard.  Vous pouvez aussi utiliser la commande ``\texttt{\#name}''
à tout moment, dans le même but, pour nommer
tous les objets d'un type donné ou seulement un objet individuel.
Lorsque vous utilisez ``\texttt{\#name}'' sur un objet déjà nommé,
indiquer une espace comme valeur supprime le nom précédent au lieu
d'en attribuer un nouveau.

%.hn 2
\subsection*{Malédictions et bénédictions}

%.pg
Tout objet que vous trouvez peut être maudit, même s'il est
par ailleurs utile.  L'effet le plus courant d'une malédiction est de vous
retrouver coincé avec (et collé à) l'objet.  Les armes maudites se soudent à votre
main lorsque vous les maniez, si bien que vous ne pouvez plus les lâcher.  Tout objet
maudit que vous portez ne peut pas être retiré par des moyens ordinaires.  De plus,
les armes et armures maudites portent généralement, mais pas toujours, des
enchantements négatifs qui les rendent moins efficaces au combat.  D'autres objets
maudits peuvent mal fonctionner ou se révéler nuisibles d'autres manières.

%.pg
Les objets peuvent aussi, à l'inverse, être bénis.
Les objets bénis fonctionnent généralement mieux ou de façon plus bénéfique que
les objets non maudits normaux.
Par exemple, une arme bénie inflige un peu plus de dégâts aux démons.

%.pg
Les objets qui ne sont ni maudits ni bénis sont dits non maudits.
On aurait tout aussi bien pu les qualifier de non bénis, mais c'est
la désignation non maudit que vous verrez dans le jeu.  Une situation de ``verre
à moitié plein ou à moitié vide'' ; à vous d'en penser ce que vous voulez.

%.pg
Il existe des moyens magiques de jeter ou de lever des malédictions sur les objets ;
même si vous êtes coincé avec l'un d'eux, vous pouvez donc encore faire lever
la malédiction et retirer l'objet.
Les Prêtres et les Prêtresses ont une
sensibilité innée à cette propriété chez tout objet ; ils peuvent donc éviter les
objets maudits plus facilement que les autres rôles.
Lâcher des objets sur un autel révèle leur état de bénédiction ou de malédiction,
à condition que vous puissiez les voir atterrir.

%.pg
Un objet dont l'état est inconnu apparaît dans votre inventaire sans préfixe.
Un objet dont vous connaissez l'état se distingue dans votre inventaire
par la présence du mot \texttt{maudit}, \texttt{non maudit} ou
\texttt{béni} dans sa description.
Dans certains cas, \texttt{non maudit} est omis car redondant lorsque
suffisamment d'autres informations sont affichées.
L'option
\textit{implicit\textunderscore uncursed}
permet de régler cela ; désactivez-la pour que \texttt{non maudit}
soit affiché même lorsque cela peut se déduire d'autres attributs.

%.pg
L'état de bénédiction ou de malédiction des objets est parfois appelé leur
attribut ``\texttt{BUC}'', pour l'état béni (Blessed), non maudit (Uncursed)
ou maudit (Cursed), ou ``\texttt{BUCX}'' pour béni, non maudit, maudit ou inconnu.
(Le terme \textit{béatitude} est aussi parfois employé.)

%.hn 2
\subsection*{Artefacts}

%.pg
Certains objets ont été dotés de pouvoirs spéciaux et sont appelés
\textit{Artefacts}.
Ils ont un type précis (comme épée longue ou dague orque) et un nom
distinct, comme \textit{Tueuse-de-géants} ou \textit{Dent-Cruelle}.
Les armes artefacts infligent généralement plus de dégâts que leurs équivalents ordinaires.
Certaines infligent des dégâts supplémentaires à tous les monstres, d'autres seulement
à certains types de monstres et ne valent donc pas mieux que des armes ordinaires
contre les autres.
Certaines confèrent des capacités défensives lorsqu'on les manie ou ont d'autres
pouvoirs qui ne sont pas décrits ici.

%.pg
Vous pouvez les trouver simplement posés par terre, y compris (mais pas
seulement) dans les boutiques, ou en recevoir un en récompense d'un ``\texttt{\#offer}''
sur un autel de votre divinité tutélaire.
Quelques-uns peuvent être lâchés par des monstres, ou obtenus à partir d'un objet
ordinaire du même type en lui donnant le bon nom.
% should we mention dipping for Excalibur here?
Vous pouvez aussi en souhaiter un, s'il vous arrive d'obtenir un vœu, mais de tels
vœux peuvent échouer.

%.pg
Certains artefacts ont un alignement précis, d'autres non.
Vous n'obtiendrez pas par des offrandes d'artefacts alignés d'un alignement
différent du vôtre, et vous risquez de recevoir une décharge si vous tentez d'en
souhaiter un ou si vous en trouvez un et tentez de l'utiliser.

%pg
Chaque rôle possède un artefact particulier qui se trouve dans la branche
\textit{Quête} du donjon.
On les appelle couramment les artefacts de quête.
Tous sont alignés et la plupart ne sont pas des armes.
On ne les trouve pas au hasard.

%.pg
Les commandes `\texttt{\textbackslash}' et `\texttt{\textasciigrave{}a}'
listent les artefacts que vous avez entièrement identifiés (connaître le nom et
le type d'objet ne suffit pas).

%.hn 2
\subsection*{Reliques}

%.pg
Il existe trois objets uniques qui portent un nom et possèdent des pouvoirs
spéciaux limités, mais qui ne sont pas classés parmi les artefacts.
Chacun est gardé par un monstre particulier et vous devrez réunir les
trois pour vous en servir en fin de partie.
% The relics are listed in the same order as the Oracle's message about
% them rather than in the order they need to be used for the invocation.
Il s'agit de \textit{la Cloche d'Ouverture},
\textit{du Livre des Morts} et
\textit{du Candélabre d'Invocation}.
Leurs descriptions respectives tant qu'ils ne sont pas identifiés sont
cloche d'argent, livre de sorts en papyrus et candélabre.

%.hn 2
\subsection*{Armes (`\texttt{)}')}

%.pg
Si on leur en laisse l'occasion, la plupart des monstres des Mazes of
Menace essaieront gratuitement de vous tuer.
Il vous faut des armes pour vous défendre (en les tuant les premiers).
Sans
arme, vous n'infligez que 1--2 points de vie de dégâts (plus les bonus éventuels).
Les Moines font exception ; ils infligent normalement plus de dégâts à mains
nues (ou gantées) qu'avec des armes.

%.pg
Il y a des armes de mêlée, comme les masses d'armes et les épées, et des armes
de jet, comme les flèches et les lances.  Pour frapper des monstres avec une arme,
vous devez la manier et les attaquer, ou la leur lancer.  Vous pouvez tout
simplement choisir de lancer une lance.
Pour tirer une flèche, vous devez d'abord manier un arc, puis lancer la flèche.
Les arbalètes tirent des carreaux d'arbalète.  Les frondes projettent des cailloux
et (d'autres) pierres (comme des gemmes).

%.pg
Les armes enchantées ont un ``plus'' (ou ``bonus au toucher'', qui peut être
positif ou négatif) qui s'ajoute à vos chances de
toucher et aux dégâts que vous infligez à un monstre.  Le seul moyen de connaître
l'enchantement d'une arme est de la faire identifier magiquement d'une façon ou d'une autre.
La plupart des armes sont sujettes à un type de dégradation comme la rouille.  Ces
dégâts d'``érosion'' peuvent être réparés.

%.pg
La probabilité qu'une attaque touche un monstre, et la quantité
de dégâts qu'infligera ce coup, dépendent de nombreux facteurs.  Parmi eux :
le type d'arme, la qualité de l'arme (enchantement et/ou érosion), le niveau
d'expérience, la force, la dextérité, l'encombrement et la maîtrise (voir plus bas).  La
classe d'armure
du monstre---un indice de défense général, qui ne vient pas forcément du port d'une armure---est
aussi un facteur ; en outre, certains monstres sont particulièrement
vulnérables à certains types d'armes.

%.pg
Beaucoup d'armes peuvent être maniées à une main ; certaines exigent les deux mains.
Lorsque vous maniez une arme à deux mains, vous ne pouvez pas porter de bouclier, et
inversement.  Lorsque vous maniez une arme à une main, vous pouvez garder une autre
arme prête à servir en l'organisant avec la commande `\texttt{x}', qui
échange votre arme principale (celle que vous maniez) et votre arme secondaire.
Et si vous maîtrisez la compétence ``combat à deux armes'', vous
pouvez manier les deux armes simultanément, comme arme principale et secondaire ;
utilisez la commande `\texttt{X}' pour activer ou désactiver ce mode.
Seuls certains types de personnages (les barbares, par exemple) disposent de la
compétence nécessaire.  Même avec cette compétence, utiliser deux armes à la fois entraîne
une pénalité aux chances de toucher votre cible par rapport à l'utilisation d'une
seule arme à la fois.

%.pg
Il se peut que vous préfériez parfois ne manier aucune arme.
Pour cela, maniez `\texttt{-}', ou utilisez la commande `\texttt{A}' qui
vous permet de lâcher l'arme maniée en plus de retirer
d'autres objets portés.

%.pg
Ceux d'entre vous, dans l'assistance, qui jouent à AD\&D doivent savoir que chaque
arme qui existait dans AD\&D inflige à peu près les mêmes dégâts aux monstres dans
\textit{NetHack}.  Certaines des armes les plus obscures (comme
l'\textit{aklys}, le \textit{marteau de Lucerne} et le \textit{bec-de-corbin}) sont définies
dans une annexe d'\textit{Unearthed Arcana}, un supplément d'AD\&D.

%.pg
Certaines interfaces prennent en charge l'option \texttt{weaponstatus}.
Lorsqu'elle est activée, un état supplémentaire est affiché, décrivant
l'arme actuellement maniée.

%.pg
Les commandes liées aux armes sont `\texttt{w}' (manier), `\texttt{t}' (lancer),
`\texttt{f}' (tirer), `\texttt{Q}' (carquois),
`\texttt{x}' (échanger), `\texttt{X}' (deux armes) et ``\texttt{\#enhance}''
(voir plus bas).

%.hn 3
\subsection*{Lancer et tirer}

%.pg
Vous pouvez lancer à peu près n'importe quoi avec la commande `\texttt{t}'.  Elle vous
demande l'objet à lancer ; choisir `\texttt{?}' liste les objets de votre inventaire
qui sont susceptibles d'être lancés, et choisir `\texttt{*}' liste
tout votre inventaire.  Une fois l'objet choisi, on vous
demande une direction plutôt qu'une cible précise.  La
distance à laquelle un objet peut être lancé dépend surtout du type d'objet
et de votre force.  Les flèches peuvent être lancées à la main, mais elles vont
beaucoup plus loin et ont plus de chances de toucher lorsqu'elles sont tirées
pendant que vous maniez un arc.

%.pg
Certaines armes reviennent une fois lancées.
Un boomerang---pourvu qu'il ne touche rien---en est un exemple évident.
Si un aklys (gourdin à lanière) est lancé alors qu'il est manié, il revient
même s'il touche quelque chose.
Un héros suffisamment fort peut lancer le marteau de guerre \textit{Mjollnir} ;
lancé par une \textit{Valkyrie}, il revient lui aussi.
Cependant, les aklys et \textit{Mjollnir} ne reviennent pas toujours.
Les objets qui reviennent ne sont pas toujours rattrapés, et frappent même parfois
le lanceur, mais une fois rattrapés ils sont de nouveau maniés.

%.pg
Vous pouvez simplifier le lancer en utilisant la commande `\texttt{Q}' pour
choisir votre ``projectile'' préféré, puis la commande `\texttt{f}' pour
le lancer.  On vous demandera une direction comme ci-dessus, mais vous n'aurez pas
à préciser l'objet à lancer chaque fois que vous utiliserez `\texttt{f}'.  Il existe
aussi une option,
\textit{autoquiver},
grâce à laquelle \textit{NetHack} choisit un autre objet pour remplir automatiquement votre
carquois (ou sacoche à projectiles, ou objet tenu prêt) lorsque l'emplacement
d'inventaire utilisé pour `\texttt{Q}' est épuisé.
Si votre carquois est vide, que \textit{autoquiver}
est désactivée et que vous maniez une arme qui revient une fois lancée,
c'est cette arme que vous lancerez au lieu de remplir le carquois.
La commande de tir offre aussi une aide supplémentaire : si \textit{fireassist}
est activée, elle essaiera de manier un lanceur adapté aux munitions du carquois.

%.pg
Certains personnages sont capables de lancer ou de tirer une salve de plusieurs
objets (d'une même pile) en une seule action.
Savoir charger plusieurs munitions à
la fois---ou tenir plusieurs projectiles en main---et
toucher quand même une cible n'est pas chose facile.
Les Rôdeurs font partie de ceux qui sont experts
dans cet exercice, de même que ceux qui ont un haut niveau de maîtrise dans la
compétence d'arme correspondante (la compétence d'arc si vous en maniez un pour
tirer des flèches, la compétence d'arbalète si vous en maniez une pour tirer des carreaux,
ou la compétence de fronde si vous en maniez une pour tirer des pierres).
Le nombre d'objets que le personnage a une chance de tirer varie d'un
tour à l'autre.  Vous pouvez limiter explicitement le nombre de tirs en utilisant un
préfixe numérique avant la commande `\texttt{t}' ou `\texttt{f}'.
Par exemple, ``\texttt{2f}'' (ou ``\texttt{n2f}'' en mode
\textit{number\textunderscore pad})
garantit qu'au plus 2 flèches sont tirées,
même si vous auriez pu en tirer 3.  Si vous indiquez
un nombre plus grand que celui qui aurait été tiré (``\texttt{4f}'' dans cet exemple),
vous tirerez simplement le même nombre (3, ici) que si aucune limite
n'avait été indiquée.  Une fois la salve partie, tous les objets
voyagent dans la même direction ; si les premiers tuent un monstre,
les suivants peuvent continuer au-delà de cette case.

%.hn 3
\subsection*{Maîtrise des armes}

%.pg
Vous aurez divers degrés de compétence dans les armes disponibles.
La maîtrise des armes, ou compétences d'armes, détermine votre aptitude à utiliser
des types d'armes particuliers, et vous pourrez améliorer vos compétences
au fil de la partie, en fonction de votre rôle, de votre niveau
d'expérience et de votre usage des armes.

%.pg
Pour la maîtrise, les armes ont
été réparties en différents groupes, comme les dagues, les épées larges et les
armes d'hast.  Chaque rôle a une limite au niveau de maîtrise qu'un personnage
peut atteindre dans chaque groupe.  Par exemple, les magiciens peuvent devenir
très compétents avec les dagues ou les bâtons, mais pas avec les épées ou les arcs.

%.pg
La commande étendue ``\texttt{\#enhance}'' permet de consulter votre maîtrise
actuelle des armes
(ainsi que des sorts) et de choisir quelle(s) compétence(s) améliorer lorsque
vous avez suffisamment utilisé une ou plusieurs compétences pour y avoir droit.  Les
niveaux de compétence sont ``aucun'' (parfois aussi appelé ``restreint'',
car vous ne pourrez pas progresser), ``non formé'', ``débutant'', ``compétent''
et ``expert''.  Les compétences restreintes n'apparaissent tout simplement pas dans la liste
affichée par ``\texttt{\#enhance}''.
(Une intervention divine peut lever la restriction d'une
compétence ; elle commence alors au niveau non formé et est limitée à débutant.)
Certains personnages peuvent améliorer leur compétence de combat à mains nues ou
d'arts martiaux au-delà d'expert, jusqu'à ``maître'' ou ``grand maître''.

%.pg
Utiliser une arme dans laquelle vous êtes restreint ou non formé
entraîne une pénalité modeste aux chances de toucher un monstre ainsi qu'aux
dégâts infligés lorsque vous touchez ; au niveau débutant, il n'y a ni
pénalité ni bonus ; au niveau compétent, vous recevez un bonus modeste aux
chances de toucher et aux dégâts infligés ; au niveau expert, le bonus est
plus élevé.  Un coup réussi a une chance de faire progresser votre entraînement vers
le niveau de compétence suivant (sauf si vous avez déjà atteint la limite pour cette
compétence).  Lorsque cet entraînement atteint le seuil du niveau suivant,
on vous dit que vous vous sentez plus sûr de vos compétences.  Vous
pouvez alors utiliser ``\texttt{\#enhance}'' pour améliorer une ou plusieurs compétences.
Ces compétences
ne sont pas améliorées automatiquement, car le total de vos compétences
est limité ; vous devez donc choisir activement celles à améliorer
et celles à négliger.

%.hn 3
\subsection*{Combat à deux armes}

%.pg
Certains personnages peuvent utiliser deux armes à la fois.  Tout organiser pour
y parvenir peut sembler fastidieux, mais cela devient naturel avec l'habitude.
Pour manier deux armes, vous devez utiliser la commande ``\texttt{\#twoweapon}''.
Mais il faut d'abord avoir une arme dans chaque main.
(Notez que vos deux armes ne sont pas tout à fait égales ; celle qui est dans la
main avec laquelle vous maniez habituellement est considérée comme principale et
l'autre comme secondaire.
La différence la plus visible apparaît lorsque vous
cessez---ou avant que vous ne commenciez, d'ailleurs---de manier
deux armes à la fois.
L'arme principale est votre arme maniée et l'arme
secondaire n'est qu'un objet de votre inventaire désigné
comme arme de rechange.)

%.pg
Si votre arme principale est maniée mais que votre autre main est vide ou tient
la mauvaise arme, utilisez la séquence `\texttt{x}', `\texttt{w}', `\texttt{x}' pour
d'abord passer votre arme
principale dans votre autre main, manier ce que vous voulez comme arme
secondaire, puis remettre les deux dans les mains voulues.
Si votre arme secondaire (ou de rechange) est la bonne mais pas votre arme
principale, utilisez simplement `\texttt{w}' pour manier l'arme principale.
Enfin, si aucune main ne tient la bonne arme,
utilisez `\texttt{w}', `\texttt{x}', `\texttt{w}'
pour d'abord manier l'arme secondaire voulue, la passer dans l'autre main, puis
manier l'arme principale.

%.pg
Tout le processus peut être simplifié grâce à l'option
\textit{pushweapon}.
Lorsqu'elle est activée, utiliser `\texttt{w}' pour manier quelque chose
fait de l'arme actuellement maniée votre arme de rechange.
La séquence `\texttt{w}', `\texttt{w}' permet donc de manier d'abord l'arme que vous
destinez au rôle d'arme secondaire, puis celle que vous voulez comme principale,
ce qui repousse la première en position secondaire.

%.pg
En mode de combat à deux armes, la commande `\texttt{X}'
ramène au mode à une seule arme.
Lancer ou lâcher l'une des deux
armes, ou se la faire voler ou détruire, vous fait aussi
revenir au combat à une seule arme.

%.hn 2
\subsection*{Armures (`\texttt{[}')}

%.pg
Bien des choses hostiles rôdent alentour ; il vous faut une armure pour vous
protéger de leurs coups.  Certains types d'armure protègent mieux
que d'autres.  Votre classe d'armure mesure cette
protection.  La classe d'armure (CA) se mesure comme dans AD\&D : 10
correspond à l'absence d'armure et plus le nombre est bas, meilleure est l'armure.
Chaque armure qui existe dans AD\&D offre la même protection dans
\textit{NetHack}.

Voici la liste des valeurs de classe d'armure fournies par les armures de corps :

\begin{center}
\begin{tabular}{lllll}
cotte d'écailles de dragon     & 1 & \makebox[20mm]{}  & armure de plates               & 3\\
armure de plates en cristal    & 3 &                   & armure de plates en bronze     & 4\\
armure d'attelles              & 4 &                   & armure à bandes                & 4\\
cotte de mithril naine         & 4 &                   & cotte de mithril elfique       & 5\\
cotte de mailles               & 5 &                   & cotte de mailles orque         & 6\\
armure d'écailles              & 6 &                   & écailles de dragon             & 7\\
armure de cuir clouté          & 7 &                   & armure d'anneaux               & 7\\
armure d'anneaux orque         & 8 &                   & armure de cuir                 & 8\\
veste de cuir                  & 9 &                   & aucune armure                  & 10\\
\end{tabular}
\end{center}

%.pg
\noindent Vous pouvez aussi porter d'autres pièces d'armure (cape par-dessus l'armure de
corps, chemise en dessous, casque, gants, bottes, bouclier) pour abaisser encore
votre classe d'armure.
%--too obvious to mention unless we include polymorph into ettin or maralith
% You can wear at most one item of each category (one suit of armor, one
% cloak, one helmet, one shield, and so on) at a time.
La plupart d'entre elles améliorent la CA d'un ou deux points (rendant la
valeur globale plus petite, et à terme négative), mais elles peuvent aussi être
enchantées.
Les chemises font exception ; elles n'offrent aucune protection si elles ne sont pas enchantées.
Certaines capes n'améliorent pas non plus la CA lorsqu'elles ne sont pas enchantées, mais toutes
les capes offrent une certaine protection contre la rouille ou la corrosion à l'armure
de corps portée dessous, ainsi que contre certaines attaques de \textit{contact} des monstres.

%.pg
Si une pièce d'armure est enchantée, sa protection est meilleure
(ou pire) que la normale, et son ``plus'' (ou moins) se soustrait de
votre classe d'armure.  Par exemple, une cotte de mailles +1 vous donnerait
une meilleure protection qu'une cotte de mailles normale, abaissant votre classe
d'armure d'une unité supplémentaire, à 4.  Lorsque vous enfilez une pièce d'armure,
vous apprenez immédiatement sa classe d'armure et les éventuels ``plus'' qu'elle fournit.
Les pièces d'armure maudites ont généralement des enchantements négatifs (des moins)
en plus d'être impossibles à retirer.

%.pg
Beaucoup de types d'armure sont sujets à une dégradation comme la rouille.  Ces
dégâts peuvent être réparés.  Certains types d'armure peuvent gêner l'incantation des sorts.

%.pg
L'option \textit{nudist}
peut être activée (avant le début de la partie) pour tenter de jouer toute la
partie sans porter la moindre armure (un défi volontaire
extrêmement difficile à relever).

%.pg
Certaines interfaces prennent en charge l'option \texttt{armorstatus}.
Lorsqu'elle est activée, un état supplémentaire est affiché, résumant
l'armure actuellement portée.

%.pg
Les commandes liées aux armures sont `\texttt{W}' (porter) et `\texttt{T}' (ôter).
La commande `\texttt{A}' permet d'ôter les armures ainsi que d'autres
objets portés.
De plus, `\texttt{P}' (mettre) et `\texttt{R}' (retirer), normalement destinées aux
accessoires, peuvent servir pour les armures, mais les pièces d'armure ne seront pas
proposées comme candidates probables lorsqu'on vous demande quoi mettre ou retirer.

%.hn 2
\subsection*{Nourriture (`\texttt{\%}')}

%.pg
La nourriture est indispensable à la survie.  Si vous restez trop longtemps sans manger,
vous vous évanouirez, puis finirez par mourir de faim.
Certains types de nourriture se gâtent et deviennent malsains à manger
s'ils ne sont pas protégés.
La nourriture conservée dans des glacières ou des boîtes de conserve
reste généralement fraîche, mais les glacières sont lourdes et les boîtes de conserve
longues à ouvrir.

%.pg
Lorsque vous tuez des monstres, ils laissent généralement des cadavres, qui sont aussi
de la ``nourriture''.  Beaucoup d'entre eux, mais pas tous, sont comestibles ; certains vous confèrent aussi
des pouvoirs spéciaux lorsque vous les mangez.  Une bonne règle empirique : ``on est
ce que l'on mange''.

%.pg
Certains rôles et certains monstres sont végétariens.  Les monstres végétariens
ne mangent généralement jamais de cadavres d'animaux, tandis que les joueurs
végétariens le peuvent, mais avec des effets secondaires assez désagréables.

%.pg
Vous pouvez donner à un aliment le nom de quelque chose que vous aimez manger grâce
à l'option \textit{fruit}.

%.pg
La commande pour manger est `\texttt{e}'.

%.hn 2
\subsection*{Parchemins (`\texttt{?}')}

%.pg
Les parchemins portent diverses étiquettes, sans doute choisies par d'anciens magiciens
pour leur valeur comique (par exemple ``READ ME'', ou ``THANX MAUD'' à l'envers).
Les parchemins disparaissent après lecture (sauf les vierges, qui ne portent
aucune formule magique).

%.pg
L'un des plus utiles est le %
\textit{parchemin d'identification}, qui
permet de déterminer ce qu'est un autre objet, s'il est maudit ou
béni, et combien d'utilisations il lui reste.
Certains objets à l'enchantement subtil sont difficiles à identifier sans eux.

%.pg
Un parchemin dont l'étiquette est connue peut être lu même si le héros est aveugle.
Si un parchemin a été découvert, il apparaît dans l'inventaire sous son type
plutôt que sous son étiquette, mais l'étiquette est connue dans ce cas même si
elle n'est pas affichée.

%.pg
Beaucoup de parchemins produisent un effet différent de l'habituel s'ils sont bénis ou
maudits, ou lus alors que le héros est confus.

%.pg
Un démon postier peut accourir pour vous remettre du courrier sous la forme d'un %
\textit{parchemin de courrier} (sur les versions compilées avec cette fonctionnalité).
Pour utiliser cette fonctionnalité sur les versions où la distribution du courrier de \textit{NetHack}
est déclenchée par l'arrivée de courrier électronique dans votre boîte aux lettres système,
vous devez indiquer à \textit{NetHack} où chercher le nouveau courrier en donnant à la
variable d'environnement \texttt{MAIL} le nom du fichier de votre boîte aux lettres.
Vous pouvez aussi donner à la variable d'environnement \texttt{MAILREADER} le
nom du fichier de votre lecteur préféré, pour que \textit{NetHack} puisse le lancer lorsque vous
lisez le parchemin.
Sur les versions de \textit{NetHack} où le courrier est généré aléatoirement
au sein du jeu, ces variables d'environnement sont ignorées.
Vous pouvez désactiver le démon postier en désactivant
l'option \textit{mail}.

%.pg
La commande pour lire un parchemin est `\texttt{r}'.

%.hn 2
\subsection*{Potions (`\texttt{!}')}

%.pg
Les potions se distinguent par la couleur du liquide contenu dans la fiole.
Elles disparaissent une fois bues.

%.pg
Les potions limpides sont des potions d'eau.  Il arrive qu'elles soient
bénies ou maudites, ce qui donne de l'eau bénite ou de l'eau maudite.  L'eau bénite est
le fléau des morts-vivants ; les potions d'eau bénite sont donc idéales pour leur
être lancées (`\texttt{t}').  Il est aussi parfois très utile de tremper
(``\texttt{\#dip}'') un objet dans une potion.

%.pg
La commande pour boire une potion est `\texttt{q}' (boire).

%.hn 2
\subsection*{Baguettes (`\texttt{/}')}

%.pg
Les baguettes ont généralement plusieurs charges magiques.
Certains types de baguettes doivent être zappées dans une direction.
Vous pouvez aussi
les zapper sur vous-même (indiquez simplement `\texttt{.}' ou `\texttt{s}' comme direction).
Soyez toutefois prévenu : c'est souvent peu judicieux.
D'autres types de baguettes
ne demandent pas de direction.  Le nombre de charges d'une
baguette est aléatoire et diminue d'une unité à chaque utilisation.

%.pg
Lorsque le nombre de charges restantes d'une baguette tombe à zéro, les tentatives
d'utilisation ne donnent généralement rien.  Il est cependant parfois
possible d'extraire les derniers points de mana d'une baguette épuisée,
ce qui la détruit.  Une baguette peut être rechargée par une magie
appropriée, mais cela risque de la faire exploser.  La probabilité
d'une telle explosion est très faible au départ et augmente à chaque
recharge de la baguette.

%.pg
Dans une situation vraiment désespérée, quand vous êtes dos au mur, vous
pouvez décider de jouer le tout pour le tout et de briser votre baguette.  Ce n'est
pas pour les âmes sensibles.  Cela provoquera presque à coup sûr une libération
catastrophique d'énergies magiques.

%.pg
Lorsque vous avez entièrement identifié une baguette, l'affichage de l'inventaire
inclut des informations supplémentaires entre parenthèses : le nombre de fois
où elle a été rechargée, suivi de deux-points puis de son nombre actuel de charges.
Un nombre de charges de \texttt{-1} est un cas particulier indiquant que la baguette
a été annulée.

%.pg
La commande pour utiliser une baguette est `\texttt{z}' (zapper).  Pour en briser une, utilisez la commande `\texttt{a}'
(appliquer).

%.hn 2
\subsection*{Anneaux (`\texttt{=}')}

%.pg
Les anneaux sont des objets très utiles, car leur magie est relativement
permanente, contrairement aux effets généralement fugaces des potions, des
parchemins et des baguettes.

%.pg
Mettre un anneau active sa magie.
Vous pouvez porter au plus deux anneaux à la fois, un à l'annulaire de
chaque main.

%.pg
La plupart des anneaux portés vous font aussi avoir faim plus vite, à un rythme
qui varie selon le type d'anneau.

%.pg
Lorsque vous portez des gants, les anneaux se portent dessous.
Si les gants sont maudits, vous ne pouvez pas mettre d'anneau et ceux déjà
portés ne peuvent pas être retirés.
Lorsque les gants portés ne sont pas maudits, vous n'avez pas à les ôter
manuellement avant de mettre ou de retirer un anneau, puis à les remettre après.
C'est fait implicitement pour éviter un ennui inutile.

%.pg
Les commandes liées aux anneaux sont `\texttt{P}' (mettre) et `\texttt{R}' (retirer).
`\texttt{A}', `\texttt{W}' et `\texttt{T}' peuvent aussi être utilisées ; voir \textit{Amulettes}.

%.hn 2
\subsection*{Livres de sorts (`\texttt{+}')}

%.pg
Les livres de sorts sont des grimoires d'une magie puissante.  Lorsqu'on les étudie avec la commande `\texttt{r}' (lire),
ils transmettent au lecteur la connaissance d'un sort (et
finissent donc par devenir
illisibles)---à moins que la tentative ne tourne mal.
Lire un livre de sorts maudit ou couvert de runes mystiques qui dépassent
votre entendement peut nuire à votre santé !

%.pg
Un sort (même appris) peut aussi mal tourner lorsque vous le lancez.  Si vous
tentez de lancer un sort bien au-dessus de votre niveau d'expérience, si vous avez
peu de compétence dans le type de sort concerné, ou si vous le lancez à
un moment où votre chance est particulièrement mauvaise, vous risquez de gaspiller
à la fois l'énergie et le temps nécessaires à l'incantation.

%.pg
Lancer un sort fait appel à des énergies magiques et les concentre par
la seule force de votre esprit.  Une partie de l'énergie magique libérée vient
de vous.
L'incantation draine temporairement votre pouvoir magique, qui se régénère
lentement, et augmente vos besoins en nourriture.
L'incantation demande aussi de la pratique.  Avec la pratique, votre
compétence dans chaque catégorie de sorts s'améliore.  Avec le temps, en revanche,
votre souvenir de chaque sort s'estompe et vous devrez le réapprendre.

%.pg
Certains sorts doivent être lancés dans une direction, comme les baguettes.
Pour en lancer un sur vous-même, indiquez simplement `\texttt{.}' ou `\texttt{s}' comme direction.
Quelques sorts vous demandent de choisir un emplacement cible plutôt qu'une
simple direction.
D'autres sorts ne demandent ni direction ni cible.

%.pg
De même que les armes sont réparties en groupes dans lesquels un personnage peut
acquérir une maîtrise (à divers degrés), les sorts sont regroupés de façon similaire.
Lancer un sort avec succès exerce son groupe de compétence ; utiliser la
commande ``\texttt{\#enhance}'' pour faire progresser une compétence suffisamment exercée
affecte tous les sorts du groupe.  Une compétence plus avancée peut augmenter la
puissance des sorts, réduire leur risque d'échec lors de l'incantation
et améliorer la précision de l'estimation de la durée pendant laquelle ils
resteront dans votre mémoire.
Les emplacements de compétence sont partagés avec les compétences d'armes.  (Voir aussi la section
``Maîtrise des armes''.)

%.pg
Lancer un sort exige aussi une liberté de mouvement, et le port de divers types
d'armure peut l'entraver.

%.pg
La commande pour lire un livre de sorts est la même que pour les parchemins, `\texttt{r}' (lire).
La commande `\texttt{+}' liste chaque sort que vous connaissez avec son niveau, sa catégorie
de compétence, son taux d'échec à l'incantation et une estimation de la solidité
de votre souvenir.
La commande `\texttt{Z}' (lancer un sort) lance un sort.

%.hn 2
\subsection*{Outils (`\texttt{(}')}

%.pg
Les outils sont des objets divers aux usages variés.  Certains outils
ont un nombre limité d'utilisations, à la manière des charges des baguettes.  Par
exemple, les lampes finissent par s'éteindre.  D'autres outils sont des contenants,
dans lesquels on peut placer des objets ou d'où l'on peut en retirer.

%.pg
Certains outils (comme un bandeau) peuvent être \textit{portés}, et se mettent et
se retirent comme les autres accessoires (anneaux, amulettes) ; voir \textit{Amulettes}.
D'autres outils (comme la pioche) peuvent être maniés comme des armes en plus
d'être appliqués pour leur usage habituel, et dans certains cas (encore la pioche)
deviennent maniés comme une arme même lorsqu'ils sont appliqués.

%.pg
% Mentioned here because of the old method of attempting "Zen" conduct:
% restart until there's a blindfold in starting inventory and put it on
% first thing.
L'option \textit{blind}
peut être activée (avant le début de la partie) pour tenter de jouer toute la
partie sans pouvoir voir (un défi volontaire
très difficile à relever).

%.pg
La commande pour utiliser un outil est `\texttt{a}' (appliquer).

%.hn 3
\subsection*{Contenants}

%.pg
Vous pourrez rencontrer des sacs, des boîtes et des coffres au cours de vos voyages.  Un outil de
ce genre peut être ouvert avec la commande étendue ``\texttt{\#loot}'' lorsque
vous vous tenez dessus (c'est-à-dire sur la même case du sol),
ou avec la commande `\texttt{a}' (appliquer) lorsque vous le portez.  Cependant,
les coffres sont souvent verrouillés, et sont de toute façon encombrants.
Vous devez en poser un avant de le déverrouiller
avec une clé ou un outil de crochetage via la commande `\texttt{a}' (appliquer),
en lui donnant un coup de pied avec la commande `\texttt{\textasciicircum D}',
ou en forçant la serrure avec une arme grâce à la commande étendue
``\texttt{\#force}''.

%.pg
Certains coffres sont piégés, et des choses désagréables se produisent lorsque
vous les déverrouillez ou les ouvrez.  Vous pouvez chercher les pièges et tenter de les désamorcer
avec la commande étendue ``\texttt{\#untrap}''.

%.pg
Lorsque le contenu d'un contenant est connu, ce contenant est
décrit à peu près comme ``sac contenant 3 objets''.
Dans cet exemple, le 3 désigne le nombre de \textit{piles} d'objets
compatibles, et non le nombre total d'objets individuels.
Ainsi, un sac contenant 2 potions bleu ciel, 7 flèches et 350 pièces d'or serait
décrit comme contenant 3 objets plutôt que 10 ou 359.
Et il vous faudrait 3 emplacements d'inventaire libres pour
tout retirer (dans le cas où les objets retirés ne se
combinent pas en piles plus grandes avec ce que vous portez déjà).

%.pg
Si un coffre ou une grande boîte est décrit comme ``à la serrure cassée'', cela
signifie qu'il ne peut plus être verrouillé, et non qu'il ne fonctionne plus comme contenant.

%.pg
Les commandes \textit{appliquer} et \textit{\#loot} vous permettent de retirer et/ou
de placer un nombre quelconque d'objets en une seule opération.
Si vous voulez tout retirer d'un contenant, vous pouvez utiliser la
commande ``\texttt{\#tip}'' pour en verser le contenu sur le sol.
C'est peut-être votre seul moyen de sortir des objets si vos mains sont collées
à une arme à deux mains maudite.
Lorsque vos mains ne sont pas collées, vous pouvez verser le
contenu dans un autre contenant.
(À l'heure où nous écrivons ces lignes, l'autre contenant doit être porté et non
posé sur le sol.)

%.hn 2
\subsection*{Amulettes (`\texttt{"}')}

%.pg
Les amulettes ressemblent beaucoup aux anneaux, et sont souvent plus puissantes.
Comme les anneaux, les amulettes ont diverses propriétés magiques, certaines bénéfiques,
d'autres nuisibles, qui s'activent lorsqu'on les met.

%.pg
On ne peut porter qu'une amulette à la fois, autour du cou.
Comme les anneaux, porter une amulette affecte votre métabolisme et vous
fait avoir faim plus vite.

%.pg
Les commandes liées aux amulettes sont les mêmes que pour les anneaux, `\texttt{P}' (mettre)
et `\texttt{R}' (retirer).
`\texttt{A}' permet de retirer divers objets portés, dont les amulettes.
De plus, '\texttt{W}' (porter) et `\texttt{T}' (ôter), normalement destinées aux
armures, peuvent servir pour les amulettes et les autres accessoires (anneaux et
lunettes), mais les accessoires ne seront pas proposés comme candidats probables
lorsqu'on vous demande quoi porter ou ôter.

%.hn 2
\subsection*{Gemmes (`\texttt{*}')}

%.pg
Certaines gemmes ont de la valeur et peuvent être vendues pour beaucoup d'or.
Elles sont aussi un moyen bien plus efficace de transporter vos richesses.
Les gemmes de valeur augmentent votre score si vous les emportez avec vous en sortant.

%.pg
D'autres petites pierres sont aussi classées parmi les gemmes, mais elles ont bien
moins de valeur.
Toutes les pierres, cependant, peuvent servir de projectiles (si vous avez une fronde).
Dans les cas les plus désespérés, vous pouvez toujours les lancer à la main.

%.hn 2
\subsection*{Grosses pierres (`\texttt{`}')}
%.pg
Les statues et les rochers ne sont pas particulièrement utiles et sont généralement lourds.
On raconte que certaines statues ne sont pas ce qu'elles semblent être.

%.pg
Les rochers bloquent parfois votre chemin.
Vous pouvez en pousser un devant vous (en tentant de marcher sur sa case)
lorsque rien ne bloque \textit{son} chemin, ou vous pouvez
le réduire en un tas de cailloux par une magie destructrice ou avec une pioche.
Il est possible d'aller sur la case d'un rocher si certaines conditions
sont remplies ; d'ordinaire, l'une de ces conditions est qu'il soit impossible
de le pousser plus loin.
Utiliser le préfixe de déplacement sans ramassage (touche par défaut `\texttt{m}')
avant la direction de déplacement permet de tenter d'aller sur la case d'un
rocher sans le pousser, en plus de l'action habituelle du préfixe qui est
d'empêcher le ramassage automatique à destination.

%.pg
Les très grands humanoïdes (les géants et leurs semblables) sont connus pour ramasser
des rochers et s'en servir comme projectiles.

%.pg
Contrairement aux rochers, les statues ne peuvent pas être poussées, mais ce n'est
pas nécessaire puisqu'elles ne bloquent pas le passage.
%.\" 'rumor' above is about statue traps; this is a hint about statue contents
On peut toutefois les réduire en cailloux.

%.pg
Dans certaines configurations du programme, les statues ne sont plus représentées
par `\texttt{`}'
mais par la lettre du monstre qu'elles représentent.

%.hn 2
\subsection*{Or (`\texttt{\$}')}

%.pg
L'or augmente votre score, et vous permet d'acheter des choses dans les boutiques.
Un certain nombre
de monstres du donjon peuvent être influencés par la quantité d'or
que vous portez (sans parler des marchands).

%.pg
Les pièces d'or sont le seul type d'objet auquel l'état de bénédiction ou de
malédiction ne s'applique pas.
Elles sont toujours non maudites, mais jamais décrites comme non maudites, même si
vous désactivez l'option ``\textit{implicit\textunderscore uncursed}''.
Vous pouvez activer l'option ``\textit{goldX}''
si vous préférez que les pièces d'or soient traitées comme ayant un état de
bénédiction ou de malédiction \textit{inconnu} plutôt que connu comme non maudit.
Cela n'a d'importance que lorsque vous utilisez une invite de sélection d'objets
qui peut filtrer selon l'état ``\texttt{BUCX}''.

%.hn 2
\subsection*{Persistance des objets}

%.pg
Normalement, si vous avez vu un objet à un endroit donné de la carte puis
allez à un autre endroit d'où vous ne pouvez plus voir directement cet objet,
il reste affiché sur votre carte.
C'est encore le cas même s'il ne s'y trouve plus
réellement---peut-être un monstre l'a-t-il ramassé, ou a-t-il pourri---jusqu'à
ce que vous puissiez de nouveau voir ou tâter cet endroit.
Une exception notable est le cas où l'objet est recouvert par le marqueur de
``monstre invisible mémorisé''.
Lorsque ce marqueur est ensuite retiré,
après que vous avez vérifié qu'aucun monstre ne s'y trouve, vous aurez oublié
qu'il y avait un objet à cet endroit, que le monstre invisible ait réellement
pris l'objet ou non.
Si l'objet est toujours là, une fois que vous verrez ou tâterez de nouveau cet
endroit, vous redécouvrirez l'objet et recommencerez à vous en souvenir.

%.pg
La situation est la même pour un tas d'objets, à ceci près que seul
l'objet du dessus du tas est affiché.
L'option
\textit{hilite\textunderscore pile}
peut être activée pour afficher différemment un objet lorsqu'il est
au sommet d'un tas.

%.hn 1
\section{Conduite}

%.pg
Comme si gagner à \textit{NetHack} n'était pas assez difficile, certains joueurs
cherchent à se mettre au défi en s'imposant des restrictions dans
leur façon de jouer.  Le jeu suit automatiquement certains de
ces défis, que l'on peut consulter à tout moment avec la commande \texttt{\#conduct}
ou en fin de partie.  Lorsque vous effectuez une action qui
rompt un défi, celui-ci n'est plus listé.  Cela donne aux
joueurs un ``droit de se vanter'' supplémentaire lorsqu'ils gagnent la partie en relevant ces
défis.  Notez qu'il est parfaitement acceptable de gagner la partie
sans recourir à ces restrictions et qu'il est rare que des
joueurs respectent des défis la première fois qu'ils gagnent.

%.pg
Plusieurs défis concernent l'alimentation.  Le plus
difficile est le défi sans nourriture.  Bien que les créatures
puissent survivre longtemps sans nourriture, l'eau est un besoin
physiologique ; il n'y a donc aucune restriction sur les boissons,
même si elles apportent un peu de valeur nutritive.
Invoquer votre dieu pour obtenir de l'aide face à la famine ne
viole non plus aucun défi alimentaire.

%.pg
Un régime végétalien strict exclut tout aliment d'origine animale.
La principale source de nutrition est constituée de fruits et de légumes.  Les
cadavres et les boîtes de conserve de blobs (`\texttt{b}'), de gelées (`\texttt{j}') et de champignons
(`\texttt{F}') sont aussi considérés comme végétaux.  Certains aliments
humains sont préparés sans produits animaux : les galettes de lembas, les rations
de cram, les rations de nourriture (gunyoki), les rations K et les rations C.
Le métal ou toute autre matière normalement indigeste mangé alors que vous êtes
métamorphosé en une créature capable de le digérer est aussi considéré comme végétalien.
Notez cependant que manger ces objets compte quand même contre la conduite sans nourriture.

%.pg
Les végétariens ne mangent pas d'animaux ;
ils sont toutefois moins regardants que les végétaliens sur les produits d'origine animale.
En plus des aliments végétaliens cités plus haut, ils peuvent manger n'importe quel
pouding (`\texttt{P}') autre que les poudings noirs,
les œufs et les aliments à base d'œufs (biscuits chinois et crêpes),
les aliments à base de lait (tartes à la crème et barres chocolatées) et les morceaux
de gelée royale.  On attend des Moines qu'ils suivent un régime végétarien.

%.pg
Manger de la viande, quelle qu'elle soit, viole les conduites végétarienne,
végétalienne et sans nourriture.  Cela inclut les rations de tripes, les cadavres ou
boîtes de conserve de tous les monstres non cités ci-dessus, et les divers autres
morceaux de viande trouvés dans le donjon.  Avaler et digérer un monstre lorsque
vous êtes métamorphosé est traité comme si vous aviez mangé le cadavre de la créature.
Manger des objets en cuir, en peau de dragon ou en os alors que vous êtes
métamorphosé en une créature capable de les digérer, ou manger des cerveaux de monstres
alors que vous êtes métamorphosé en flagelleur mental, est considéré comme manger
un animal, alors que la cire n'est qu'un produit d'origine animale.

%.pg
Quelle que soit la conduite, certains objets sont indigestes
et d'autres dangereux à manger.  Utiliser une attaque d'avalement et de
digestion contre un monstre équivaut à manger son cadavre.
Notez que le terme ``végétalien'' n'est employé ici que dans le contexte de
l'alimentation.  Vous restez libre de choisir de ne pas utiliser ni porter d'objets d'origine
animale (cuir, peau de dragon, os, cornes, corail, etc.), mais le
jeu n'en tiendra pas le compte pour vous.  Notez aussi que les potions ``laiteuses''
peuvent être d'un blanc translucide, mais qu'elles ne contiennent pas de lait ;
elles sont donc compatibles avec un régime végétalien.  Les moisissures gluantes ou
les ``fruits'' définis par le joueur, bien qu'ils puissent être n'importe quoi,
des ``cerises'' aux ``côtelettes de porc'', sont aussi supposés végétaliens.

%.pg
Un athée est quelqu'un qui rejette la religion.  Cela signifie que vous ne pouvez pas
utiliser \texttt{\#pray}, faire des sacrifices (\texttt{\#offer}) à un dieu,
repousser les morts-vivants (\texttt{\#turn}) ni utiliser \texttt{\#chat} avec un prêtre.
Les lecteurs particulièrement pointilleux pourraient objecter que jouer un Moine ou
un Prêtre devrait violer cette conduite ; ce choix est laissé au
joueur.  Offrir l'Amulette de Yendor à votre dieu est nécessaire pour
gagner la partie et n'est pas compté contre cette conduite.  Vous n'êtes pas non
plus pénalisé lorsqu'un dieu en colère, un prêtre (ou une prêtresse) ou une
autre figure religieuse vous adresse la parole ; un véritable athée entendrait les
mots sans leur attacher de signification particulière.

%.pg
Un démuni commence la partie sans possessions, sans sorts et sans compétences
d'armes ni de sorts (et si vous jouez un chevalier, votre poney n'aura pas de selle).
Cette conduite ne peut être engagée qu'en commençant une nouvelle partie avec \texttt{OPTIONS=pauper}
dans votre fichier de configuration ou dans la variable d'environnement \texttt{NETHACKOPTIONS}.
Une fois la partie lancée, vous pouvez acquérir et utiliser objets, sorts et
compétences de la manière habituelle.

%.pg
La plupart des joueurs combattent avec une arme maniée (ou un outil destiné à être
manié comme une arme).  Un autre défi consiste à gagner la partie sans
utiliser une telle arme maniée.  Vous avez toujours le droit de lancer,
de tirer et de frapper du pied des armes ; d'utiliser une baguette, un sort ou un autre
type d'objet ; ou de combattre avec vos mains et vos pieds.

%.pg
Dans \textit{NetHack}, un pacifiste refuse de causer la mort de tout autre monstre
(c'est-à-dire chaque fois que cette mort vous rapporterait de l'expérience).  C'est
un défi particulièrement difficile, même s'il reste possible de gagner de
l'expérience par d'autres moyens.

%.pg
Un personnage illettré ne lit pas et n'écrit pas.  Cela inclut la lecture
d'un parchemin, d'un livre de sorts, du message d'un biscuit chinois ou d'un T-shirt ;
l'écriture d'un parchemin ; ou la gravure de quoi que ce soit d'autre qu'un simple ``X'' (la
signature traditionnelle d'une personne illettrée).  Lire une gravure,
ou tout objet absolument nécessaire pour gagner la partie, n'est pas compté
contre cette conduite.  L'identité des parchemins et des livres de sorts (et la
connaissance des sorts) de votre inventaire de départ est supposée avoir été
apprise auprès de vos maîtres avant le début de la partie et n'est pas
comptée.

%.pg
Il existe une branche latérale du donjon principal appelée ``Sokoban'', brièvement
décrite plus haut dans la section sur les \textit{Pièges}.
Comme indiqué alors, le but est de pousser des rochers dans des fosses et/ou des trous
pour les combler, afin à la fois de dégager les rochers et de pouvoir
franchir les pièges.
Certaines ``règles'' spéciales s'appliquent dans cette branche
du donjon.
Certaines règles ne peuvent pas être contournées, comme l'impossibilité de pousser
un rocher en diagonale.
D'autres le peuvent, comme l'interdiction de briser les rochers par la magie ou
avec des outils, mais les enfreindre vous inflige une pénalité de chance.
Aucun message ne vous le signale sur le moment, mais c'est suivi comme une conduite.
La commande \texttt{\#conduct} et le récapitulatif de fin de partie indiquent si
vous avez respecté les règles spéciales de Sokoban et, sinon, combien de
fois vous les avez enfreintes, ce qui vous permet de découvrir quelles
actions portent malheur et de mieux décider s'il faut
éviter ou non de répéter ces actions à l'avenir.
(Remarque : la conduite de Sokoban n'est affichée que si vous êtes
entré dans la branche Sokoban du donjon au cours de la partie.
Dès lors, elle fait partie des conduites révélées même si
vous n'y avez rien fait d'intéressant.
Terminer la partie avec la conduite ``n'a enfreint aucune des règles spéciales
de Sokoban'' a surtout du sens si vous parvenez aussi à réaliser
l'exploit ``a obtenu le trésor de Sokoban''
(voir \textit{Exploits} plus bas).)

%.pg
Le jeu suit plusieurs autres défis.  Il est possible
d'éliminer une ou plusieurs espèces de monstres par génocide ; jouer sans
cette possibilité est considéré comme un défi.  Lorsque le jeu vous offre
l'occasion de génocider des monstres, vous pouvez répondre par le type de monstre
``aucun'' si vous voulez refuser.  Vous pouvez changer la forme d'un objet en
un autre objet du même type (``polypiling'') ou la forme de votre propre
corps en une autre créature (``polyself'') au moyen d'une baguette, d'un sort ou
d'une potion de métamorphose ; éviter chacun de ces effets constitue un défi.
Métamorphoser des monstres, animaux de compagnie compris, ne rompt aucun de ces
deux défis.
Enfin, il peut vous arriver de recevoir des vœux ; une partie sans tentative de
souhaiter le moindre objet est un défi, tout comme une partie sans souhaiter
d'artefact (même si l'artefact disparaît aussitôt).  Lorsque le
jeu vous offre l'occasion de faire un vœu, vous pouvez
choisir ``rien'' si vous voulez refuser.

%.hn 2
\subsection*{Exploits}

%.pg
Le récapitulatif de fin de partie affiche aussi divers exploits
représentant votre progression vers l'ascension finale, si vous en avez
accompli.
Ils ne sont pas directement liés à la \textit{conduite}, mais ils sont regroupés avec
elle parce qu'ils relèvent de la même catégorie de ``droit de se vanter''
et pour limiter le nombre de questions lors du récapitulatif.
Ils sont listés ici à peu près par ordre de difficulté, et pas forcément dans l'ordre
dans lequel vous pourriez les accomplir.

% [length stuff copied from paranoid_confirmation]
\newlength{\achwidth}
%.PS "Mines'\~End\~"
\settowidth{\achwidth}{\texttt{Fin~des~Mines~}}
\addtolength{\achwidth}{\labelsep}
\begin{description}[leftmargin=\achwidth, topsep=1mm, itemsep=0mm, labelwidth=*, font=\ttfamily, align=right]
%.PL Shop
\item[<Rang>]
A atteint le titre de rang \textit{Rang}.
\item[Boutique]
Est entré dans une boutique.
\item[Temple]
Est entré dans un temple.
\item[Mines]
Est entré dans les Mines des Gnomes.
\item[Minebourg]
Est entré dans Minebourg.
\item[Oracle]
A consulté l'Oracle de Delphes.
\item[Roman]
A lu un passage d'un roman du Disque-monde.
\item[Sokoban]
Est entré dans Sokoban.
\item["Grande~Salle"]
Est entré dans la Grande Salle.
\item["Prix-Soko"]
A exploré Sokoban jusqu'en haut et y a trouvé un objet spécial.
\item[Fin~des~Mines]
A exploré les Mines des Gnomes jusqu'en bas et y a trouvé un objet spécial.
\item[Méduse]
A vaincu Méduse.
\item[Air]
A découvert l'air qui permet d'ouvrir et de fermer le pont-levis du
niveau du Château.
\item[Cloche]
A obtenu la Cloche d'Ouverture.
\item[Géhenne]
Est entré dans la Géhenne.
\item[Candélabre]
A obtenu le Candélabre d'Invocation.
\item[Livre]
A obtenu le Livre des Morts.
\item[Invocation]
A obtenu l'accès au niveau le plus profond de la Géhenne.
\item[Amulette]
A obtenu la légendaire Amulette de Yendor.
\item[Fin~de~partie]
A atteint les Plans élémentaires.
\item[Astral]
A atteint le niveau du Plan astral.
\item[Aveugle]
Aveugle de naissance.
\item[Sourd]
Sourd de naissance.
\item[Nudiste]
N'a jamais porté d'armure.
\item[Démuni]
A commencé sans possessions.
\item[Ascension]
A porté l'Amulette à sa destination finale.
\end{description}
%.PE
%.ED

%.lp "Notes:  "
\noindent
Remarques :

%.pg
Les exploits sont enregistrés puis rapportés dans l'ordre où ils
se produisent au cours de votre partie plutôt que dans l'ordre listé ici.

%.pg
Il existe neuf titres de \textit{Rang} pour chaque rôle, attribués aux niveaux
d'expérience 1, 3, 6, 10, 14, 18, 22, 26 et 30.
Celui du niveau d'expérience 1 n'est pas enregistré comme exploit.
Perdre assez de niveaux pour redescendre à un ou plusieurs rangs inférieurs ne
supprime pas les exploits correspondants.

%.pg
Il n'y a pas de \textit{Roman} garanti ; l'exploit consistant à en lire un peut donc
ne pas toujours être réalisable (sauf peut-être par un vœu).
De même, le niveau de la \textit{Grande Salle} n'est pas toujours présent.
Contrairement au roman, il n'existe aucun moyen de souhaiter cette occasion.

%.pg
Les ``objets spéciaux'' cachés dans la \textit{Fin~des~Mines} et dans \textit{Sokoban}
ne sont pas uniques, mais sont considérés comme des prix ou des récompenses
pour avoir exploré ces niveaux, puisque cela n'est pas nécessaire pour terminer
la partie.
Trouver d'autres exemplaires des mêmes objets n'enregistre pas
l'exploit correspondant.

%.pg
L'exploit \textit{Méduse} est enregistré si elle meurt pour quelque raison que ce soit,
même si vous n'en êtes pas directement responsable, et seulement si elle meurt.

%.pg
L'\textit{air} de 5 notes peut être appris par essais et erreurs avec un instrument
de musique joué suffisamment
près---mais pas trop !---du
pont-levis du niveau du Château, ou vous être donné en réponse à une prière.

%.pg
\textit{Aveugle}, \textit{Sourd}, \textit{Nudiste} et \textit{Démuni} sont aussi des conduites, et elles ne peuvent être
activées qu'en réglant l'option correspondante dans \texttt{NETHACKOPTIONS}
ou dans le fichier de configuration avant le début de la partie.
Dans le cas d'\textit{Aveugle} et de \textit{Sourd}, l'option impose aussi la
conduite.
Ce ne sont pas vraiment des accomplissements notables tant que vous n'avez pas
progressé de façon substantielle dans le donjon.

%.hn 1
\section{Options}

%.pg
Les goûts de chacun et les idées sur la façon dont \textit{NetHack} devrait
faire les choses étant variés, il existe des options que vous pouvez
régler pour modifier le comportement de \textit{NetHack}.

%.hn 2
\subsection*{Régler les options}

%.pg
Les options peuvent être réglées de plusieurs façons. En cours de partie, la commande `\texttt{O}'
vous permet de consulter toutes les options et de modifier la plupart d'entre elles.
Vous pouvez aussi régler les options automatiquement en les plaçant dans un fichier
de configuration, ou dans la variable d'environnement ``\texttt{NETHACKOPTIONS}''.
Certaines versions de \textit{NetHack} disposent aussi de programmes frontaux qui
permettent de régler les options avant de lancer la partie, ou d'une configuration
globale destinée aux administrateurs système.

%.hn 2
\subsection*{Utiliser un fichier de configuration}

%.pg
Le nom et l'emplacement par défaut du fichier de configuration varient selon le
système d'exploitation.\\

%.lp ""
Sous UNIX, Linux et macOS, il s'agit de \mbox{``.nethackrc''} dans le répertoire
personnel de l'utilisateur.
Ce fichier peut ne pas exister, mais c'est un fichier texte ASCII ordinaire qui
peut être créé avec n'importe quel éditeur de texte.\\

%.lp ""
Sous Windows, il s'agit de \mbox{``.nethackrc''}, situé dans le dossier
\mbox{{``\%USERPROFILE\%\textbackslash NetHack\textbackslash''}}.
Ce fichier peut ne pas exister,
mais c'est un fichier texte ASCII ordinaire qui peut être créé avec n'importe quel
éditeur de texte.
Après avoir lancé \textit{NetHack} pour la première fois, vous devriez trouver un modèle
par défaut du fichier de configuration nommé \mbox{``.nethackrc.template''} dans
\mbox{{``\%USERPROFILE\%\textbackslash NetHack\textbackslash''}}.
Si vous n'avez pas créé de fichier de configuration, \textit{NetHack} en créera
un pour vous à partir de ce modèle par défaut.\\

%.lp ""
Sous MS-DOS, il s'agit de \mbox{``defaults.nh''} dans le même dossier que
\mbox{\textit{nethack.exe}}.\\

%.lp ""
Toute ligne du fichier de configuration commençant par `\texttt{\#}' est traitée
comme un commentaire et ignorée.
Les lignes vides sont ignorées.

Toute ligne commençant par `\texttt{[}' et se terminant par `\texttt{]}'
est un marqueur de section (le `\texttt{]}' fermant peut être suivi
d'espaces puis d'un commentaire quelconque commençant par `\texttt{\#}').
Le texte entre les crochets est le nom de la section.
Les marqueurs de section ne sont valables qu'après une directive CHOOSE, et leurs noms
ne tiennent pas compte de la casse.
Les lignes qui suivent un marqueur de section appartiennent à cette section jusqu'au
début d'une autre section, jusqu'à un marqueur sans nom ou jusqu'à la fin du fichier.
Les lignes d'une section sont ignorées, sauf si une directive CHOOSE a sélectionné
cette section.

%.pg
Vous pouvez utiliser différentes directives de configuration dans le fichier, dont
certaines peuvent être utilisées plusieurs fois.
En général, les directives sont
écrites en majuscules, suivies d'un signe égal, puis de
réglages propres à cette directive.

%.pg
Voici la liste des directives autorisées :

%.lp
\begin{description}
\item[OPTIONS]
Il existe deux types d'options : les options booléennes et les options composées.
Les options booléennes activent ou désactivent un réglage, tandis que les options
composées acceptent des valeurs plus variées.
Faites précéder une option booléenne de `no' ou de `!' pour la désactiver.
Pour les options composées, le nom de l'option et sa valeur sont séparés par deux-points.
Certaines options sont persistantes et ne s'appliquent qu'aux nouvelles parties.
Vous pouvez indiquer plusieurs directives OPTIONS, et plusieurs options
séparées par des virgules dans une même directive OPTIONS.
(Les options séparées par des virgules sont traitées de droite à gauche.)

%.lp ""
Exemple :
%.sd
\begin{verbatim}
    OPTIONS=dogname:Fido
    OPTIONS=!legacy,autopickup,pickup_types:$"=/!?+
\end{verbatim}
%.ed

%.lp
\item[HACKDIR]
{[}Non disponible sous Unix]\\
Emplacement par défaut des fichiers dont \textit{NetHack} a besoin. Sous Windows, HACKDIR
correspond par défaut à l'emplacement du fichier \textit{NetHack.exe} ou \textit{NetHackw.exe} ;
il n'est donc généralement ni nécessaire ni recommandé de régler HACKDIR pour le remplacer.
%.lp
\item[LEVELDIR]
{[}Non disponible sous Unix]\\
L'emplacement où sont stockés les fichiers des niveaux en cours. Par défaut HACKDIR ;
doit être accessible en écriture.
%.lp
\item[SAVEDIR]
{[}Non disponible sous Unix]\\
L'emplacement où sont conservées les parties sauvegardées. Par défaut HACKDIR ; doit être
accessible en écriture.
%.lp
\item[BONESDIR]
{[}Non disponible sous Unix]\\
L'emplacement où sont conservés les fichiers d'ossements (bones). Par défaut HACKDIR ; doit être
accessible en écriture.
%.lp
\item[LOCKDIR]
{[}Non disponible sous Unix]\\
L'emplacement où sont stockés les verrous de synchronisation de fichiers. Par défaut
HACKDIR ; doit être accessible en écriture.
%.lp
\item[TROUBLEDIR]
{[}Non disponible sous Unix]\\
L'emplacement où est conservé un relevé des abandons de partie et des problèmes
autodiagnostiqués. Par défaut HACKDIR ; doit être accessible en écriture.
%
% config file entries beyond this point are shown alphabetically
%
%.lp
\item[AUTOCOMPLETE]
Active ou désactive l'autocomplétion d'une commande étendue.
L'autocomplétion n'a aucun effet avec l'interface X11.
Vous pouvez indiquer plusieurs autocomplétions. Pour activer
l'autocomplétion, indiquez la commande étendue. Faites précéder la
commande de ``\texttt{!}'' pour désactiver l'autocomplétion
de cette commande.

%.lp ""
Exemple :
%.sd
\begin{verbatim}
   AUTOCOMPLETE=zap,!annotate
\end{verbatim}
%.ed

%.lp
\item[AUTOPICKUP\textunderscore EXCEPTION]
Définit des exceptions à l'option {\textit{pickup\textunderscore types}}.
Voir la section ``Configurer les exceptions de ramassage automatique''.
%.lp
\item[BINDINGS]
Modifie l'association des touches de certaines touches spéciales, des raccourcis de menu,
des commandes étendues ou des boutons de la souris. Vous pouvez indiquer plusieurs associations.
Le format est la touche suivie de la commande, séparées par deux-points.
Voir la section ``Modifier les associations de touches'' pour plus d'informations.

%.lp ""
Exemple :
%.sd
\begin{verbatim}
   BIND=^X:getpos.autodescribe
\end{verbatim}
%.ed

%.lp
\item[CHOOSE]
Choisit au hasard l'un des paramètres séparés par des virgules comme nom de
section active.
Les lignes des autres sections sont ignorées.

%.lp ""
Exemple :
%.sd
\begin{verbatim}
   OPTIONS=color
   CHOOSE=char A,char B
   [char A]
   OPTIONS=role:arc,race:dwa,align:law,gender:fem
   [char B]
   OPTIONS=role:wiz,race:elf,align:cha,gender:mal
   [] #fin du CHOOSE
   OPTIONS=!rest_on_space
\end{verbatim}
%.ed

%.lp ""
Si \texttt{[]} est présent, la section précédente est fermée et aucune nouvelle
section ne commence ; ce qui suit sera commun à toutes les sections.
Sinon, la dernière section s'étend jusqu'à la fin du fichier d'options.

%.lp
\item[MENUCOLOR]
Met en valeur des lignes de menu avec différentes couleurs.
Voir la section ``Configurer les couleurs des menus''.
%.lp
\item[MSGTYPE]
Modifie la façon dont les messages sont affichés sur la ligne du haut.
Voir la section ``Configurer les types de messages''.
%.lp
\item[ROGUESYMBOLS]
Symboles personnalisés pour le jeu de symboles du niveau rogue.
Voir \textit{SYMBOLS} ci-dessous.
%.lp
\item[SOUND]
Définit une association de son.
Voir la section ``Configurer les sons de l'utilisateur''.
%.lp
\item[SOUNDDIR]
Définit le répertoire qui contient les fichiers sonores.
Voir la section ``Configurer les sons de l'utilisateur''.
%.lp
\item[SYMBOLS]
Remplace un ou plusieurs symboles du jeu de symboles utilisé pour tous les niveaux
du donjon, sauf le niveau spécial rogue.
Voir la section ``Modifier les symboles de \textit{NetHack}''.
%.pg

%.lp ""
Exemple :
%.sd
\begin{verbatim}
   # remplacer la petite ponctuation (apostrophes) par des chiffres
   SYMBOLS=S_golem:7
\end{verbatim}
%.ed

%.lp
\item[WIZKIT]
Mode débogage uniquement : objets supplémentaires à ajouter à l'inventaire initial.
La valeur est le nom d'un fichier texte contenant une liste de noms d'objets,
un par ligne, jusqu'à un maximum de 128 lignes.
Chaque ligne est traitée par la fonction qui gère les vœux.
Les entrées sont ajoutées à l'historique des vœux ; voir la commande wizwish.

%.lp ""
Exemple :
%.sd
\begin{verbatim}
   WIZKIT=~/wizkit.txt
\end{verbatim}
%.ed
\end{description}

%.lp ""
%.pg
Voici un exemple de contenu de fichier de configuration :
%.sd
\begin{verbatim}
    # Régler le rôle, la race, le sexe et l'alignement de votre personnage.
    OPTIONS=role:Valkyrie, race:Human, gender:female, align:lawful

    # Activer le ramassage automatique, régler les types d'objets ramassés
    OPTIONS=autopickup,pickup_types:$"=/!?+

    # Personnalisation de la carte
    OPTIONS=color           # Afficher en couleur si possible
    OPTIONS=lit_corridor    # Distinguer les couloirs éclairés
    OPTIONS=hilite_pet,hilite_pile
    # Remplacer la petite ponctuation (apostrophes) par des chiffres
    OPTIONS=boulder:0
    SYMBOLS=S_golem:7

    # Pas d'écran d'accueil au démarrage. Interface graphique Windows seulement.
    OPTIONS=!splash_screen
\end{verbatim}
%.ed
%.BR 2

%.hn 2
\subsection*{Utiliser la variable d'environnement \texttt{NETHACKOPTIONS}}

%.pg
La variable NETHACKOPTIONS est une liste, séparée par des virgules, de valeurs
initiales pour les différentes options. Certaines ne peuvent qu'être activées ou désactivées.
Vous activez l'une d'elles en ajoutant le nom de l'option à la liste,
et la désactivez en tapant un `\texttt{!}' ou ``\texttt{no}'' devant le nom.
D'autres prennent une
chaîne de caractères pour valeur. Vous pouvez régler les options de type chaîne en tapant
le nom de l'option, deux-points ou un signe égal, puis la valeur de la chaîne.
La valeur se termine à la virgule suivante ou à la fin de la chaîne.

%.pg
Par exemple, pour définir une variable d'environnement de sorte que
\textit{color} soit \texttt{on},
\textit{legacy} soit \texttt{off},
le \textit{name} du personnage soit ``\texttt{Blue Meanie}'',
et le \textit{fruit} nommé soit ``\texttt{lime}'',
vous saisiriez la commande
%.SD i
\begin{verbatim}
    setenv NETHACKOPTIONS "color,\!leg,name:Blue Meanie,fruit:lime"
\end{verbatim}
%.ED

\noindent sous \textit{csh}
(notez la nécessité d'échapper le `!', car il est spécial
pour ce shell), ou la paire de commandes
%.SD i
\begin{verbatim}
    NETHACKOPTIONS="color,!leg,name:Blue Meanie,fruit:lime"
    export NETHACKOPTIONS
\end{verbatim}
%.ED

\noindent sous \textit{sh}, \textit{ksh} ou \textit{bash}.

%.pg
La valeur de NETHACKOPTIONS équivaut en pratique à une seule directive OPTIONS
dans un fichier de configuration.
Le préfixe ``OPTIONS='' est implicite et les options séparées par des virgules sont
traitées de droite à gauche.
Les autres types de directives de configuration, comme BIND ou MSGTYPE, ne
sont pas autorisés.

%.pg
Au lieu d'une liste d'options séparées par des virgules,
NETHACKOPTIONS peut contenir le nom complet d'un fichier de configuration que
vous voulez utiliser.
Si ce nom complet ne commence pas par une barre oblique, faites-le précéder de `\texttt{@}'
(arobase) pour indiquer à NetHack que le reste est un nom de fichier.
S'il commence par `\texttt{/}', l'arobase est facultative.

%.hn 2
\subsection*{Options de personnalisation}

%.pg
Voici l'explication de ce que font les différentes options.
Les chaînes de caractères trop longues peuvent être tronquées.
Certaines des options listées peuvent être inactives dans votre donjon.

%.pg
Certaines options sont persistantes : elles sont sauvegardées et rechargées avec
la partie. Modifier une option persistante dans le fichier de configuration
ne s'applique qu'aux nouvelles parties.

\begin{description}[font=\mdseries\itshape]
%.lp
\item[accessiblemsg]
Ajoute des informations de position ou de direction aux messages (désactivée par défaut).
%.lp
\item[acoustics]
Active les messages sur ce qu'entend votre personnage (activée par défaut).
Notez que cela n'a rien à voir avec les capacités audio de votre ordinateur.
Persistante.
%.lp
\item[alignment]
Votre alignement de départ (\texttt{align:lawful}, \texttt{align:neutral}
ou \texttt{align:chaotic}).
Vous pouvez n'indiquer que la première lettre.
De nombreux rôles et les races non humaines restreignent les alignements autorisés.
Voir \textit{role}
pour savoir comment utiliser la négation afin d'exclure des choix.
%.lp ""
\\
Si \texttt{align} n'est pas indiqué, il n'y a pas de valeur par défaut ;
le joueur sera interrogé, sauf si le rôle ou la race impose un alignement.
Ne peut pas être réglée avec la commande `\texttt{O}'. Persistante.
%.lp
\item[armorstatus]
Affiche un état supplémentaire qui résume l'armure actuellement portée
(désactivée par défaut, non prise en charge par toutes les interfaces).
%.lp ""
\\
Dans le cas habituel où plusieurs pièces d'armure sont portées, une liste de
lettres est affichée dans l'ordre suivant :
\\
\texttt{G} - gants ;\\
\texttt{C} - cape ;\\
\texttt{A} - armure de corps ;\\
\texttt{U} - chemise ;\\
\texttt{H} - casque ;\\
\texttt{B} - bottes ;\\
\texttt{S} - bouclier.
\item[autodescribe]
Décrit automatiquement le terrain sous le curseur lorsqu'il faut indiquer une position
sur la carte (activée par défaut).
L'option \textit{whatis\textunderscore coord}
détermine si la description inclut les coordonnées sur la carte.
%.lp
\item[autodig]
Creuse automatiquement si vous maniez un outil de creusage et que vous avancez vers un endroit
qui peut être creusé (désactivée par défaut). Persistante.
%.lp
\item[autoopen]
Marcher contre une porte fermée tente de l'ouvrir (activée par défaut).
Persistante.
%.lp
\item[autopickup]
Ramasse automatiquement les objets sur lesquels vous vous déplacez (désactivée par défaut).
Persistante.
\\
%.lp ""
Voir ``\textit{pickup\textunderscore types}'' ainsi que
``\textit{autopickup\textunderscore exception}'' pour affiner ce comportement.
\\
%.lp ""
Remarque : avant la version 5.0.0, la valeur par défaut de \textit{autopickup} était \textit{on}.
%.lp
\item[autoquiver]
Cette option détermine ce qui se passe lorsque vous tentez la commande `\texttt{f}' (tirer)
alors que rien n'est placé dans le carquois ni préparé (désactivée par défaut).
Si elle est activée, l'ordinateur remplira
votre carquois ou votre sac de carquois, ou préparera une arme adaptée.
Notez qu'il ne tiendra pas
compte de l'état béni/maudit, de l'enchantement, des dégâts ni de la
qualité de l'arme ; vous restez libre de remplir vous-même votre carquois
ou votre sac de carquois, ou de préparer une arme,
avec la commande `\texttt{Q}'.
Si aucune arme n'est trouvée ou si l'option est
désactivée, c'est la commande `\texttt{t}' (lancer) qui est exécutée à la place. Persistante.
%.lp
\item[autounlock]
%\hyphenation{apply\-key}%this needs to be tested...
Détermine l'action à entreprendre lorsque vous tentez d'entrer dans une porte verrouillée
ou de piller un conteneur verrouillé.
Accepte une liste de valeurs séparées par des signes plus :
% paranoid
% au => autounlock
\newlength{\auwidth}
%.PS Apply-Key
\settowidth{\auwidth}{\texttt{Apply-Key}}
\addtolength{\auwidth}{\labelsep}
\begin{description}[leftmargin=\auwidth, topsep=1mm, itemsep=0mm, labelwidth=*, font=\ttfamily, align=right]
%.PL Untrap
\item[Untrap]
demander s'il faut chercher un piège ;
la recherche peut échouer même s'il y en a un ; si elle en trouve un, on
vous demandera si vous voulez tenter de le désamorcer ; si vous refusez,
votre personnage oubliera que la porte ou la boîte est piégée ;
%.PL Apply-Key
\item[Apply-Key]
si vous portez une clé ou un autre outil de déverrouillage, proposer de l'utiliser ;
%.PL Kick
\item[Kick]
donner un coup de pied dans la porte (si vous omettez untrap ou refusez de tenter
untrap, et que vous omettez apply-key, n'avez pas de clé ou refusez d'utiliser la clé ;
sans effet sur les conteneurs) ;
%.PL Force
\item[Force]
tenter de forcer le couvercle d'un conteneur avec l'arme que vous maniez
actuellement (si vous omettez untrap ou refusez de tenter untrap, et que
vous omettez apply-key, n'avez pas de clé ou refusez d'utiliser la clé ;
sans effet sur les portes) ;
%.PL None
\item[None]
aucune des actions ci-dessus ; ne peut pas être combinée avec les autres choix.
%.PE
\end{description}
Omettre la valeur équivaut à \texttt{autounlock:apply-key}.
Faire précéder \texttt{autounlock} de `\texttt{!}' ou de ``\texttt{no}'' équivaut à
\texttt{autounlock:none}.
\\
%.lp ""
Appliquer une clé peut déclencher un piège si la porte ou le conteneur est piégé.
Un coup de pied réussi dans une porte la brise et réveille les monstres proches.
Forcer avec succès un conteneur brise sa serrure et peut aussi
détruire une partie de son contenu, endommager votre arme, ou les deux.
\\
%.lp ""
La valeur par défaut est Apply-Key.
Persistante.
%.lp
\item[blind]
Rend le personnage aveugle en permanence dès le départ (désactivée par défaut). Persistante.
%.lp
\item[bones]
Autorise la sauvegarde et le chargement des fichiers d'ossements (activée par défaut). Persistante.
%.lp
\item[boulder]
Définit le caractère utilisé pour afficher les rochers (par défaut, le symbole de la
classe des ``grosses pierres'', `\texttt{`}').
%.lp
\item[catname]
Donne un nom à votre chat de départ (par exemple ``\texttt{catname:Morris}'').
Ne peut pas être réglée avec la commande `\texttt{O}'.
%.lp character
\item[character]
Synonyme de ``\texttt{role}'' pour choisir le type de votre personnage
(par exemple ``\texttt{character:Monk}''). Voir \textit{role} pour plus de détails.
%.lp
\item[checkpoint]
Sauvegarde l'état de la partie après chaque changement de niveau, pour une éventuelle
récupération après un plantage du programme (activée par défaut). Persistante.
%.lp
\item[cmdassist]
Le jeu fournit une aide supplémentaire sur les commandes aux nouveaux
joueurs s'il détecte certaines erreurs prévisibles (activée par défaut).
%.lp
\item[confirm]
Demande à l'utilisateur de confirmer les attaques contre les familiers, les commerçants et les autres
créatures pacifiques (activée par défaut). Persistante.
%.lp
\item[dark\textunderscore room]
Affiche les zones hors de vue des pièces éclairées (activée par défaut). Persistante.
%.lp
\item[deaf]
Rend le personnage sourd en permanence dès le départ (désactivée par défaut). Persistante.
%.lp
\item[dropped\textunderscore nopick]
Si cette option est activée, les objets que vous avez posés ne seront pas ramassés automatiquement,
même si ``\textit{autopickup}'' est aussi activée et qu'ils figurent dans
``\textit{pickup\textunderscore types}'' ou correspondent à une exception de ramassage automatique positive
(activée par défaut). Persistante.
%.lp
\item[disclose]
Détermine quelles informations le programme révèle à la fin de la partie.
La valeur est une liste, séparée par des espaces, de paires demande/catégorie
(par défaut `\texttt{ni na nv ng nc no}',
demander avec `\texttt{n}' comme réponse par défaut pour chaque catégorie).
Persistante.
Les possibilités sont :

%.sd
%.si
\texttt{i} --- révéler votre inventaire ;\\
\texttt{a} --- révéler vos attributs ;\\
\texttt{v} --- résumer les monstres vaincus ;\\
\texttt{g} --- lister les espèces de monstres génocidées ;\\
\texttt{c} --- afficher votre conduite, ainsi que vos accomplissements éventuels ;\\
\texttt{o} --- afficher le survol du donjon.
%.ei
%.ed

Chaque possibilité de révélation peut éventuellement être précédée d'un préfixe qui
vous permet d'affiner son comportement. Voici les préfixes valides :

%.sd
%.si
\texttt{y} --- vous demander, avec oui comme réponse par défaut ;\\
\texttt{n} --- vous demander, avec non comme réponse par défaut ;\\
\texttt{+} --- révéler sans demander ;\\
\texttt{-} --- ne pas révéler et ne pas demander.
%.ei
%.ed

La liste des monstres vaincus peut être triée ;
il existe donc deux choix supplémentaires pour `\texttt{v}' :
Les listes des monstres vaincus et des espèces génocidées peuvent être triées ;
il existe donc deux choix supplémentaires pour `v' et `g' :
%.sd
%.si
\texttt{?} --- vous demander, avec la demande d'ordre comme réponse par défaut ;\\
{\tt\#} --- révéler sans demander, mais demander l'ordre de tri.
%.ei
%.ed

Demander l'ordre consiste à choisir l'un des ordres de tri dans un menu.
Le choix `\texttt{+}' (révéler sans demander),
ou le fait d'être interrogé et de répondre `\texttt{y}' plutôt que `\texttt{a}',
affiche par défaut les monstres dans l'ordre indiqué par l'option
\textit{sortvanquished}.
\\
%.lp ""
Les catégories omises sont implicitement ajoutées avec le préfixe `\texttt{n}'.
Les catégories indiquées sans préfixe utilisent implicitement le préfixe `\texttt{+}'.
L'ordre des catégories n'a pas d'importance : l'affichage du programme
en fin de partie suit une séquence fixe.

%.lp ""
(par exemple ``\texttt{disclose:yi na +v -g o}'')
Cet exemple règle
\texttt{inventaire} sur \textit{demander} avec \textit{oui} par défaut,
\texttt{attributs} sur \textit{demander} avec \textit{non} par défaut,
\texttt{monstres vaincus} sur \textit{révéler sans demander},
\texttt{génocides} sur \textit{ne pas révéler} et \textit{ne pas demander},
\texttt{conduite} implicitement sur \textit{demander} avec \textit{non} par défaut,
\texttt{survol} sur \textit{révéler sans demander}.

%.lp ""
Notez que la liste des monstres vaincus comprend tous les monstres tués par
des pièges et entre eux, ainsi que ceux que vous avez tués.
Et le survol du donjon montre tous les niveaux que vous avez visités, mais ne
révèle pas ce que vous n'y aviez pas découvert.
%.lp
\item[dogname]
Donne un nom à votre chien de départ (par exemple ``\texttt{dogname:Fang}'').
Ne peut pas être réglée avec la commande `\texttt{O}'.
%.lp
\item[extmenu]
Modifie l'interface des commandes étendues pour faire apparaître un menu des commandes disponibles.
Elle reste compatible au niveau des touches avec l'interface traditionnelle, sauf qu'il n'est
pas nécessaire d'appuyer sur Entrée.
Elle est implémentée pour l'interface tty (désactivée par défaut).
.lp ""
Pour l'interface X11, qui utilise toujours un menu pour choisir une commande
étendue, elle détermine si le menu affiche toutes les commandes disponibles (activée)
ou seulement le sous-ensemble des commandes traditionnellement considérées comme
étendues (désactivée).
%.lp
\item[female]
Synonyme obsolète de ``\texttt{gender:female}''. Ne peut pas être réglée avec la
commande `\texttt{O}'.
%.lp
\item[fireassist]
Cette option détermine ce qui se passe lorsque vous tentez la commande `\texttt{f}' (tirer)
sans manier de lanceur approprié, comme un arc ou une fronde.
Si elle est activée, vous manierez automatiquement le lanceur. Activée par défaut.
%.lp
\item[fixinv]
La lettre d'inventaire d'un objet lui reste attachée quand il est posé (activée par défaut).
Si cette option est désactivée, poser un objet décale toutes les lettres d'inventaire restantes.
Persistante.
%.lp
\item[force\textunderscore invmenu]
Les commandes qui demandent un objet de l'inventaire affichent un menu au lieu
d'une question textuelle avec les lettres possibles. Désactivée par défaut.
%.lp
\item[fruit]
Donne à un fruit le nom de quelque chose que vous aimez manger (par exemple ``\texttt{fruit:mango}'')
(par défaut ``\texttt{moisissure gluante}''). Il s'agit au fond d'une fantaisie nostalgique que
\textit{NetHack} utilise de temps en temps. Vous devriez choisir quelque chose que vous
trouvez plus appétissant qu'une moisissure gluante. Les pommes, oranges, poires, bananes et
melons existent déjà dans \textit{NetHack} ; ne les utilisez donc pas.
%.lp
\item[gender]
Votre sexe de départ (\texttt{gender:male} ou \texttt{gender:female}).
Vous pouvez n'indiquer que la première lettre.
Bien que vous puissiez
encore indiquer votre sexe avec les options obsolètes
``\textit{male}'' et ``\textit{female}'',
c'est l'option ``\textit{gender}'' qui l'emporte.
Voir \textit{role}
pour savoir comment utiliser la négation afin d'exclure des choix.
%.lp ""
\\
Si \texttt{gender} n'est pas indiqué, il n'y a pas de valeur par défaut ;
le joueur sera interrogé, sauf si le rôle ou la race impose un sexe.
Ne peut pas être réglée avec la commande `\texttt{O}'. Persistante.
%.lp
\item[goldX]
Lors du filtrage des objets selon leur état béni/maudit (BUCX), indique s'il faut
traiter les pièces d'or comme \texttt{X} (état béni/maudit inconnu, si `on')
ou comme \texttt{U} (connues comme non maudites, si `off', la valeur par défaut).
L'or n'est jamais béni ni maudit, mais il n'est pas décrit comme ``non maudit'',
même quand l'option \textit{implicit\textunderscore uncursed} est `off'.
%.lp
\item[help]
Si davantage d'informations sont disponibles pour un objet examiné
avec la commande `\texttt{/}', vous demande si vous voulez les voir (activée par défaut).
Désactiver l'aide rend le simple examen des choses plus rapide, puisque vous n'êtes pas
interrompu par la question ``\texttt{Plus d'informations sur ... ?}'', mais cela signifie aussi que vous
risquez de manquer des informations intéressantes ou importantes. Persistante.
%.lp
\item[herecmd\textunderscore menu]
Avec une interface qui gère la souris, cliquer sur vous-même ou
à côté de vous affiche un menu des actions possibles pour cet emplacement.
Identique aux commandes ``\texttt{\#herecmdmenu}'' et ``\texttt{\#therecmdmenu}''.
%.lp
\item[hilite\textunderscore pet]
Distingue visuellement les familiers des animaux semblables (désactivée par défaut).
Le comportement de cette option dépend du type de fenêtrage utilisé.
En fenêtrage texte, on utilise souvent la mise en valeur du texte ou la vidéo inverse ;
avec les tuiles, un symbole de cœur est généralement affiché près des familiers.

%.lp ""
Avec l'interface tty ou curses, l'option \textit{petattr}
détermine comment mettre en valeur les familiers, et la régler active ou désactive
l'option \textit{hilite\textunderscore pet} selon le cas.
%.lp
\item[hilite\textunderscore pile]
Distingue visuellement les tas d'objets des objets isolés (désactivée par défaut).
Le comportement de cette option dépend du type de fenêtrage utilisé.
En fenêtrage texte, on utilise souvent la mise en valeur du texte ou la vidéo inverse ;
avec les tuiles, un petit symbole plus est généralement affiché à côté de l'objet
situé en haut du tas.
%.lp
\item[hitpointbar]
Affiche une barre graphique des points de vie derrière votre nom et votre titre dans
l'affichage d'état (désactivée par défaut).
\\
%.lp
L'interface ``curses'' la prend en charge même si la fonctionnalité de mise en valeur
de l'état a été désactivée lors de la compilation du programme.
Les interfaces ``tty'' et ``mswin'' (alias ``Windows GUI'') ne la
prennent en charge que si la mise en valeur de l'état est restée activée à la compilation.
Vous n'avez besoin de définir aucune règle de mise en valeur pour afficher
la barre.
S'il en existe une active pour les points de vie et qu'elle indique une couleur, cette
couleur sera utilisée pour la barre.
En revanche, si elle indique des attributs vidéo, ils seront ignorés au
profit de \textit{inverse}.
Pour tty et curses, \textit{blink} sera aussi utilisé si la valeur actuelle
des points de vie est inférieure ou égale au seuil des \textit{PV critiques}.
\\
%.lp
L'interface ``Qt'' prend aussi en charge hitpointbar, en dessinant
une barre pleine au-dessus du nom et du titre, avec des couleurs codées en dur.
(À l'heure où ces lignes sont écrites, activer la barre empêche involontairement
de redimensionner le panneau d'état.
Pour le redimensionner, utilisez la commande \texttt{\#optionsfull} pour désactiver
l'option \textit{hitpointbar}, redimensionnez pendant qu'elle est désactivée, puis
utilisez la même commande pour la réactiver.)
%.lp
\item[horsename]
Donne un nom à votre cheval de départ (par exemple ``\texttt{horsename:Trigger}'').
Ne peut pas être réglée avec la commande `\texttt{O}'.
%.lp
\item[ignintr]
Ignore les signaux d'interruption, y compris Break (désactivée par défaut). Persistante.
%.lp
\item[implicit\textunderscore uncursed]
Omet ``non maudit'' des descriptions d'objets quand on peut le déduire
d'autres éléments de la description (activée par défaut).
Persistante.

%.lp ""
Si vous utilisez la coloration des menus, vous voudrez peut-être la désactiver.
%.lp
\item[legacy]
Affiche un message d'introduction au début de la partie (activée par défaut).
Persistante.
%.lp
\item[lit\textunderscore corridor]
Affiche comme éclairées les cases de couloir vues grâce à la vision nocturne ou à une source
de lumière tenue par votre personnage (désactivée par défaut). Persistante.
%.lp
\item[lootabc]
Lorsqu'un menu sert à interagir avec un conteneur,
utilise les anciens raccourcis clavier `\texttt{a}', `\texttt{b}' et `\texttt{c}'
plutôt que les mnémoniques `\texttt{o}', `\texttt{i}' et `\texttt{b}' (désactivée par défaut).
Persistante.
%.lp
\item[mail]
Active la distribution du courrier pendant la partie (activée par défaut). Persistante.
%.lp
\item[male]
Synonyme obsolète de ``\texttt{gender:male}''. Ne peut pas être réglée avec la
commande `\texttt{O}'.
%.lp
\item[mention\textunderscore decor]
Signale le passage sur divers éléments du donjon comme les escaliers,
les fontaines ou les autels, qui ne sont d'ordinaire décrits que lorsqu'ils sont recouverts
par un ou plusieurs objets (désactivée par défaut). Persistante.
%.lp
\item[mention\textunderscore map]
Signale les changements des emplacements intéressants de la carte (désactivée par défaut).
%.lp
\item[mention\textunderscore walls]
Signale quand vous marchez contre un mur (désactivée par défaut). Persistante.
%.lp
\item[menucolors]
Active la coloration des lignes de menu (désactivée par défaut).
Voir ``\textit{Configurer les couleurs des menus}'' pour savoir comment configurer les couleurs.
%.lp
\item[menustyle]
Détermine la méthode utilisée quand vous devez choisir plusieurs objets (en
réponse à la commande \texttt{Drop} (alias \texttt{droptype}), par exemple).
La valeur indiquée doit être la première lettre de l'une des suivantes :
traditional, combination, full ou partial.
La valeur par défaut est \texttt{full}.
Persistante.
\\
%.lp ""
Traditional était la seule méthode disponible dans les toutes
premières versions ; elle consiste en une question demandant les caractères des classes d'objets,
suivie d'une question objet par objet pour tous les objets correspondant aux classes
sélectionnées.
Combination commence par une question demandant les classes d'objets
qui vous intéressent, puis affiche un menu des objets correspondants au lieu
de les proposer un par un.
Full affiche un menu des
classes d'objets au lieu d'une question par caractère, puis un menu des objets
correspondants à sélectionner.
(Choisir son option `A' (tout sélectionner automatiquement) saute le second menu.
Pour éviter de la choisir par accident,
réglez \textit{paranoid\textunderscore confirm:AutoAll} pour exiger une confirmation.)
Partial saute le filtrage par classe d'objets et
affiche immédiatement un menu de tous les objets.
\item[menu\textunderscore deselect\textunderscore all]
Touche pour désélectionner tous les éléments d'un menu.
Implémentée par les versions Amiga, Gem, X11 et tty.
Par défaut `-'.
\item[menu\textunderscore deselect\textunderscore page]
Touche pour désélectionner tous les éléments de cette page d'un menu.
Implémentée par les versions Amiga, Gem et tty.
Par défaut `\textbackslash'.
\item[menu\textunderscore first\textunderscore page]
Touche pour aller à la première page d'un menu.
Implémentée par les versions Amiga, Gem et tty.
Par défaut `\textasciicircum'.
\item[menu\textunderscore headings]
Détermine la mise en valeur des en-têtes d'un menu.
Accepte un attribut de texte, ou une couleur de texte et un attribut séparés par une esperluette.
Pour les attributs et couleurs autorisés, voir ``\textit{Configurer les couleurs des menus}''.
Toutes les versions ne peuvent pas afficher tous les types.
\item[menu\textunderscore invert\textunderscore all]
Touche pour inverser la sélection de tous les éléments d'un menu.
Implémentée par les versions Amiga, Gem, X11 et tty.
Par défaut `@'.
\item[menu\textunderscore invert\textunderscore page]
Touche pour inverser la sélection de tous les éléments de cette page d'un menu.
Implémentée par les versions Amiga, Gem et tty.
Par défaut `\textasciitilde'.
\item[menu\textunderscore last\textunderscore page]
Touche pour aller à la dernière page d'un menu.
Implémentée par les versions Amiga, Gem et tty.
Par défaut `|'.
\item[menu\textunderscore next\textunderscore page]
Touche pour passer à la page suivante du menu.
Implémentée par les versions Amiga, Gem et tty.
Par défaut `>'.
\item[menu\textunderscore objsyms]
% [originally menu_objsyms was a boolean]
% Show object symbols in menu headings in menus where
% the object symbols act as menu accelerators (default off).
L'inventaire et les autres menus d'objets sont normalement séparés par classe d'objets
(armes, armures, etc.), avec une ligne d'en-tête au début
de chaque groupe.
Vous pouvez faire ajouter par les menus le symbole d'affichage de la classe d'objets sur
chaque ligne d'en-tête.
Vous pouvez aussi ajouter le symbole d'affichage de chaque objet dans chaque entrée
de menu.
Pour une carte en tuiles, il s'agirait d'une petite reproduction de la tuile de l'objet.
Pour une carte en texte, c'est le même caractère que celui utilisé pour la
classe de l'objet, ce qui est surtout utile quand aucun en-tête ne sépare
les objets par classe.
% FIXME! -- three column table may be simpler than mashing first two together
% Possible values are
% .PS "5\ -\ One-or-other"
% .PL "0\ -\ None"
% no symbols for either header lines or menu entries;
% .PL "1\ -\ Headers"
% show symbols on header lines but not entries;
% .PL "2\ -\ Entries"
% show symbols on menu entry lines but not headers;
% .PL "3\ -\ Both"
% show symbols on headers and entries;
% .PL "4\ -\ Conditional"
% only show symbols for entries if there are no headers;
% .PL "5\ -\ One-or-other"
% show symbols on headers, or on entries if no headers.
% .PE

Prise en charge par tty et curses.
Lors du réglage, la valeur peut être indiquée par son chiffre ou par son mot-clé.
La valeur par défaut est \texttt{Conditional} (4).
\item[menu\textunderscore overlay]
N'efface pas l'écran avant de dessiner les menus, et aligne
les menus sur le bord droit de l'écran. Uniquement pour la version tty.
(activée par défaut)
\item[menu\textunderscore previous\textunderscore page]
Touche pour revenir à la page précédente du menu.
Implémentée par les versions Amiga, Gem et tty.
Par défaut `<'.
\item[menu\textunderscore search]
Touche pour rechercher un texte et inverser l'état de sélection des éléments de menu correspondants.
Par défaut `:'.
\item[menu\textunderscore select\textunderscore all]
Touche pour sélectionner tous les éléments d'un menu.
Implémentée par les versions Amiga, Gem, X11 et tty.
Par défaut `.'.
\item[menu\textunderscore select\textunderscore page]
Touche pour sélectionner tous les éléments de cette page d'un menu.
Implémentée par les versions Amiga, Gem et tty.
Par défaut `,'.

%.lp
\item[menu\textunderscore shift\textunderscore left]
Touche pour faire défiler vers la gauche un menu---qui a été
défilé vers la droite---afin de le ramener.
Implémentée pour \textit{perm\textunderscore invent} uniquement par curses et X11.
Par défaut `\texttt{\textbraceleft}'.

%.lp
\item[menu\textunderscore shift\textunderscore right]
Touche pour faire défiler vers la droite un menu dont le texte dépasse le
bord droit.
Implémentée pour \textit{perm\textunderscore invent} uniquement par curses et X11.
Par défaut `\texttt{\textbraceright}'.
% %.lp
% \item[menu\textunderscore tab\textunderscore sep]
% Format menu entries using TAB to separate columns (default off).
% Only applicable to some menus, and only useful to some interfaces.
% Debug mode only.
%.lp
\item[mon\textunderscore movement]
Affiche un message quand le héros remarque le déplacement d'un monstre (désactivée par défaut).
%.lp
\item[monpolycontrol]
Demande la nouvelle forme chaque fois qu'un monstre change de forme (désactivée par défaut).
Mode débogage uniquement.
%.lp
\item[montelecontrol]
Demande la destination chaque fois qu'un monstre est téléporté (désactivée par défaut).
Mode débogage uniquement.
%.lp
\item[mouse\textunderscore support]
Autorise l'utilisation de la souris pour la saisie et le voyage.
Les réglages valides sont :

%.sd
%.si
\texttt{0} --- désactivée\\
\texttt{1} --- activée, avec des ajustements du système pour gérer la souris\\
\texttt{2} --- comme \texttt{1}, mais sans aucun ajustement du système\\
%.ei
%.ed

Omettre la valeur équivaut à indiquer \texttt{1},
et nier
\textit{mouse\textunderscore support}
équivaut à indiquer \texttt{0}.
%.lp
\item[msghistory]
Le nombre de messages de la ligne du haut à conserver (et que l'on peut rappeler
avec `\texttt{\textasciicircum P}') (20 par défaut).
Ne peut pas être réglée avec la commande `\texttt{O}'.
%.lp
\item[msg\textunderscore window]
Permet de modifier la façon dont les messages rappelés sont affichés.
Actuellement prise en charge uniquement par tty (les quatre choix) et par curses
(choix `\texttt{f}' et `\texttt{r}', `\texttt{r}' par défaut).
Les valeurs possibles sont :

%.sd
%.si
\texttt{s} --- un seul message (par défaut ; seul choix avant la version 3.4.0) ;\\
\texttt{c} --- combinaison : deux messages en mode \textit{single}, puis en mode \textit{full} ;\\
\texttt{f} --- fenêtre complète, message le plus ancien en premier ;\\
\texttt{r} --- fenêtre complète inversée, message le plus récent en premier.
%.ei
%.ed

Pour la compatibilité ascendante, aucune valeur n'est nécessaire (ce qui
équivaut à \textit{full}), et l'option peut être niée (ce qui équivaut
à \textit{single}).
%.lp
\item[name]
Définit le nom de votre personnage (par défaut, votre nom d'utilisateur). Vous pouvez aussi
définir le rôle de votre personnage en ajoutant un tiret et une ou plusieurs lettres du
rôle (c'est-à-dire en ajoutant l'un des suffixes
``\texttt{-A -B -C -H -K -M -P -Ra -Ro -S -T -V -W}'').
Si ``\texttt{-@}'' est utilisé pour le rôle, un rôle aléatoire sera
automatiquement choisi.
%.lp
\\
Sur certains systèmes, la valeur par défaut est le nom d'utilisateur du joueur ;
sur d'autres, il n'y a pas de valeur par défaut et le joueur sera interrogé.
On peut obtenir le second comportement avec le premier en indiquant un nom générique
comme ``player''. Le réglage GENERICUSERS= de sysconf détermine quels
noms sont considérés comme génériques.
Ne peut pas être réglée avec la commande `\texttt{O}'.
%.lp
\item[news]
Lit le fichier de nouvelles de \textit{NetHack}, s'il existe (activée par défaut).
Comme les nouvelles sont affichées au début de la partie, il est inutile
de régler cette option avec la commande `\texttt{O}'.
%.lp
\item[nudist]
Fait commencer le personnage sans armure (désactivée par défaut). Persistante.
%.lp
\item[null]
Envoie des caractères nuls de remplissage au terminal (activée par défaut). Persistante.
%.lp
\item[number\textunderscore pad]
Utilise les touches numériques au lieu des lettres pour se déplacer (0 ou désactivée par défaut).\\
Les réglages valides sont :

%.sd
%.si
\newlength{\mwidth}
\settowidth{\mwidth}{\texttt{-0}}
\newcommand{\numbox}[1]{\makebox[\mwidth][r]{\texttt{#1}}}
\numbox{0} --- déplacement par lettres ; `\texttt{yuhjklbn}'\\
\numbox{1} --- déplacement par chiffres ; le chiffre `\texttt{5}' sert de préfixe de déplacement `\texttt{G}'\\
\numbox{2} --- comme \texttt{1}, mais `\texttt{5}' sert de préfixe `\texttt{g}' au lieu de `\texttt{G}'\\
\numbox{3} --- par chiffres avec la disposition d'un clavier de téléphone ; \texttt{123} en haut, \texttt{789} en bas\\
\numbox{4} --- combine \texttt{3} et \texttt{2} ; disposition téléphone plus compatibilité MS-DOS\\
\numbox{-1} --- par lettres, mais `\texttt{z}' pour aller au nord-ouest et `\texttt{y}' pour zapper les baguettes
%.ei
%.ed

Pour la compatibilité ascendante, omettre la valeur équivaut à indiquer \texttt{1},
et nier
\textit{number\textunderscore pad}
équivaut à indiquer \texttt{0}.
(Les réglages \texttt{2} et \texttt{4} servent à la compatibilité avec MS-DOS ou l'ancien PC Hack ;
en plus du comportement différent de `\texttt{5}', `\texttt{Alt-5}' agit comme `\texttt{G}'
et `\texttt{Alt-0}' comme `\texttt{I}'.
Le réglage \texttt{-1} est destiné à certains claviers QWERTZ sur lesquels les
touches `\texttt{y}' et `\texttt{z}' sont interverties.)
En déplacement par chiffres, pour saisir un préfixe de répétition pour les commandes
qui en acceptent un (comme ``\texttt{12s}'' pour fouiller douze fois), faites-le précéder
de la lettre `\texttt{n}' (``\texttt{n12s}'').
%.lp
\item[packorder]
Indique l'ordre dans lequel les types d'objets sont listés (par défaut
``\texttt{")[\%?+!=/(*\textasciigrave 0\textunderscore}''). La valeur de cette option doit être une chaîne
contenant les symboles des différents types d'objets. Les types omis
sont ajoutés à la fin, dans l'ordre précédent.
%.lp
\item[paranoid\textunderscore confirmation]
Une liste, séparée par des espaces, de situations précises dans lesquelles une
demande de confirmation différente est souhaitée.
La valeur par défaut est ``\textit{paranoid\textunderscore confirmation:pray swim trap}''.
%.sd
%.si
\newlength{\pcwidth}
\settowidth{\pcwidth}{\texttt{Were-change}}
\addtolength{\pcwidth}{\labelsep}
\begin{description}[leftmargin=\pcwidth, topsep=1mm, itemsep=0mm, labelwidth=*, font=\ttfamily, align=right]
\item[Confirm]
pour toutes les questions réglées pour exiger ``yes'' plutôt que `y',
exiger aussi ``no'' pour refuser, au lieu de considérer toute réponse autre que oui
comme un non ; pray et AutoAll exigent alors aussi ``yes'' ou ``no'' ;
\item[quit]
exiger ``\texttt{yes}'' plutôt que `\texttt{y}' pour confirmer l'abandon de
la partie ou le passage en mode exploration (sans score) ;
\item[die]
exiger ``\texttt{yes}'' plutôt que `\texttt{y}' pour confirmer la mort (sans
utilité en jeu normal ; s'applique au mode exploration) ;
\item[bones]
exiger ``\texttt{yes}'' plutôt que `\texttt{y}' pour confirmer la sauvegarde
des données d'ossements lors d'une mort en mode débogage ;
\item[attack]
exiger ``\texttt{yes}'' plutôt que `\texttt{y}' pour confirmer l'attaque
d'un monstre pacifique ;
\item[wand-break]
exiger ``\texttt{yes}'' plutôt que `\texttt{y}' pour confirmer le bris
d'une baguette avec la commande \textit{apply} ;
\item[eating]
exiger ``\texttt{yes}'' plutôt que `\texttt{y}' pour confirmer s'il faut
continuer à manger ;
\item[Were-change]
exiger ``\texttt{yes}'' plutôt que `\texttt{y}' pour confirmer un changement de forme dû
à la lycanthropie quand le héros contrôle sa métamorphose ;
\item[pray]
exiger `\texttt{y}' pour confirmer une tentative de prière au lieu
de prier immédiatement ; activé par défaut ;
(pour exiger ``yes'' plutôt que simplement `y', activez aussi Confirm) ;
\item[trap]
exiger `\texttt{y}' pour confirmer une tentative d'entrer dans un piège connu ou de marcher dessus,
sauf si c'est considéré comme sans danger ;
lorsqu'elle est activée, cette confirmation s'applique aussi à l'entrée dans des
régions de nuage de gaz visibles ;
(pour exiger ``yes'' plutôt que simplement `y', activez aussi Confirm) ;
la confirmation peut être évitée avec le préfixe de déplacement `\texttt{m}' ;
\item[swim]
empêcher d'entrer dans l'eau ou la lave ; activé par défaut ; (pour marcher
délibérément sur un tel terrain quand ce choix est actif, utilisez le préfixe
de déplacement `\texttt{m}' depuis une case adjacente) ;
\item[AutoAll]
exiger une confirmation quand le choix `A' (tout sélectionner automatiquement) est sélectionné
dans les menus de filtrage par classe d'objets avec \textit{menustyle:Full} ;
(pour exiger ``yes'' plutôt que simplement `y', activez aussi Confirm) ;
\item[Remove]
exiger une sélection dans l'inventaire pour les commandes `\texttt{R}' et `\texttt{T}'
même quand un seul objet concerné est porté ;
\item[all]
activer tous les choix ci-dessus.
\end{description}
%.ei
%.ed
Par défaut, les choix pray, swim et trap sont activés et les autres désactivés.
Pour les désactiver sans activer
aucun autre choix, utilisez ``\textit{paranoid\textunderscore confirmation:none}''.
Pour les garder activés tout en activant d'autres choix, vous pouvez
les inclure dans la liste, par exemple
``\textit{par\-a\-noid\textunderscore con\-fir\-ma\-tion:attack~pray~swim~Remove}'',
ou faire précéder la première entrée de la liste d'un signe plus,
``\textit{paranoid\textunderscore confirmation:+attack~Remove}''.
Pour retirer une entrée précédemment activée sans retirer les autres,
faites précéder la première entrée de la liste d'un signe moins,
``\textit{paranoid\textunderscore confirmation:-swim}''.
Pour à la fois ajouter de nouvelles entrées et en retirer d'anciennes, vous pouvez utiliser
plusieurs réglages de l'option \textit{paranoid\textunderscore confirmation}, ou bien
utiliser la forme `\texttt{+}' et lister les entrées à ajouter par leur nom
et les entrées à retirer par `\texttt{!}' suivi du nom.
Les entrées positives (sans `!') et négatives (avec `!')
peuvent être mélangées.
%.lp
\item[pauper]
Fait commencer le personnage sans aucune possession (désactivée par défaut). Persistante.
%.lp ""
\\
Il commence aussi sans sorts ni compétences, ceux-ci étant liés à l'équipement de départ.
N'empêche pas d'acquérir et d'utiliser objets, sorts et compétences une fois la partie
commencée.
%.lp
\item[perm\textunderscore invent]
Si elle est activée, affiche en permanence votre inventaire actuel dans une fenêtre (désactivée par défaut).
%.lp ""
\\
Cela n'a de sens
que pour les interfaces à fenêtres qui implémentent cette fonctionnalité.
Pour celles-ci, l'option
\texttt{perminv\textunderscore mode}
permet d'affiner ce qui est affiché
pour \textit{perm\textunderscore invent}.
Lui donner une valeur autre que \textit{none}
alors que \textit{perm\textunderscore invent} est désactivée activera cette dernière.
%.lp
\item[perminv\textunderscore mode]
Complète l'option
\texttt{perm\textunderscore invent}.
La valeur est l'une des suivantes :
%.PS "\f(CRin-use\fP"
\settowidth{\pcwidth}{\texttt{in-use}}  %reuse the paranoid_confirm width
\addtolength{\pcwidth}{\labelsep}
\begin{description}[leftmargin=\pcwidth, topsep=1mm, itemsep=0mm, labelwidth=*, font=\mdseries, align=right]
%.PL
\item[\texttt{none}]
se comporter comme si \textit{perm\textunderscore invent} était désactivée ;
\item[{all}]
afficher tout l'inventaire sauf l'or ;
\item[{full}]
afficher l'inventaire complet, or compris ;
\item[{in-use}]
n'afficher que les objets utilisés (portés, maniés, lampe allumée).
%.PE
\end{description}
La valeur par défaut est \textit{none}, mais si \textit{perm\textunderscore invent} est activée
alors qu'elle vaut \textit{none}, elle passera à \textit{all}.
%.lp ""
\\
Remarque : si de l'or a été placé dans le carquois ou la bourse à munitions, il sera
inclus pour \textit{all}, bien que ce mode omette normalement l'or.
%.lp
%.\" petattr is a wincap option but we'll document it here...
\item[petattr]
Indique un ou plusieurs attributs de mise en valeur du texte à utiliser pour afficher
les familiers sur la carte.
Il s'agit en pratique d'une extension de l'option booléenne \textit{hilite\textunderscore pet}.
Interfaces curses ou tty uniquement ; la valeur est l'une de
none, bold, dim, underline, italic, blink et inverse.
Certains de ces choix peuvent ne pas fonctionner,
selon le matériel du terminal ou le logiciel d'émulation de terminal.

%.lp
\item[pettype]
Indique le type de votre familier initial, si vous jouez une classe de personnage
qui utilise plusieurs types de familiers ; ou choisissez de n'avoir aucun familier initial.
Les valeurs possibles sont ``\texttt{cat}'', ``\texttt{dog}'', ``\texttt{horse}''
et ``\texttt{none}''.
Si le choix n'est pas autorisé pour le rôle que vous jouez actuellement,
il sera ignoré sans avertissement. Par exemple, ``\texttt{horse}'' ne sera
respecté que si vous jouez un chevalier.
Ne peut pas être réglée avec la commande `\texttt{O}'.
%.lp
\item[pickup\textunderscore burden]
Lorsque vous ramassez un objet qui vous ferait dépasser ce niveau
de charge (Unencumbered, Burdened, streSsed, straiNed, overTaxed
ou overLoaded, soit non chargé, Chargé, Stressé, Tendu, Exténué ou Surchargé),
on vous demandera si vous voulez continuer.
(`S' par défaut). Persistante.
%.lp
\item[pickup\textunderscore stolen]
Si cette option est activée et que ``\textit{autopickup}'' l'est aussi, tente de ramasser
les objets qu'un monstre vous a volés, même s'ils ne figurent pas dans
``\textit{pickup\textunderscore types}'' ou
ne correspondent à aucune exception de ramassage automatique.
Activée par défaut.
Persistante.
%.lp
\item[pickup\textunderscore thrown]
Si cette option est activée et que ``\textit{autopickup}'' l'est aussi, tente de ramasser
les objets que vous avez lancés, même s'ils ne figurent pas dans
``\textit{pickup\textunderscore types}'' ou
ne correspondent à aucune exception de ramassage automatique.
Activée par défaut.
Persistante.
%.lp
\item[pickup\textunderscore types]
Indique les types d'objets à ramasser quand ``\textit{autopickup}''
est activée.
Par défaut, tous les types.
Persistante.
\\
%.lp ""
La valeur est une liste de symboles d'objets, par exemple
\texttt{pickup\textunderscore types:\textdollar ?!\&} pour ramasser l'or, les parchemins et les potions.
Vous pouvez utiliser des lignes
``\textit{autopickup\textunderscore exception}''
du fichier de configuration pour affiner encore le comportement de ``\textit{autopickup}''.
\\
%.lp ""
Il n'existe aucun moyen de régler \textit{pickup\textunderscore types} sur ``\textit{aucun}''.
(Lui donner une valeur vide revient à ``\textit{tous}''.)
Si vous voulez éviter de ramasser automatiquement le moindre type d'objet tout en
gardant \textit{autopickup} activée afin que les réglages
\textit{autopickup\textunderscore exceptions} déterminent ce que vous ramassez ou non,
vous pouvez régler \textit{pickup\textunderscore types} sur `\texttt{.}'.
C'est le symbole du type \textit{venin}, et comme vous ne croiserez
aucun objet de type venin, vous n'en ramasserez pas par mégarde.
%.lp
\item[pile\textunderscore limit]
Lorsque vous traversez un tas d'objets sur le sol, seuil à partir duquel
le message ``Il y a quelques/plusieurs/de nombreux objets ici.'' est donné au lieu
d'afficher une liste de ces objets. Une valeur de 0 signifie ``aucune limite''
(toujours lister les objets) ; une valeur de 1 signifie en pratique ``ne jamais montrer
les objets'', puisque le tas aura toujours au moins cette taille ;
la valeur par défaut est 5. Persistante.
%.lp
\item[playmode]
Les valeurs sont \textit{normal}, \textit{explore} ou \textit{debug}.
Permet de choisir le mode exploration (aussi appelé mode découverte) ou le mode
débogage (aussi appelé mode sorcier) au lieu du jeu normal.
Le mode débogage peut n'être autorisé que pour une personne connectée sous un
nom d'utilisateur particulier (sur les systèmes multi-utilisateurs) ou indiquant un nom
de personnage particulier (sur les systèmes mono-utilisateur), ou être entièrement désactivé.
Le demander quand il n'est pas autorisé ou pas possible mène au mode exploration à la place.
Par défaut, jeu normal.
%.lp
\item[price\textunderscore quotes]
Chaque fois que le jeu mentionne le nom d'un objet que vous n'avez pas encore identifié,
il mentionne aussi la fourchette des prix d'achat et de vente que vous avez vus pour cet
objet (pour vous aider à deviner ce qu'il peut être).
Le prix affiché est le prix unitaire d'un objet (même si vous regardez
une pile de plusieurs objets).
Beaucoup de joueurs voudront l'activer pendant l'identification des objets, puis
la désactiver pour le jeu courant.
Désactivée par défaut.
%.lp
\item[pushweapon]
Utiliser la commande `\texttt{w}' (manier) alors que vous maniez déjà
quelque chose place l'ancien objet dans votre emplacement d'arme secondaire (désactivée par défaut).
De même pour la commande `\texttt{a}' (appliquer) si elle fait manier l'objet
appliqué. Persistante.
%.lp
\item[query\textunderscore menu]
Utilise un menu, au lieu d'une question, pour certaines questions oui/non.
%.lp
\item[quick\textunderscore farsight]
Lorsqu'elle est activée, empêche généralement le message ``Vous percevez les alentours.'',
avec lequel le jeu fait une pause pour vous laisser parcourir la carte chaque fois que la
clairvoyance s'active au hasard.
Certaines situations, comme être sous l'eau ou englouti, ignorent cette option.
Elle n'affecte pas le sort de clairvoyance, pour lequel la pause permettant d'examiner les objets
ou monstres révélés est moins gênante.
Désactivée par défaut. Persistante.
%.lp
\item[race]
Les choix sont \texttt{human}, \texttt{dwarf}, \texttt{elf}, \texttt{gnome} et
\texttt{orc}, mais la plupart des rôles restreignent les races non humaines autorisées.
Voir \textit{role}
pour savoir comment utiliser la négation afin d'exclure des choix.
%.lp ""
\\
Si \texttt{race} n'est pas indiqué, il n'y a pas de valeur par défaut ;
le joueur sera interrogé, sauf si le rôle impose une race.
Ne peut pas être réglée avec la commande `\texttt{O}'. Persistante.
%.lp
\item[reroll]
Permet de retirer au sort l'inventaire de départ et les attributs de votre personnage (désactivée
par défaut). Persistante.
%.lp ""
\\
Notez que retirer votre personnage n'est pas une façon recommandée de jouer si vous visez
simplement la victoire (un départ chanceux influe bien moins sur votre victoire
que vos actions ultérieures dans la partie). Cette option existe
en partie parce que l'on reconnaît que certains joueurs insisteront pour le faire de toute façon,
et en partie parce que retirer peut être nécessaire pour certains types de parties
à défi.
%.lp
\item[rest\textunderscore on\textunderscore space]
Fait de la barre d'espace un synonyme de la commande `\texttt{.}' (\#wait) (désactivée par défaut).
Persistante.
%.lp
\item[role]
Choisit votre type de personnage (par exemple ``\texttt{role:Samurai}'') ;
synonyme de ``\textit{character}''.
Voir ``\textit{name}'' pour une autre méthode permettant d'indiquer votre rôle.
%.\" Normally only the first letter of the
%.\" value is examined; `r' is an exception with ``\texttt{Rogue}'',
%.\" ``\texttt{Ranger}'', and ``\texttt{random}'' values.
%.lp ""
Cette option peut aussi servir à limiter le choix quand le rôle est tiré
au hasard.
Utilisez une liste de rôles séparés par des espaces, et niez chacun d'eux ou niez
plutôt l'option elle-même.
La négation s'effectue de la même manière que pour les \textit{options booléennes},
en faisant précéder l'option ou sa ou ses valeurs de `\texttt{!}' ou de ``\texttt{no}''.
%.BR 0
\\
Exemples :
\\
\begin{verbatim}
OPTIONS=role:!arc !bar !kni
OPTIONS=!role:arc bar kni
\end{verbatim}
Il peut y avoir plusieurs occurrences de l'option \textit{role}
si ce sont toutes des négations.
%.\" Only one positive value is allowed, and if present, it overrides any
%.\" negations.
%.lp ""
\\
Si \texttt{role} n'est pas indiqué, il n'y a pas de valeur par défaut ;
le joueur sera interrogé.
Ne peut pas être réglée avec la commande `\texttt{O}'. Persistante.
%.lp
\item[roguesymset]
Cette option permet de choisir l'un des jeux de symboles nommés figurant dans
\texttt{symbols} pour modifier les symboles affichés à l'écran sur le
niveau rogue.
%.lp
\item[rlecomp]
Lors de l'écriture d'un fichier de sauvegarde, compresse la carte par plages (run length).
Toutes les versions ne prennent pas en charge la compression par plages. Cela n'a aucun
effet sur la lecture d'un fichier de sauvegarde existant.
%.lp
\item[runmode]
Détermine la fréquence de mise à jour de la fenêtre de la carte pendant un
déplacement sur plusieurs tours (course avec \texttt{Maj}+direction
ou \texttt{Ctrl}+direction,
etc., ou avec la commande de voyage ou un clic de souris).
Les valeurs possibles sont :

%.sd
%.si
\texttt{teleport} --- mettre à jour la carte une fois le déplacement terminé ;\\
\texttt{run} --- mettre à jour la carte tous les sept pas environ ;\\
\texttt{walk} --- mettre à jour la carte après chaque pas ;\\
\texttt{crawl} --- comme \textit{walk}, mais avec une brève pause après chaque pas.
%.ei
%.ed

Cette option n'affecte que l'affichage du jeu à l'écran, pas le
résultat réel du déplacement. La valeur par défaut est \textit{run} ; les versions antérieures à 3.4.1
n'utilisaient que \textit{teleport}. Que l'effet soit perceptible ou non
dépend de l'interface utilisée ou du type de terminal. Persistante.
%.lp
\item[safe\textunderscore pet]
Vous empêche d'attaquer (sciemment) vos familiers (activée par défaut). Persistante.
%.lp
\item[safe\textunderscore wait]
Vous empêche d'attendre ou de fouiller quand vous êtes à côté d'un monstre hostile
(activée par défaut). Persistante.
%.lp
\item[sanity\textunderscore check]
Vérifie les monstres, les objets et la carte avant chaque tour (désactivée par défaut).
Mode débogage uniquement.
%.lp
\item[scores]
Détermine quelles parties de la liste des scores vous sont montrées à la fin (par exemple
``\texttt{scores:5top scores/4around my score/own scores}''). Seule la première
lettre de chaque catégorie (`\texttt{t}', `\texttt{a}' ou `\texttt{o}') est nécessaire.
Persistante.
%.lp
\item[showdamage]
Chaque fois que votre personnage subit des dégâts, affiche un message indiquant les dégâts subis
et le nombre de points de vie restants.
%.lp
\item[showexp]
Affiche vos points d'expérience accumulés sur la ligne du bas (désactivée par défaut).
Persistante.
%.lp
\item[showrace]
Vous affiche avec le glyphe de votre race plutôt qu'avec celui
de votre rôle (désactivée par défaut). Notez que ce réglage n'affecte que
l'apparence de l'affichage, pas la façon dont le jeu vous traite.
Persistante.
%.lp
\item[showscore]
Affiche votre score approximatif accumulé sur la ligne du bas (désactivée par défaut).
Par défaut, cette fonctionnalité est supprimée lors de la compilation du programme.
Persistante.
%.lp
\item[showvers]
Inclut le numéro de version du jeu sur les lignes d'état (désactivée par défaut).
Potentiellement utile si vous passez d'une version ou d'une variante à l'autre,
ou si vous faites des captures d'écran ou diffusez des vidéos.
L'utilisation de l'option
\textit{statuslines:3}
est recommandée afin de disposer de plus de place pour les
informations d'état, sauf si vous utilisez l'interface \textit{Qt} de NetHack
ou si la fenêtre de votre émulateur de terminal affiche moins de 25 lignes.
Persistante.
\\
%.lp ""
L'option \textit{versinfo}
offre un contrôle limité sur les informations
qu'affiche \textit{showvers}.
%.lp
\item[silent]
Supprime les bips du terminal (activée par défaut). Persistante.
%.lp
\item[sortdiscoveries]
Détermine le tri de l'affichage des commandes `\texttt{\textbackslash}'
et `\texttt{\`{}}'.
Persistante.
\\
%.lp ""
Les valeurs possibles sont :
\\
%.PS
\texttt{o} --- lister les types d'objets par classe, dans l'ordre de découverte au sein de chaque classe ;
par défaut ;
\\
\texttt{s} --- lister les types d'objets selon la classification de \textit{sortloot} :
par classe, puis par sous-classe au sein des classes qui
comportent des regroupements importants (comme casques, bottes, gants, etc.
pour les armures), les types d'objets partiellement découverts par un nom attribué
venant avant les types entièrement identifiés ;
\\
\texttt{c} --- lister par classe, par ordre alphabétique au sein de chaque classe ;\\
\texttt{a} --- lister par ordre alphabétique, toutes classes confondues.\\
%.PE
Peut être réglée interactivement avec la commande `\texttt{O}' ou en utilisant
le préfixe `\texttt{m}' devant la commande `\texttt{\textbackslash}'
ou `\texttt{\textasciigrave}'.
%.lp
\item[sortloot]
Détermine le tri des listes de ramassage pour l'inventaire,
la commande \#loot et quelques autres. Persistante.
\\
Les valeurs possibles sont :
\\
%.sd
%.si
\texttt{full} --- toujours trier les listes ;\\
\texttt{loot} --- ne trier que les listes qui n'utilisent pas les lettres
       d'inventaire, comme avec les commandes \#loot et de ramassage ;\\
\texttt{none} --- afficher les listes de manière traditionnelle, sans tri ; par défaut.
%.ei
%.ed
%.lp
\item[sortpack]
Trie le contenu du sac par type lors de l'affichage de l'inventaire (activée par défaut).
Persistante.
%.lp
\item[sortvanquished]
Détermine le tri de l'affichage de la commande \texttt{\#vanquished}
ainsi que de la commande \texttt{\#genocided}.
Persistante.
\\
%.lp ""
Les valeurs possibles sont :
\\
%.PS
\texttt{t} ---
traditionnel : tri par niveau du monstre ; les égalités sont départagées par l'index
interne du monstre ;
par défaut ;
\\
\texttt{d} ---
tri par indice de difficulté du monstre ; égalités départagées par l'index interne ;
\\
\texttt{a} ---
tri alphabétique, d'abord les monstres uniques, puis tous les autres ;
\\
%note: 'A' and 'C' can be set in RC file or NETHACKOPTIONS but not by 'O'
% \texttt{A} ---
% order alphabetically, unique monsters intermixed with other monsters;
% \\
% \texttt{C} ---
% order by monster class, by high to low level within each class;
% \\
\texttt{c} ---
tri par classe de monstre, par niveau croissant au sein de chaque classe ;
\\
\texttt{n} ---
tri par nombre, décroissant ; les égalités sont départagées par l'index interne du monstre ;
\\
\texttt{z} ---
tri par nombre, croissant ; égalités départagées par l'index interne.
\\
%.PE
Peut être réglée interactivement avec la commande `\texttt{m O}' ou en utilisant
le préfixe `\texttt{m}' devant la commande \texttt{\#vanquished}
ou la commande \texttt{\#genocided}.
%.lp
\item[sounds]
Autorise l'émission de sons par une bibliothèque sonore intégrée (activée par défaut).
%.lp
\item[sparkle]
Affiche un effet scintillant quand un monstre (vous compris) est touché par une
attaque à laquelle il résiste (activée par défaut). Persistante.
%.lp
\item[spot\textunderscore monsters]
Affiche un message quand le héros remarque un monstre (désactivée par défaut).
%.lp
\item[standout]
Affiche en gras les monstres et ``\texttt{--Plus--}'' (désactivée par défaut). Persistante.
%.lp
\item[statushilites]
Détermine pendant combien de tours les comportements de mise en valeur de l'état mettent
le champ en valeur. Si elle est niée ou réglée sur zéro, la mise en valeur de l'état est désactivée.
Voir ``\textit{Configurer la mise en valeur de l'état}'' pour plus d'informations.
%.lp
\item[status\textunderscore updates]
Autorise les mises à jour des lignes d'état en bas de l'écran (activée par défaut).
%.lp
\item[suppress\textunderscore alert]
Cette option peut être réglée sur un numéro de version de \textit{NetHack} pour supprimer
les messages d'alerte concernant les changements de fonctionnalités de cette version
et des précédentes (par exemple ``\texttt{suppress\textunderscore alert:3.3.1}'')
%.lp
\item[symset]
Cette option permet de choisir l'un des jeux de symboles nommés figurant dans
\texttt{symbols} pour modifier les symboles affichés à l'écran.
Utilisez ``\texttt{symset:default}'' pour choisir explicitement les symboles par défaut.
%.lp
\item[terrainstatus]
Affiche un état supplémentaire décrivant l'endroit situé sous les pieds du
héros (désactivée par défaut, non prise en charge par toutes les interfaces).
%.lp
\item[time]
Affiche le temps de jeu écoulé, en tours, sur la ligne du bas (désactivée par défaut). Persistante.
%.lp
\item[timed\textunderscore delay]
Lors des brèves pauses destinées aux effets d'affichage, comme les explosions et
les objets en mouvement, utilise une minuterie plutôt que d'envoyer des caractères
supplémentaires à l'écran. (S'applique uniquement aux interfaces ``tty'' et ``curses'' ; l'interface ``X11''
utilise toujours un délai fondé sur une minuterie. Activée par défaut si elle a été configurée
dans le programme.) Persistante.
%.lp
\item[tips]
Affiche quelques conseils utiles pendant la partie (activée par défaut). Persistante.
%.lp
\item[tombstone]
Dessine une pierre tombale à votre mort (activée par défaut). Persistante.
%.lp
\item[toptenwin]
Place l'affichage de fin de partie dans une fenêtre de \textit{NetHack} plutôt que sur stdout (désactivée par défaut).
Activer cette option rend la liste des scores visible quand une version à fenêtres
de \textit{NetHack} est lancée sans fenêtre parente, mais la liste des scores ne reste
plus affichée après la fin de la partie dans un terminal ou une fenêtre d'émulation.
%.lp
\item[travel]
Autorise la commande de voyage par clic de souris (activée par défaut).
Désactiver cette option empêche le jeu de tenter des déplacements non voulus
si vous cliquez par inadvertance dans la fenêtre de la carte.
N'affecte pas le voyage avec la commande `\texttt{\textunderscore}' (``\texttt{\#travel}'').
Persistante.
% %.lp
% \item[ib{travel\textunderscore debug}]
% Display intended path during each step of travel (default off).
% Debug mode only.
%.lp
\item[tutorial]
Joue un niveau didacticiel au début de la partie.
Activer ou désactiver cette option dans le fichier de configuration évite la question.
%.lp
\item[verbose]
Fournit davantage de commentaires pendant la partie (activée par défaut). Persistante.
%.lp
\item[versinfo]
Détermine ce qu'affiche l'option \textit{showvers} sur les lignes d'état.
%.lp
\item[weaponstatus]
Affiche un état supplémentaire qui décrit l'arme actuellement maniée
(désactivée par défaut, non prise en charge par toutes les interfaces).
%.lp ""
\\
Quelques valeurs pouvant être affichées :
\\
\texttt{Mains-nues} - ni arme ni gants ;\\
\texttt{Mains-vides} - pas d'arme, mais des gants portés ;\\
\texttt{Deux-armes} - maniement de deux armes.
%.lp
\item[whatis\textunderscore coord]
Lorsque vous utilisez les commandes `\texttt{/}' ou `\texttt{;}' pour examiner la carte avec
``\texttt{autodescribe}''
activée, affiche les coordonnées après la description.
Fonctionne aussi dans les autres situations où l'on vous demande de choisir une position.\\

%.lp ""
Les réglages possibles sont :

%.sd
%.si
\texttt{c} --- \texttt{boussole ('est', '3s' ou '2n,4o')} ;\\
\texttt{f} --- \texttt{boussole complète ('est', '3sud' ou '2nord,4ouest')} ;\\
\texttt{m} --- \texttt{carte <x,y> (la colonne x=0 de la carte n'est pas utilisée)} ;\\
\texttt{s} --- \texttt{écran [ligne,colonne] (la ligne est décalée pour correspondre à l'usage de tty)} ;\\
\texttt{n} --- \texttt{aucun (pas de coordonnées affichées) [par défaut]}.
%.ei
%.ed

%.lp ""
L'option
\textit{whatis\textunderscore coord}
est aussi utilisée avec
les sous-commandes `\texttt{/m}', `\texttt{/M}', `\texttt{/o}' et `\texttt{/O}'
de `\texttt{/}',
où le réglage `\textit{none}' est remplacé par `\textit{map}'.
%.lp
\item[whatis\textunderscore filter]
Lorsqu'il faut indiquer une position sur la carte et que vous utilisez les touches pour
passer à la cible suivante ou précédente, permet de filtrer les cibles possibles.
(aucun filtrage par défaut)\\
%.lp ""
Les réglages possibles sont :

%.sd
%.si
\texttt{n} --- \texttt{aucun filtrage} ;\\
\texttt{v} --- \texttt{seulement ce qui est en vue} ;\\
\texttt{a} --- \texttt{seulement dans la même zone (pièce, couloir, etc.)}.
%.ei
%.ed
%.lp ""
Le filtre par zone essaie d'être légèrement
prédictif---si
vous vous tenez dans l'embrasure d'une porte, il considère la zone du côté
de la porte vers lequel vous vous déplaciez en dernier.\\
%.lp ""
Le filtrage peut aussi être modifié lors de l'indication d'une position avec
la touche ``getpos.filter''.
%.lp
\item[whatis\textunderscore menu]
Lorsqu'il faut indiquer une position sur la carte et que vous utilisez une touche pour
passer à la cible suivante ou précédente, utilise plutôt un menu pour choisir une cible.
(désactivée par défaut)
%.lp
\item[whatis\textunderscore moveskip]
Lorsqu'il faut indiquer une position sur la carte et que vous utilisez les touches de déplacement
avec Maj ou les touches méta+chiffre pour vous déplacer rapidement, au lieu d'avancer de 8 cases à la fois,
avance en sautant les glyphes identiques.
(désactivée par défaut)
%.lp
\item[windowtype]
Lorsque le programme a été compilé pour prendre en charge plusieurs interfaces,
choisit celle à utiliser, comme ``\texttt{tty}'' ou ``\texttt{X11}''
(la valeur par défaut dépend des réglages de compilation ; utilisez ``\texttt{\#version}'' pour la connaître).
Ne peut pas être réglée avec la commande `\texttt{O}'.

%.lp ""
Si elle est utilisée, ce doit être la première option réglée, car sa valeur peut
activer ou désactiver la disponibilité de diverses autres options.
Pour plusieurs lignes dans un fichier de configuration, ce serait la première
ligne qui n'est pas un commentaire.
Pour une liste séparée par des virgules dans NETHACKOPTIONS ou dans une ligne OPTIONS d'un
fichier de configuration, ce serait l'option la plus à \textit{droite} de la liste.
%.lp
\item[wizweight]
Ajoute le poids des objets à leur description (désactivée par défaut).
Mode débogage uniquement.
%.lp
\item[zerocomp]
Lors de l'écriture d'un fichier de sauvegarde, compresse le contenu par
suppression des zéros (zero-comp). Toutes les versions ne prennent pas en charge cette compression.
Cela n'a aucun effet sur la lecture d'un fichier de sauvegarde existant.
\end{description}

%.hn 2
\subsection*{Options de personnalisation de l'interface (window port)}

%.pg
Voici l'explication des diverses options qui servent
à personnaliser et à modifier les caractéristiques du
type d'interface (windowtype) que vous avez choisi.
Les chaînes de caractères trop longues peuvent être tronquées.
Toutes les interfaces ne tiennent pas compte de tous les réglages
listés ici.  Vous pouvez sans crainte ajouter n'importe laquelle de ces
options à votre fichier de configuration : si l'interface est capable de
s'adapter à vos préférences, elle essaiera de le faire. Sinon,
elle l'ignorera silencieusement.  Pour savoir si une
option est prise en charge par l'interface que vous utilisez
actuellement, vérifiez si elle apparaît dans la liste des options.
Certaines options sont dynamiques et peuvent être modifiées en cours
de partie avec la commande `\texttt{O}'.

\begin{description}[font=\mdseries\itshape]
%.lp
\item[align\textunderscore message]
 Où aligner ou placer la fenêtre des messages (top, bottom, left ou right)
%.lp
\item[align\textunderscore status]
 Où aligner ou placer la fenêtre d'état (top, bottom, left ou right).
%.lp
\item[ascii\textunderscore map]
%.hw DECgraphics IBMgraphics \% don't hyphenate these
\hyphenation{DECgraphics IBMgraphics}
Si \textit{NetHack} le peut, il affiche la carte avec de simples
caractères (lettres et ponctuation) plutôt qu'avec des graphismes en \textit{tuiles}.
Dans certains cas, ces caractères peuvent être complétés par des symboles de
tracé de lignes ; utilisez l'option \texttt{symset}
pour choisir un jeu de symboles tel que \textit{DECgraphics}
ou \textit{IBMgraphics} si votre affichage les prend en charge.
Régler \texttt{ascii\textunderscore map} à \textit{True} force
\texttt{tiled\textunderscore map} à \textit{False}.
%.lp
\item[color]
Si \textit{NetHack} le peut, il affiche en couleur les différents monstres,
objets et éléments du donjon (activée par défaut).
%.lp
\item[eight\textunderscore bit\textunderscore tty]
Si \textit{NetHack} le peut, il transmet tels quels à votre terminal les caractères sur huit bits (par exemple ceux indiqués avec
l'option \textit{traps}) (désactivée par défaut).
%.lp
\item[font\textunderscore map]
Si \textit{NetHack} le peut, il utilise la police du nom choisi pour la
fenêtre de la carte.
%.lp
\item[font\textunderscore menu]
Si \textit{NetHack} le peut, il utilise la police du nom choisi pour les fenêtres
de menu.
%.lp
\item[font\textunderscore message]
Si \textit{NetHack} le peut, il utilise la police du nom choisi pour la fenêtre des messages.
%.lp
\item[font\textunderscore status]
Si \textit{NetHack} le peut, il utilise la police du nom choisi pour la fenêtre d'état.
%.lp
\item[font\textunderscore text]
Si \textit{NetHack} le peut, il utilise la police du nom choisi pour les fenêtres de texte.
%.lp
\item[font\textunderscore size\textunderscore map]
Si \textit{NetHack} le peut, il utilise cette taille de police pour la fenêtre de la carte.
%.lp
\item[font\textunderscore size\textunderscore menu]
Si \textit{NetHack} le peut, il utilise cette taille de police pour les fenêtres de menu.
%.lp
\item[font\textunderscore size\textunderscore message]
Si \textit{NetHack} le peut, il utilise cette taille de police pour la fenêtre des messages.
%.lp
\item[font\textunderscore size\textunderscore status]
Si \textit{NetHack} le peut, il utilise cette taille de police pour la fenêtre d'état.
%.lp
\item[font\textunderscore size\textunderscore text]
Si \textit{NetHack} le peut, il utilise cette taille de police pour les fenêtres de texte.
%.lp
\item[fullscreen]
Si \textit{NetHack} le peut, il essaie d'utiliser tout l'écran plutôt qu'une fenêtre.
%.lp
\item[guicolor]
Utilise du texte en couleur et/ou des attributs de mise en évidence pour
afficher certaines données hors carte (comme les lettres de sélection des menus).
Interface curses uniquement ; activée par défaut.
%.lp
\item[large\textunderscore font]
Si \textit{NetHack} le peut, il utilise une grande police.
%.lp
\item[map\textunderscore mode]
Si \textit{NetHack} le peut, il affiche la carte de la manière indiquée.
%.lp
\item[player\textunderscore selection]
Si \textit{NetHack} le peut, il ouvre des boîtes de dialogue, ou pose des questions, pour le choix du personnage.
%.lp
\item[popup\textunderscore dialog]
Si \textit{NetHack} le peut, il ouvre des boîtes de dialogue pour les saisies.
%.lp
\item[preload\textunderscore tiles]
Si \textit{NetHack} le peut, il précharge les tuiles en mémoire.
Par exemple, dans la version MS-DOS en mode protégé, cette option détermine
si les tuiles sont chargées en RAM au début de la partie.  Cela
améliore les performances de l'affichage en tuiles, mais consomme davantage
de mémoire (activée par défaut).
Ne peut pas être modifiée avec la commande `\texttt{O}'.
%.lp
\item[scroll\textunderscore amount]
Si \textit{NetHack} le peut, il fait défiler l'affichage de ce nombre de cases
lorsque le héros atteint la marge scroll\textunderscore margin.
%.lp
\item[scroll\textunderscore margin]
Si \textit{NetHack} le peut, il fait défiler l'affichage lorsque le héros ou le curseur
se trouve à ce nombre de cases du bord de la fenêtre.
%.lp
\item[selectsaved]
Si \textit{NetHack} le peut, il affiche au lancement un menu des parties sauvegardées
existantes parmi lesquelles le joueur peut choisir. Toutes les versions ne prennent pas en charge cette option.
%.lp
\item[softkeyboard]
Si \textit{NetHack} le peut, il affiche un clavier à l'écran.
Ce sont surtout les appareils portables qui prennent en charge cette option.
%.lp
\item[splash\textunderscore screen]
Si \textit{NetHack} le peut, il affiche un écran d'accueil à son
lancement (activée par défaut).
%.lp
\item[statuslines]
Nombre de lignes de l'affichage d'état traditionnel, sous la carte.
Les valeurs acceptées sont \texttt{2} et \texttt{3} (\texttt{2} par défaut).

%.lp ""
Avec la valeur \texttt{3}, l'interface \texttt{tty} déplace certains champs et
affiche surtout les états (conditions) sur une ligne à part.
Un écran capable d'afficher au moins 25 lignes est recommandé.
La valeur peut être modifiée dans un sens ou dans l'autre en cours de partie avec la commande `\texttt{O}'.

%.lp ""
L'interface \texttt{curses} fait de même si l'option
\textit{align\textunderscore status}
vaut \textit{top} ou \textit{bottom}, mais ignore
\textit{statuslines}
si elle vaut \textit{left} ou \textit{right}.

%.lp ""
L'interface \texttt{Qt} affiche déjà plus de 3 lignes d'état
et utilise donc la valeur de
\textit{statuslines}
différemment.
La valeur \texttt{3} affiche l'état dans le format d'origine de
l'interface \texttt{Qt}, avec une fenêtre d'état étalée verticalement.
La valeur \texttt{2} condense légèrement l'état, en déplaçant certains
champs sur d'autres lignes pour supprimer une ligne entière, ce qui réduit
la hauteur nécessaire.
(Si NetHack a été compilé avec une version de \texttt{Qt}
antérieure à \texttt{qt-5.9},
\textit{statuslines}
ne peut être réglée que dans le fichier de configuration ou via NETHACKOPTIONS,
et non en cours de partie avec la commande `\texttt{O}'.)
%.lp
\item[term\textunderscore cols \textrm{et}]
%.lp
\item[term\textunderscore rows]
Interface curses uniquement.
Nombre de colonnes et de lignes à utiliser pour l'affichage.
Curses essaiera de se redimensionner aux valeurs indiquées, mais se contentera
de tailles plus petites si elles sont trop grandes.
Par défaut, la taille actuelle de la fenêtre.
%.lp
\item[tile\textunderscore file]
Indique le nom d'un autre fichier de tuiles à utiliser à la place de celui par défaut.
\\
%.lp ""
Remarque : l'interface X11 utilise les ressources X plutôt que les options de
NetHack pour choisir un autre fichier de tuiles.
Voyez \texttt{NetHack.ad}, le fichier d'exemple de ``valeurs par défaut de l'application'' X.
%.lp
\item[tile\textunderscore height]
Indique la hauteur souhaitée de chaque tuile, pour une interface gérant les tuiles.
%.lp
\item[tile\textunderscore width]
Indique la largeur souhaitée de chaque tuile, pour une interface gérant les tuiles.
%.lp
\item[tiled\textunderscore map]
Si \textit{NetHack} le peut, il affiche la carte avec des graphismes en \textit{tuiles}
plutôt qu'avec de simples caractères (lettres et ponctuation, éventuellement
complétées par des symboles de tracé de lignes).
Régler \texttt{tiled\textunderscore map} à \textit{True} force
\texttt{ascii\textunderscore map} à \textit{False}.
%.lp
\item[use\textunderscore darkgray]
Utilise du noir en gras plutôt que du bleu pour les glyphes noirs (TTY uniquement).
%.lp
\item[use\textunderscore inverse]
Si \textit{NetHack} le peut, il affiche en vidéo inverse lorsque le jeu le demande.
%.lp
\item[use\textunderscore menu\textunderscore glyphs]
Si \textit{NetHack} le peut, il affiche les glyphes à côté des objets dans
l'inventaire.
%.lp
\item[vary\textunderscore msgcount]
Si \textit{NetHack} le peut, il affiche ce nombre de messages à la fois
dans la fenêtre des messages.
%.lp
\item[windowborders]
Indique s'il faut tracer des cadres autour de la carte, de la zone d'état, de la
zone des messages et de la fenêtre d'inventaire permanent si elle est activée.
Interface curses uniquement.
Les valeurs acceptées sont

%.sd
%.si
\texttt{0} --- désactivé, jamais de bordures\\
\texttt{1} --- activé, toujours des bordures\\
\texttt{2} --- auto, activé si l'écran fait au moins
$(24+2)\times(80+2)$ [défaut]\\
\texttt{3} --- activé, sauf désactivation forcée pour perm\textunderscore invent\\
\texttt{4} --- auto, sauf désactivation forcée pour perm\textunderscore invent\\
%.ei
%.ed

%.lp ""
(Le seuil de taille 26x82 de la valeur `2' correspond au nombre de lignes et
de colonnes de l'affichage.
Une largeur d'au moins 110 colonnes ($80+2+26+2$) est nécessaire si
\textit{align\textunderscore status}
vaut \texttt{left} ou \texttt{right}.)

%.lp ""
La fenêtre d'inventaire permanent, lorsqu'elle est activée, peut grandir jusqu'à
devenir trop grande pour la plupart des écrans, ce qui tronque son contenu.
Si les bordures sont forcées (1) ou si l'écran est assez grand pour les afficher (2),
choisir plutôt la valeur 3 ou 4 conserve les bordures des fenêtres de la carte,
des messages et d'état, tout en laissant la place à deux lignes d'inventaire
supplémentaires et en élargissant chaque ligne d'inventaire de deux colonnes.
%.lp
\item[windowcolors]
Si \textit{NetHack} le peut, il affiche toutes les fenêtres d'un style donné
avec les couleurs de premier plan et d'arrière-plan indiquées.
Interface graphique Windows et interface curses uniquement.
Le format est\\
\texttt{~~~~OPTION=windowcolors:}\textit{style premier\textunderscore plan}\texttt{/}\textit{arrière\textunderscore plan}\\
où \textit{style} est \texttt{menu}, \texttt{message}, \texttt{status}
ou \texttt{text}, et où
\textit{premier\textunderscore plan} et \textit{arrière\textunderscore plan} sont des couleurs, soit numériques
(un dièse suivi de trois paires de chiffres hexadécimaux, \textit{\#rrggbb}),
soit l'une des couleurs nommées (\textit{black}, \textit{red}, \textit{green}, \textit{brown},
\textit{blue}, \textit{magenta}, \textit{cyan}, \textit{orange},
\textit{bright-green}, \textit{yellow}, \textit{bright-blue}, \textit{bright-magenta},
\textit{bright-cyan}, \textit{white}, \textit{gray}, \textit{purple},
\textit{silver}, \textit{maroon}, \textit{fuchsia}, \textit{lime}, \textit{olive},
\textit{navy}, \textit{teal}, \textit{aqua}),
soit (sous Windows uniquement) l'une des couleurs de l'interface Windows (\textit{trueblack},
\textit{activeborder}, \textit{activecaption}, \textit{appworkspace}, \textit{background},
\textit{btnface}, \textit{btnshadow}, \textit{btntext}, \textit{captiontext},
\textit{graytext}, \textit{greytext}, \textit{highlight},
\textit{highlighttext}, \textit{inactiveborder}, \textit{inactivecaption}, \textit{menu},
\textit{menutext}, \textit{scrollbar}, \textit{window}, \textit{windowframe},
\textit{windowtext}).

%.lp
\item[wraptext]
Si \textit{NetHack} le peut, il coupe les longues lignes de texte qui ne tiennent pas
dans la zone visible de la fenêtre.
\end{description}

%.hn 2
\subsection*{Options de rapport de plantage}
%.pg

Veuillez noter que NetHack n'envoie {\textbf aucune} information hors de votre
ordinateur, sauf si vous cliquez vous-même sur le bouton d'envoi d'un formulaire.
%.si
\begin{description}
%.lp
\item[OPTION=crash\textunderscore email:\textit{adresse\textunderscore email}]
%.lp
\item[OPTION=crash\textunderscore name:\textit{votre\textunderscore nom}]
%.ei
\end{description}
Ces options servent uniquement à vous épargner un peu de saisie dans les
formulaires de rapport de plantage et de \#bugreport.
%.si
\begin{description}
%.lp
\item[OPTION=crash\textunderscore urlmax:\textit{octets}]
%.ei
\end{description}
Cette option sert à limiter la longueur des URL générées ; elle n'est
utile que si votre navigateur ne gère pas les URL de longueur arbitraire.

%.hn 2
\subsection*{Options de personnalisation propres à une plateforme}

%.pg
Voici l'explication des options utilisées par certaines plateformes
ou versions pour personnaliser et modifier leur comportement.

\begin{description}[font=\mdseries\itshape]
%.lp
\item[altkeyhandling]
Choisit une autre manière de gérer les frappes au clavier (\textit{NetHack tty Win32} uniquement).
Le nom du type de gestion est \textit{default}, \textit{ray} ou \textit{340}
%.\" \item[altmeta]
%.\" On Amiga, this option controls whether typing ``Alt'' plus another key
%.\" functions as a meta-shift for that key (default on).
%.lp
\item[altmeta]
%.\" On other (non-Amiga) systems where this option is available, it can be
Sur les systèmes où cette option est disponible, elle peut être
activée pour que \textit{NetHack} convertisse une séquence de deux caractères commençant
par ESC en la version ``meta'' du second caractère (désactivée par défaut).

%.lp ""
Cette conversion n'est faite que pour les commandes, pas pour les autres saisies.
Notez que la saisie d'un ou plusieurs chiffres comme préfixe de nombre avant une
commande---précédés de \texttt{n} si l'option \textit{number\textunderscore pad}
est activée---est
également soumise à cette conversion : si vous tentez
d'annuler le nombre en tapant ESC, \textit{NetHack} attendra un autre
caractère pour compléter la séquence de deux caractères.
Tapez un second ESC pour achever l'annulation du nombre.
Pour les autres questions, un seul ESC suffit.
%.lp
\item[BIOS]
Utilise les appels BIOS pour mettre à jour rapidement l'affichage et pour lire le clavier
(ce qui permet d'utiliser les touches fléchées pour se déplacer) sur les machines dotées
d'une ROM BIOS compatible IBM PC (désactivée par défaut ; \textit{NetHack OS/2, PC {\rm et} ST} uniquement).
%.lp
\item[rawio]
Force le mode brut (non cbreak) pour un affichage plus rapide et une saisie
plus robuste (sans cela, MS-DOS traite parfois `\texttt{\textasciicircum P}' comme une bascule
de l'imprimante) (désactivée par défaut ; \textit{NetHack OS/2, PC {\rm et} ST} uniquement).
Remarque : les DEC Rainbow se bloquent si elle est activée.
Ne peut pas être modifiée avec la commande `\texttt{O}'.
%.lp
\item[subkeyvalue]
(\textit{NetHack tty Win32} uniquement).
Permet de modifier la valeur des frappes que le système d'exploitation
renvoie à \textit{NetHack}, pour compenser les problèmes de claviers
internationaux.
OPTIONS=subkeyvalue:171/92
renverra 92 à \textit{NetHack} si 171 devait être renvoyé à l'origine.
Vous pouvez utiliser plusieurs affectations subkeyvalue dans le fichier de
configuration si nécessaire.
Ne peut pas être modifiée avec la commande `\texttt{O}'.
%.lp
\item[video]
Règle le mode vidéo utilisé (\textit{NetHack PC} uniquement).
Les valeurs sont \textit{autodetect}, \textit{default}, \textit{vga} ou \textit{vesa}.
Avec \textit{vesa}, le jeu affiche des tuiles en exploitant toutes les
capacités du matériel VGA.
Avec \textit{vga}, le jeu affiche des tuiles en 640x480 fixe
et 16 couleurs, un mode compatible avec tout le matériel VGA. Les jeux
de tuiles tiers ne fonctionneront probablement pas.
Avec \textit{autodetect}, le jeu essaie \textit{vesa}, puis \textit{vga}, et
choisit finalement \textit{default} si aucun de ces modes ne fonctionne.
Ne peut pas être modifiée avec la commande `\texttt{O}'.
%.lp
\item[video\textunderscore height]
Règle la hauteur de la résolution du mode VGA (MS-DOS uniquement, avec video:vesa)
%.lp
\item[video\textunderscore width]
Règle la largeur de la résolution du mode VGA (MS-DOS uniquement, avec video:vesa)
%.lp
\item[videocolors]
\begin{sloppypar}
Règle la palette de couleurs des systèmes PC utilisant NO\textunderscore TERMS
(par défaut 4-2-6-1-5-3-15-12-10-14-9-13-11 ; \textit{NetHack PC} uniquement).
L'ordre des couleurs est : rouge, vert, brun, bleu, magenta, cyan,
blanc vif, rouge vif, vert vif, jaune, bleu vif,
magenta vif et cyan vif.
Ne peut pas être modifiée avec la commande `\texttt{O}'.
\end{sloppypar}
%.lp
\item[videoshades]
Règle le niveau d'intensité des trois niveaux de gris disponibles
(par défaut dark normal light ; \textit{NetHack PC} uniquement).
Si l'affichage du jeu est difficile à lire, essayez d'ajuster ces niveaux ;
si cela ne corrige pas le problème, essayez \texttt{!color}.
Ne peut pas être modifiée avec la commande `\texttt{O}'.
\end{description}

%.hn 2
\subsection*{Expressions rationnelles}

%.pg
Les expressions rationnelles sont normalement des expressions rationnelles
étendues POSIX. Il est possible de compiler \textit{NetHack} sans prise en charge des
expressions rationnelles sur
une plateforme dépourvue de bibliothèque adéquate.
Bien que ce ne soit le cas d'aucune plateforme moderne, si votre \textit{NetHack} a été
compilé ainsi, les motifs sont alors des motifs de type ``glob'' ; quoi qu'il en
soit, ce document appelle les uns comme les autres ``expressions rationnelles''.
Cela concerne les exceptions de ramassage automatique, les types de messages,
les couleurs des menus et les sons de l'utilisateur.

%.hn 2
\subsection*{Configurer les exceptions de ramassage automatique}

%.pg
Vous pouvez affiner le comportement de l'option ``\texttt{autopickup}''
au-delà de ce que permet l'option ``\texttt{pickup\textunderscore types}''.

%.pg
En plaçant des lignes ``\texttt{autopickup\textunderscore exception}'' dans votre fichier de
configuration, vous pouvez définir des motifs à vérifier lorsque le jeu
s'apprête à ramasser automatiquement quelque chose.

\begin{description}
%.lp
\item[autopickup\textunderscore exception]
Définit une exception à l'option ``\textit{pickup\textunderscore types}''.
L'option \textit{autopickup\textunderscore exception} doit être suivie d'une expression
rationnelle qui sert de motif, comparé à la forme au singulier de la
description d'un objet se trouvant sur votre case.

De plus, certains caractères ont un sens particulier lorsqu'ils sont le premier
caractère du motif, à savoir :

%.sd
%.si
\texttt{<} --- toujours ramasser un objet correspondant au reste du motif ;\\
\texttt{>} --- ne jamais ramasser un objet correspondant au reste du motif.
%.ei
%.ed

Les règles \textit{autopickup\textunderscore exception} sont traitées dans l'ordre
où elles apparaissent dans votre fichier de configuration, ce qui permet à une
règle ultérieure d'outrepasser une règle antérieure.

%.lp ""
Les exceptions peuvent être définies avec la commande `\texttt{O}', mais comme elles ne
figurent pas dans votre fichier de configuration, elles ne seront plus en vigueur
si vous sauvegardez puis restaurez votre partie.
Les règles \textit{autopickup\textunderscore exception} ne sont pas sauvegardées avec la partie.
\end{description}

%.lp "Here are some examples:"
Voici quelques exemples :
\begin{verbatim}
    autopickup_exception="<*flèche*"
    autopickup_exception=">*cadavre*"
    autopickup_exception=">* maudit*"
\end{verbatim}

%.pg
Le premier exemple ci-dessus provoque le ramassage automatique de tout type de flèche.
Le deuxième exclut tout cadavre du ramassage automatique.
Le dernier exclut du ramassage automatique les objets que l'on sait maudits
(attention : il reconnaît aussi ``non maudit'').

%.lp

%.hn 2
\subsection*{Modifier les associations de touches}

%.pg
Il est possible de modifier les touches associées par défaut à certaines commandes
spéciales, aux raccourcis des menus et aux commandes étendues, à l'aide de
directives BIND dans le fichier de configuration. Le format est la touche, suivie
de la commande à lui associer, séparées par un deux-points. La touche peut être un caractère seul (``\texttt{x}''),
une touche Contrôle (``\texttt{\textasciicircum X}'', ``\texttt{C-x}''), une touche Meta (``\texttt{M-x}''),
un bouton de la souris, ou un code ASCII décimal à trois chiffres.

%.pg
Par exemple :

\begin{verbatim}
    BIND=^X:getpos.autodescribe
    BIND=\:menu_first_page
    BIND=v:loot
\end{verbatim}

\begin{description}[font=\mdseries\ttfamily]
%.lp "Extended command keys"
\item[Touches des commandes étendues]
Vous pouvez associer plusieurs touches à la même commande étendue. Pour libérer
une touche, associez-la à la commande étendue ``\texttt{nothing}''. Vous pouvez aussi associer
les touches ``\texttt{<esc>}'', ``\texttt{<enter>}'' et ``\texttt{<space>}''.

%.lp "Menu accelerator keys"
\item[Raccourcis des menus]
Les touches de contrôle ou raccourcis des menus peuvent aussi être réassociés
via des lignes OPTIONS dans le fichier de configuration.
Vous ne pouvez pas associer des symboles d'objets ou des lettres de sélection
à des raccourcis de menu.
Certaines interfaces ne gèrent qu'une partie des raccourcis de menu.

%.lp "Mouse buttons"
\item[Boutons de la souris]
Vous pouvez associer ``mouse1'' ou ``mouse2'' à ``\texttt{nothing}'',
``\texttt{therecmdmenu}'', ``\texttt{clicklook}'' ou ``\texttt{mouseaction}''.

%.lp "Special command keys"
\item[Touches des commandes spéciales]
Voici les commandes spéciales que vous pouvez réassocier. Certaines peuvent être associées aux
mêmes touches sans problème ; d'autres partagent le même ``contexte'' et, si elles sont
associées aux mêmes touches, une seule d'entre elles sera disponible. Une commande spéciale
ne peut être associée qu'à une seule touche.
\end{description}

%.pg
\begin{description}[itemindent=10mm, labelwidth=15mm, rightmargin=15mm]
%.lp
\item[count]
Touche préfixe pour commencer un nombre, afin de répéter une commande autant de fois.
Avec \textit{number\textunderscore pad} uniquement. Par défaut,~`\texttt{n}'.
%.lp
\item[getdir.help]
Lorsqu'une direction est demandée, touche pour afficher l'aide. Par défaut,~`\texttt{?}'.
%.lp
\item[getdir.mouse]
Lorsqu'une direction est demandée, touche pour simuler un clic de souris.
Il vous sera demandé de choisir un emplacement.
Utilisez les touches de déplacement pour déplacer le curseur sur la carte, puis tapez
la touche getpos.pick.once (`\texttt{,}' par défaut)
ou la touche getpos.pick (`\texttt{.}' par défaut)
pour terminer comme par un clic gauche ou droit.
N'est utile qu'avec la commande \texttt{\#therecmdmenu}.
Par défaut,~`\texttt{\textunderscore}'.
%.lp
\item[getdir.self]
Lorsqu'une direction est demandée, touche pour vous viser vous-même. Par défaut,~`\texttt{.}'.
%.lp
\item[getdir.self2]
Lorsqu'une direction est demandée, autre touche pour vous viser vous-même.
Par défaut,~`\texttt{s}'.
%.lp
\item[getpos.autodescribe]
Lorsqu'un emplacement est demandé, touche pour activer ou désactiver \textit{autodescribe}.
Par défaut,~`\texttt{\#}'.
%.lp
\item[getpos.all.next]
Lorsqu'un emplacement est demandé, touche pour aller à la chose intéressante suivante la plus proche.
Par défaut,~`\texttt{a}'.
%.lp
\item[getpos.all.prev]
Lorsqu'un emplacement est demandé, touche pour aller à la chose intéressante précédente la plus proche.
Par défaut,~`\texttt{A}'.
%.lp
\item[getpos.door.next]
Lorsqu'un emplacement est demandé, touche pour aller à la porte ou embrasure suivante la plus proche.
Par défaut,~`\texttt{d}'.
%.lp
\item[getpos.door.prev]
Lorsqu'un emplacement est demandé, touche pour aller à la porte ou embrasure précédente la plus proche.
Par défaut,~`\texttt{D}'.
%.lp
\item[getpos.help]
Lorsqu'un emplacement est demandé, touche pour afficher l'aide. Par défaut,~`\texttt{?}'.
%.lp
\item[getpos.mon.next]
Lorsqu'un emplacement est demandé, touche pour aller au monstre suivant le plus proche.
Par défaut,~`\texttt{m}'.
%.lp
\item[getpos.mon.prev]
Lorsqu'un emplacement est demandé, touche pour aller au monstre précédent le plus proche.
Par défaut,~`\texttt{M}'.
%.lp
\item[getpos.obj.next]
Lorsqu'un emplacement est demandé, touche pour aller à l'objet suivant le plus proche.
Par défaut,~`\texttt{o}'.
%.lp
\item[getpos.obj.prev]
Lorsqu'un emplacement est demandé, touche pour aller à l'objet précédent le plus proche.
Par défaut,~`\texttt{O}'.
%.lp
\item[getpos.menu]
Lorsqu'un emplacement est demandé et que vous utilisez l'une des touches suivant ou
précédent pour parcourir les cibles, affiche plutôt un menu. Par défaut,~`\texttt{!}'.
%.lp
\item[getpos.moveskip]
Lorsqu'un emplacement est demandé et que vous utilisez les touches de déplacement
avec Maj ou les touches meta-chiffre pour vous déplacer rapidement, saute les
glyphes identiques au lieu d'avancer de 8 cases.
Par défaut,~`\texttt{*}'.
%.lp
\item[getpos.filter]
Lorsqu'un emplacement est demandé, change le mode de filtrage utilisé par les
touches suivant ou précédent pour parcourir les cibles. Alterne entre aucun
filtrage, uniquement ce qui est en vue, et uniquement la même zone. Par défaut,~`\texttt{"}'.
%.lp
\item[getpos.pick]
Lorsqu'un emplacement est demandé, touche pour choisir l'emplacement, et éventuellement
demander plus d'informations.
Lors de la simulation d'un clic de souris après une demande de direction (voir
getdir.mouse ci-dessus), touche pour répondre comme par un clic droit.
Par défaut,~`\texttt{.}'.
%.lp
\item[getpos.pick.once]
Lorsqu'un emplacement est demandé, touche pour choisir l'emplacement sans
demander plus d'informations.
Lors de la simulation d'un clic de souris après une demande de direction,
touche pour répondre comme par un clic gauche.
Par défaut,~`\texttt{,}'.
%.lp
\item[getpos.pick.quick]
Lorsqu'un emplacement est demandé, touche pour choisir l'emplacement, sans demander
plus d'informations, et quitter la boucle de demande d'emplacement. Par défaut,~`\texttt{;}'.
%.lp
\item[getpos.pick.verbose]
Lorsqu'un emplacement est demandé, touche pour choisir l'emplacement et afficher plus
d'informations sans le demander. Par défaut,~`\texttt{:}'.
%.lp
\item[getpos.self]
Lorsqu'un emplacement est demandé, touche pour aller à votre emplacement.
Par défaut,~`\texttt{@}'.
%.lp
\item[getpos.unexplored.next]
Lorsqu'un emplacement est demandé, touche pour aller à l'emplacement inexploré suivant le plus proche.
Par défaut,~`\texttt{x}'.
%.lp
\item[getpos.unexplored.prev]
Lorsqu'un emplacement est demandé, touche pour aller à l'emplacement inexploré
précédent le plus proche. Par défaut,~`\texttt{X}'.
%.lp
\item[getpos.valid]
Lorsqu'un emplacement est demandé, touche pour afficher les emplacements cibles valides.
Par défaut,~`\texttt{\$}'.
%.lp
\item[getpos.valid.next]
Lorsqu'un emplacement est demandé, touche pour aller à l'emplacement valide suivant le plus proche.
Par défaut,~`\texttt{z}'.
%.lp
\item[getpos.valid.prev]
Lorsqu'un emplacement est demandé, touche pour aller à l'emplacement valide précédent le plus proche.
Par défaut,~`\texttt{Z}'.
\end{description}


%.hn 2
\subsection*{Configurer les types de messages}

%.pg
Vous pouvez modifier la façon dont les messages sont affichés dans la zone des
messages lorsqu'ils correspondent à un motif défini par l'utilisateur.

%.pg
En général, les entrées du fichier de configuration décrivant les types de
messages ressemblent à ceci :
\begin{verbatim}
    MSGTYPE=type "motif"
\end{verbatim}
\begin{description}[font=\mdseries\itshape]
%.lp
\item[type]
la façon dont le message doit être affiché :
%.sd
%.si
\\
\texttt{show}  --- afficher le message normalement.\\
\texttt{hide}  --- ne jamais afficher le message.\\
\texttt{stop}  --- attendre l'utilisateur avec l'invite \texttt{-{}-Plus-{}-}.\\
\texttt{norep} --- afficher le message une fois, mais pas de nouveau si aucun autre message n'a été
affiché entre-temps.
%.ei
%.ed
%.lp
\item[motif]
le motif à reconnaître. Le motif doit être une expression rationnelle.
\end{description}

%.lp ""
Voici un exemple de types de messages utilisant le mécanisme interne de
reconnaissance de motifs de \textit{NetHack} :

\begin{verbatim}
    MSGTYPE=stop "Vous avez faim."
    MSGTYPE=hide "Vous échangez votre place avec *."
\end{verbatim}

indique que, chaque fois que le message ``Vous avez faim.'' est affiché,
l'utilisateur doit valider l'invite \texttt{-{}-Plus-{}-}, et qu'un message correspondant à
``Vous échangez votre place avec <quelque chose>.'' n'est pas affiché du tout.

%.lp
L'ordre des lignes MSGTYPE définies est important : c'est la dernière règle
correspondante qui est utilisée. Placez d'abord le cas général, puis les exceptions.

%.pg

%.lp
%.hn 2
\subsection*{Configurer les couleurs des menus}

%.pg
Certaines plateformes vous permettent de définir les couleurs des lignes de menu
lorsque la ligne correspond à un motif défini par l'utilisateur.
Actuellement, les interfaces tty, curses, win32tty et
win32gui le permettent.

%.pg
En général, les entrées du fichier de configuration décrivant les couleurs des
menus ressemblent à ceci :
\begin{verbatim}
    MENUCOLOR="motif"=couleur&attribut
\end{verbatim}

\begin{description}[font=\mdseries\itshape]
%.lp
\item[motif]
le motif à reconnaître ;
%.lp
\item[couleur]
la couleur à utiliser pour les lignes correspondant au motif ;
%.lp
\item[attribut]
l'attribut à utiliser pour les lignes correspondant au motif. L'attribut est
facultatif ; si vous l'omettez, vous devez aussi omettre l'esperluette qui le précède.
Si aucun attribut n'est défini, aucun attribut n'est utilisé.
\end{description}

%.lp ""
Le motif doit être une expression rationnelle.

%.lp ""
Les couleurs autorisées sont \textit{black}, \textit{red}, \textit{green}, \textit{brown},
\textit{blue}, \textit{magenta}, \textit{cyan}, \textit{gray}, \textit{orange},
\textit{light-green}, \textit{yellow}, \textit{light-blue}, \textit{light-magenta},
\textit{light-cyan} et \textit{white}.
Plus \textit{no-color}, la couleur de premier plan par défaut, qui n'est pas
forcément identique à l'une des autres couleurs.

%.lp ""
Les attributs autorisés sont \textit{none}, \textit{bold}, \textit{dim}, \textit{italic},
\textit{underline}, \textit{blink} et \textit{inverse}.
\textit{Normal} est un synonyme de \textit{none}.
Notez que la plateforme utilisée peut interpréter les attributs comme bon
lui semble.

%.lp ""
Voici un exemple de couleurs de menus utilisant le mécanisme interne de
reconnaissance de motifs de \textit{NetHack} :

\begin{verbatim}
    MENUCOLOR="* béni*"=green
    MENUCOLOR="* maudit*"=red
    MENUCOLOR="* maudit*(porté*)"=red&underline
\end{verbatim}

indique que toute ligne de menu contenant ``~béni''
sera affichée en vert, que les lignes contenant ``~maudit'' seront
affichées en rouge, et que les lignes contenant ``~maudit'' suivi de ``(porté)''
sur la même ligne seront affichées en rouge et soulignées.
Vous pouvez avoir plusieurs entrées MENUCOLOR dans votre fichier de configuration ;
c'est la dernière ligne MENUCOLOR correspondant
à une ligne de menu qui est utilisée pour celle-ci.

%.pg
Notez que si vous voulez qu'une ou plusieurs spécifications de couleur reconnaissent
``~non~maudit'' (que le motif ``~maudit'' ci-dessus reconnaît d'ailleurs
aussi), vous voudrez sans doute désactiver l'option
\textit{implicit\textunderscore uncursed}
afin que tous les objets que l'on sait non maudits soient effectivement
affichés avec la mention ``non maudit''.

%.lp
%.hn 2
\subsection*{Configurer les sons de l'utilisateur}

%.pg
Certaines plateformes vous permettent de définir des fichiers son à jouer lorsqu'un
message correspondant à un motif défini par l'utilisateur est envoyé à la fenêtre
des messages.
Actuellement, l'interface Qt et les interfaces win32tty et win32gui gèrent
les sons de l'utilisateur.

%.pg
Les entrées suivantes du fichier de configuration servent à associer des sons
utilisateur aux messages :

\begin{description}[font=\mdseries\itshape]
%.lp
\item[SOUNDDIR]
Le répertoire qui contient les fichiers son à jouer.
%.lp
\item[SOUND]
Une entrée qui associe un fichier son à un motif de message défini par l'utilisateur.
Chaque entrée SOUND se décompose ainsi :

%.sd
%.si
\texttt{MESG} --- association à la fenêtre des messages (la seule prise en charge depuis la 5.0.0) ;\\
\texttt{msgtype} --- facultatif ; type de message à utiliser, voir ``Configurer les types de messages''\\
\texttt{motif} --- le motif à reconnaître ;\\
\texttt{fichier son} --- le fichier son à jouer ;\\
\texttt{volume} --- le volume à utiliser pour jouer le fichier son ;\\
\texttt{index du son} --- facultatif ; l'index correspondant à un fichier son.
%.ei
%.ed
\end{description}

%.lp ""
Le motif doit être une expression rationnelle.

Par exemple :

\begin{verbatim}
    SOUNDDIR=C:\nethack\sounds
    SOUND=MESG "Cette porte est verrouillée" "lock.wav" 100
    SOUND=MESG hide "^Vous ratez " "swing.wav" 75
\end{verbatim}
%.pg

%.lp
%.hn 2
\subsection*{Configurer la mise en valeur de l'état}

%.pg
Votre copie de \textit{NetHack} a peut-être été compilée avec la prise en charge
de la \textit{mise en valeur de l'état} (\textit{Status Hilites}).
Si c'est le cas, vous pouvez personnaliser l'affichage du jeu en définissant des
seuils qui changent la couleur ou l'apparence des champs de l'affichage d'état.

Le format de définition des couleurs d'état est :\\
\begin{verbatim}
OPTION=hilite_status:nom-du-champ/comportement/couleur&attributs
\end{verbatim}

Par exemple, la ligne suivante dans votre fichier de configuration affichera
le champ des points de vie en rouge si vos points de vie
tombent à 30~\% ou moins :\\
\begin{verbatim}
OPTION=hilite_status:hitpoints/<=30%/red/normal
\end{verbatim}
(Cet exemple spécifie en fait \texttt{red\& normal} pour  \texttt{<=30\%}
et \texttt{no-color\& normal} pour \texttt{>30\%}.)\\

Autre exemple : la ligne suivante dans votre fichier de configuration affichera
la sagesse en rouge si elle baisse et en vert si elle augmente :\\
\begin{verbatim}
OPTION=hilite_status:wisdom/down/red/up/green
\end{verbatim}

Les couleurs autorisées sont black, red, green, brown, blue, magenta, cyan, gray,
orange, light-green, yellow, light-blue, light-magenta, light-cyan et white.
Plus \textit{no-color}, la couleur de premier plan par défaut de l'affichage, qui
n'est pas forcément identique au noir, au blanc ou à l'une des autres couleurs.

Les attributs autorisés sont none, bold, dim, underline, italic, blink et inverse.
``Normal'' est un synonyme de ``none'' ; ils ne doivent pas être combinés
avec l'un des autres attributs.

Pour indiquer à la fois une couleur et un attribut, combinez-les avec `\&'.
Pour indiquer plusieurs attributs, combinez-les avec `+'.

%.lp ""
Par exemple : \texttt{magenta\& inverse+dim}.

Notez que l'affichage peut remplacer ou ignorer certains attributs
selon ses capacités et, de manière générale, interpréter les
attributs comme bon lui semble.
Par exemple, sur certains systèmes d'affichage, demander du gras peut donner
du clignotement, ou inversement.
Sur d'autres, demander un attribut alors qu'un autre est déjà
en place remplace l'attribut précédent au lieu de s'y ajouter.
Comme NetHack envoie les demandes d'attributs les unes après les autres (du moins
avec l'interface \textit{tty}) plutôt que toutes à la fois, la seule façon de
maîtriser une telle situation est de n'indiquer qu'un seul attribut.

Vous pouvez modifier l'affichage des champs d'état suivants :
%.sd
\begin{center}
\begin{tabular}{lll}
%TABLE_START
title & dungeon-level & experience-level\\
strength & gold & experience\\
dexterity & hitpoints & HD\\
constitution & hitpoints-max & time\\
intelligence & power & hunger\\
wisdom & power-max & carrying-capacity\\
charisma & armor-class & condition\\
alignment &  & score\\
%TABLE_END  Do not delete this line.
\end{tabular}
\end{center}
%.ed
%.lp ""
Le pseudo-champ `characteristics' permet de régler les six
caractéristiques Fo, Dx, Co, In, Sa et Ch d'un coup.  `HD' désigne les
`dés de vie' (\textit{hit dice}), une approximation du niveau d'expérience affichée
lorsque vous êtes métamorphosé.
`experience', `time' et `score' ne sont affichés que
selon vos autres réglages d'options.

%.lp ""
Au lieu d'un comportement, `condition' accepte les indicateurs d'état suivants :
\textit{stone}, \textit{slime}, \textit{strngl}, \textit{foodpois}, \textit{termill},
\textit{blind}, \textit{deaf}, \textit{stun}, \textit{conf}, \textit{hallu},
\textit{lev}, \textit{fly} et \textit{ride}.
Vous pouvez utiliser `major\textunderscore troubles' comme alias
de stone à termill, `minor\textunderscore troubles' de blind à hallu,
`movement' de lev, fly et ride, et `all' de tous les états.

%.lp ""
Les comportements autorisés sont ``always'', ``up'', ``down'', ``changed'', un
seuil en pourcentage ou en valeur absolue, ou un texte à reconnaître.
Pour le champ \textit{hitpoints}, le comportement supplémentaire ``criticalhp''
est disponible.
Il l'emporte sur les autres règles de comportement si les
points de vie sont au niveau du seuil de \textit{problème majeur} ou en dessous
(seuil qui varie selon les points de vie maximum et le niveau d'expérience).

\begin{description}[font=\mdseries\ttfamily]
%.lp "*"
\item[always] définit les attributs par défaut de ce champ.
%.lp "*"
\item[up \textrm{\textmd{et}} down] définissent les attributs du champ
lorsque sa valeur augmente ou diminue. Cet attribut
expire au bout de \texttt{statushilites} tours.
%.lp "*"
\item[changed] définit l'attribut du champ lorsque sa valeur
change. Cet attribut expire au bout de \texttt{statushilites} tours.
(Si un champ possède à la fois une règle ``changed'' et une règle ``up''
ou ``down'' qui correspond à un changement de sa valeur,
c'est la règle ``up'' ou ``down'' qui l'emporte.)
%.lp "*"
\item[pourcentage] définit l'attribut du champ lorsque sa valeur
correspond au pourcentage.
Il s'écrit sous la forme d'un nombre entre 0 et 100, suivi de `\texttt{\%}'
(signe pour cent).
Si le pourcentage est précédé de `\texttt{<=}' ou `\texttt{>=}',
il correspond aussi lorsque la valeur est inférieure ou supérieure au pourcentage.
Utilisez le préfixe `\texttt{<}' ou `\texttt{>}' pour une valeur strictement inférieure ou supérieure.
(La limite numérique est un peu assouplie dans ce cas : \texttt{>-1\%}
et \texttt{<101\%} sont autorisés.)
Seuls quatre champs acceptent les règles en pourcentage.
Les pourcentages de ``\textit{hitpoints}'' et ``\textit{power}'' sont
simples : ils se basent sur le champ maximum correspondant.
Les règles en pourcentage sont aussi autorisées pour le ``\textit{niveau d'expérience}''
et les ``\textit{points d'expérience}'' (valables lorsque l'option
\textit{showexp}
est activée).
Pour ceux-ci, le pourcentage se base sur la progression entre le début du
niveau d'expérience actuel et le début du niveau suivant.
Ainsi, si le niveau 2 commence à 20 points et le niveau 3 à 40 points,
avoir 30 points correspond à 50~\% et 35 points à 75~\%.
100~\% est inaccessible pour l'expérience, car vous gagnerez un niveau et
les calculs repartiront de zéro pour ce nouveau niveau, mais une règle
\texttt{=100\%} est autorisée et correspond au cas particulier où il vous manque
exactement 1 point d'expérience pour atteindre le niveau suivant.
% (If you manage to reach level 30, there is no next level and the
% percentage will remain at 0\% no matter have many additional experience
% points you earn.)
%.lp "*"
\item[valeur] absolue : définit l'attribut lorsque la valeur du champ
correspond à ce nombre.
Le nombre doit être supérieur ou égal à 0, sauf pour ``\textit{armor-class}'' qui
accepte les valeurs négatives, et peut éventuellement être précédé de `\texttt{=}'.
Si le nombre est précédé de `\texttt{<=}' ou `\texttt{>=}' à la place,
il correspond aussi lorsque la valeur est inférieure ou supérieure.
Si le préfixe est `\texttt{<}' ou `\texttt{>}', il ne correspond que pour une valeur
strictement supérieure ou inférieure.
%.lp "*"
\item[criticalhp] ne s'applique qu'au champ hitpoints, et seulement
lorsque les points de vie actuels sont sous un seuil (qui varie selon les
points de vie maximum et le niveau d'expérience).
Lorsque le seuil est atteint, une règle criticalhp l'emporte sur toutes les
autres règles de hitpoints.
%.lp "*"
\item[texte] : définit l'attribut lorsque la valeur du champ
correspond à ce texte.
La reconnaissance de texte n'est possible que pour ``\textit{alignment}'',
``\textit{carrying-capacity}'', ``\textit{hunger}'', ``\textit{dungeon-level}''
et ``\textit{title}''.
Pour title, seul le titre de rang du rôle
est testé ; le nom du personnage est ignoré.
%.ei
\end{description}

Le menu des options du jeu peut vous aider à trouver la syntaxe correcte pour un
fichier de configuration.

Toute la fonctionnalité peut être désactivée en réglant l'option \textit{statushilites} à 0.

Exemples de mises en valeur :
\begin{verbatim}
    OPTION=hilite_status: gold/up/yellow/down/brown
    OPTION=hilite_status: characteristics/up/green/down/red
    OPTION=hilite_status: hitpoints/100%/gray&normal
    OPTION=hilite_status: hitpoints/<100%/green&normal
    OPTION=hilite_status: hitpoints/<66%/yellow&normal
    OPTION=hilite_status: hitpoints/<50%/orange&normal
    OPTION=hilite_status: hitpoints/<33%/red&bold
    OPTION=hilite_status: hitpoints/<15%/red&inverse
    OPTION=hilite_status: condition/major/orange&inverse
    OPTION=hilite_status: condition/lev+fly/red&inverse
\end{verbatim}

%.lp
%.hn 2
\subsection*{Modifier les symboles de \textit{NetHack}}

%.pg
\textit{NetHack} peut charger des jeux de symboles complets depuis le fichier des symboles.

%.pg
Les options qui servent à choisir un jeu de symboles particulier dans le
fichier des symboles sont :

\begin{description}[font=\mdseries\itshape]
%.lp
\item[symset]
Indique le nom du jeu de symboles que vous voulez charger.
\textit{symbols}.

%.lp
\item[roguesymset]
Indique le nom du jeu de symboles que vous voulez charger pour l'affichage
du niveau Rogue.
\end{description}

Vous pouvez aussi redéfinir un ou plusieurs symboles à l'aide des options
\textit{SYMBOLS} et \textit{ROGUESYMBOLS} du fichier de configuration.
Les symboles s'indiquent sous forme de paires \textit{nom:valeur}.
Notez que \textit{NetHack} interprète les séquences d'échappement de
la chaîne \textit{valeur} à la manière habituelle du C.
Cela signifie que `\texttt{\textbackslash}' est un préfixe qui fait prendre littéralement
le caractère suivant.
Ainsi, `\texttt{\textbackslash}' doit s'écrire `\texttt{\textbackslash\textbackslash}'.
La forme de préfixe spéciale
`\texttt{\textbackslash m}' active le bit meta dans la valeur du symbole, et le
préfixe `\texttt{\textasciicircum}' fait traiter le caractère suivant comme un caractère
de contrôle.

\small
\begin{longtable}{lll}
\caption[]{Symboles de NetHack}\\
Défaut                      & Nom du symbole                & Description\\
\hline \hline
\endhead
\space @ & S\textunderscore air                     &	(air)\\
\textunderscore  & S\textunderscore altar                   &	(autel)\\
" & S\textunderscore amulet                  &	(amulette)\\
A & S\textunderscore angel                   &	(être angélique)\\
a & S\textunderscore ant                     &	(fourmi ou autre insecte)\\
\textasciicircum & S\textunderscore anti\textunderscore magic\textunderscore trap       &	(champ antimagie)\\
{[} & S\textunderscore armor                   &	(armure ou pièce d'armure)\\
{[} & S\textunderscore armour                  &	(armure ou pièce d'armure)\\
\textasciicircum & S\textunderscore arrow\textunderscore trap             &	(piège à flèches)\\
0 & S\textunderscore ball                    &	(boulet de fer)\\
\# & S\textunderscore bars                    &	(barreaux de fer)\\
B & S\textunderscore bat                     &	(chauve-souris ou oiseau)\\
\textasciicircum & S\textunderscore bear\textunderscore trap              &	(piège à ours)\\
- & S\textunderscore blcorn                  &	(coin inférieur gauche)\\
b & S\textunderscore blob                    &	(blob)\\
+ & S\textunderscore book                    &	(livre de sorts)\\
) & S\textunderscore boomleft                &	(boomerang ouvert à gauche)\\
( & S\textunderscore boomright               &	(boomerang ouvert à droite)\\
\textasciigrave & S\textunderscore boulder                 &	(rocher)\\
- & S\textunderscore brcorn                  &	(coin inférieur droit)\\
> & S\textunderscore brdnladder              &	(échelle d'embranchement descendante)\\
> & S\textunderscore brdnstair               &	(escalier d'embranchement descendant)\\
< & S\textunderscore brupladder              &	(échelle d'embranchement montante)\\
< & S\textunderscore brupstair               &	(escalier d'embranchement montant)\\
C & S\textunderscore centaur                 &	(centaure)\\
\textunderscore & S\textunderscore chain                   &	(chaîne de fer)\\
\# & S\textunderscore cloud                   &	(nuage)\\
c & S\textunderscore cockatrice              &	(cocatrix)\\
\textdollar & S\textunderscore coin                    &	(tas de pièces)\\
\# & S\textunderscore corr                    &	(couloir)\\
- & S\textunderscore crwall                  &	(mur)\\
- & S\textunderscore darkroom                &	(pièce sombre)\\
\textasciicircum & S\textunderscore dart\textunderscore trap              &	(piège à fléchettes)\\
\& & S\textunderscore demon                   &	(démon majeur)\\
* & S\textunderscore digbeam                 &	(rayon de creusement)\\
> & S\textunderscore dnladder                &	(échelle descendante)\\
> & S\textunderscore dnstair                 &	(escalier descendant)\\
d & S\textunderscore dog                     &	(chien ou autre canidé)\\
D & S\textunderscore dragon                  &	(dragon)\\
; & S\textunderscore eel                     &	(monstre marin)\\
E & S\textunderscore elemental               &	(élémentaire)\\
\# & S\textunderscore engrcorr                      &	(gravure dans un couloir)\\
\textasciigrave & S\textunderscore engroom                 &	(gravure dans une pièce)\\
/ & S\textunderscore expl\textunderscore tl          &	(explosion en haut à gauche)\\
- & S\textunderscore expl\textunderscore tc          &	(explosion en haut au centre)\\
\textbackslash & S\textunderscore expl\textunderscore tr          &	(explosion en haut à droite)\\
| & S\textunderscore expl\textunderscore ml          &	(explosion au milieu à gauche)\\
~ & S\textunderscore expl\textunderscore mc          &	(explosion au milieu au centre)\\
| & S\textunderscore expl\textunderscore mr          &	(explosion au milieu à droite)\\
\textbackslash & S\textunderscore expl\textunderscore bl          &	(explosion en bas à gauche)\\
- & S\textunderscore expl\textunderscore bc          &	(explosion en bas au centre)\\
/ & S\textunderscore expl\textunderscore br          &	(explosion en bas à droite)\\
e & S\textunderscore eye                     &	(œil ou sphère)\\
\textasciicircum & S\textunderscore falling\textunderscore rock\textunderscore trap     &	(piège à chute de pierres)\\
f & S\textunderscore feline                  &	(chat ou autre félin)\\
\textasciicircum & S\textunderscore fire\textunderscore trap              &	(piège de feu)\\
! & S\textunderscore flashbeam               &	(rayon éclair)\\
\% & S\textunderscore food                    &	(morceau de nourriture)\\
\textbraceleft & S\textunderscore fountain                &	(fontaine)\\
F & S\textunderscore fungus                  &	(champignon ou moisissure)\\
* & S\textunderscore gem                     &	(gemme ou caillou)\\
~ & S\textunderscore ghost                   &	(fantôme)\\
H & S\textunderscore giant                   &	(humanoïde géant)\\
G & S\textunderscore gnome                   &	(gnome)\\
\textquotesingle & S\textunderscore golem                   &	(golem)\\
| & S\textunderscore grave                   &	(tombe)\\
g & S\textunderscore gremlin                 &	(gremlin)\\
- & S\textunderscore hbeam                   &	(mur)\\
\# & S\textunderscore hcdbridge               &	(pont-levis levé horizontal)\\
+ & S\textunderscore hcdoor                  &	(porte fermée)\\
. & S\textunderscore hodbridge               &	(pont-levis baissé horizontal)\\
| & S\textunderscore hodoor                  &	(porte ouverte)\\
\textasciicircum & S\textunderscore hole                    &	(trou)\\
@ & S\textunderscore human                   &	(humain ou elfe)\\
h & S\textunderscore humanoid                &	(humanoïde)\\
- & S\textunderscore hwall                   &	(mur horizontal)\\
. & S\textunderscore ice                     &	(glace)\\
i & S\textunderscore imp                     &	(diablotin ou démon mineur)\\
I & S\textunderscore invisible               &	(monstre invisible)\\
J & S\textunderscore jabberwock              &	(jabberwock)\\
j & S\textunderscore jelly                   &	(gelée)\\
k & S\textunderscore kobold                  &	(kobold)\\
K & S\textunderscore kop                     &	(Keystone Kop)\\
\textasciicircum & S\textunderscore land\textunderscore mine              &	(mine terrestre)\\
\textbraceright & S\textunderscore lava                    &	(lave en fusion)\\
\textbraceright & S\textunderscore lavawall                &	(mur de lave)\\
1 & S\textunderscore leprechaun              &	(leprechaun)\\
\textasciicircum & S\textunderscore level\textunderscore teleporter       &	(téléporteur de niveau)\\
L & S\textunderscore lich                    &	(liche)\\
y & S\textunderscore light                   &	(lumière)\\
\# & S\textunderscore litcorr                 &	(couloir éclairé)\\
: & S\textunderscore lizard                  &	(lézard)\\
\textbackslash & S\textunderscore lslant                  &	(mur)\\
\textasciicircum & S\textunderscore magic\textunderscore portal           &	(portail magique)\\
\textasciicircum & S\textunderscore magic\textunderscore trap             &	(piège magique)\\
m & S\textunderscore mimic                   &	(mimique)\\
{]} & S\textunderscore mimic\textunderscore def              &	(mimique)\\
M & S\textunderscore mummy                   &	(momie)\\
N & S\textunderscore naga                    &	(naga)\\
. & S\textunderscore ndoor                   &	(embrasure de porte)\\
n & S\textunderscore nymph                   &	(nymphe)\\
O & S\textunderscore ogre                    &	(ogre)\\
o & S\textunderscore orc                     &	(orque)\\
p & S\textunderscore piercer                 &	(perceur)\\
\textasciicircum & S\textunderscore pit                     &	(fosse)\\
\# & S\textunderscore poisoncloud             &	(nuage toxique)\\
\textasciicircum & S\textunderscore polymorph\textunderscore trap         &	(piège de métamorphose)\\
\textbraceright & S\textunderscore pool                    &	(eau)\\
! & S\textunderscore potion                  &	(potion)\\
P & S\textunderscore pudding                 &	(pouding ou limon)\\
q & S\textunderscore quadruped               &	(quadrupède)\\
Q & S\textunderscore quantmech               &	(mécanicien quantique)\\
= & S\textunderscore ring                    &	(anneau)\\
\textasciigrave & S\textunderscore rock                    &	(rocher ou statue)\\
r & S\textunderscore rodent                  &	(rongeur)\\
\textasciicircum & S\textunderscore rolling\textunderscore boulder\textunderscore trap  &	(piège à rocher roulant)\\
. & S\textunderscore room                    &	(sol d'une pièce)\\
/ & S\textunderscore rslant                  &	(mur)\\
\textasciicircum & S\textunderscore rust\textunderscore trap              &	(piège à rouille)\\
R & S\textunderscore rustmonst               &	(monstre de rouille ou désenchanteur)\\
? & S\textunderscore scroll                  &	(parchemin)\\
\# & S\textunderscore sink                    &	(évier)\\
\textasciicircum & S\textunderscore sleeping\textunderscore gas\textunderscore trap     &	(piège à gaz soporifique)\\
S & S\textunderscore snake                   &	(serpent)\\
s & S\textunderscore spider                  &	(arachnide ou mille-pattes)\\
\textasciicircum & S\textunderscore spiked\textunderscore pit             &	(fosse à pieux)\\
\textasciicircum & S\textunderscore squeaky\textunderscore board          &	(planche grinçante)\\
0 & S\textunderscore ss1                     &	(bouclier magique 1 sur 4)\\
\# & S\textunderscore ss2                     &	(bouclier magique 2 sur 4)\\
@ & S\textunderscore ss3                     &	(bouclier magique 3 sur 4)\\
* & S\textunderscore ss4                     &	(bouclier magique 4 sur 4)\\
\textasciicircum & S\textunderscore statue\textunderscore trap            &	(piège à statue)\\
~ & S\textunderscore stone                   &	(roche massive)\\
{]} & S\textunderscore strange\textunderscore obj      &	(objet étrange)\\
- & S\textunderscore sw\textunderscore bc                  &	(engloutissement en bas au centre)\\
\textbackslash & S\textunderscore sw\textunderscore bl                  &	(engloutissement en bas à gauche)\\
/ & S\textunderscore sw\textunderscore br                  &	(engloutissement en bas à droite	)\\
| & S\textunderscore sw\textunderscore ml                  &	(engloutissement au milieu à gauche)\\
| & S\textunderscore sw\textunderscore mr                  &	(engloutissement au milieu à droite)\\
- & S\textunderscore sw\textunderscore tc                  &	(engloutissement en haut au centre)\\
/ & S\textunderscore sw\textunderscore tl                  &	(engloutissement en haut à gauche)\\
\textbackslash & S\textunderscore sw\textunderscore tr                  &	(engloutissement en haut à droite)\\
- & S\textunderscore tdwall                  &	(mur)\\
\textasciicircum & S\textunderscore teleportation\textunderscore trap     &	(piège de téléportation)\\
\textbackslash & S\textunderscore throne                  &	(trône opulent)\\
- & S\textunderscore tlcorn                  &	(coin supérieur gauche)\\
| & S\textunderscore tlwall                  &	(mur)\\
( & S\textunderscore tool                    &	(objet utile (pioche, clé, lampe...))\\
\textasciicircum & S\textunderscore trap\textunderscore door              &	(trappe)\\
t & S\textunderscore trapper                 &	(piégeur ou rôdeur d'en haut)\\
- & S\textunderscore trcorn                  &	(coin supérieur droit)\\
\# & S\textunderscore tree                    &	(arbre)\\
T & S\textunderscore troll                   &	(troll)\\
| & S\textunderscore trwall                  &	(mur)\\
- & S\textunderscore tuwall                  &	(mur)\\
U & S\textunderscore umber                   &	(mastodonte des ombres)\\
~ & S\textunderscore unexplored              &	(terrain inexploré)\\
u & S\textunderscore unicorn                 &	(licorne ou cheval)\\
< & S\textunderscore upladder                &	(échelle montante)\\
< & S\textunderscore upstair                 &	(escalier montant)\\
V & S\textunderscore vampire                 &	(vampire)\\
| & S\textunderscore vbeam                   &	(mur)\\
\# & S\textunderscore vcdbridge               &	(pont-levis levé vertical)\\
+ & S\textunderscore vcdoor                  &	(porte fermée)\\
. & S\textunderscore venom                   &	(giclée de venin)\\
\textasciicircum & S\textunderscore vibrating\textunderscore square       &	(carré vibrant)\\
. & S\textunderscore vodbridge               &	(pont-levis baissé vertical)\\
- & S\textunderscore vodoor                  &	(porte ouverte)\\
v & S\textunderscore vortex                  &	(vortex)\\
| & S\textunderscore vwall                   &	(mur vertical)\\
/ & S\textunderscore wand                    &	(baguette)\\
\textbraceright & S\textunderscore water                   &	(eau)\\
\textbraceright & S\textunderscore weapon                  &	(arme)\\
" & S\textunderscore web                     &	(toile d'araignée)\\
w & S\textunderscore worm                    &	(ver)\\
\textasciitilde & S\textunderscore worm\textunderscore tail              &	(queue de ver long)\\
W & S\textunderscore wraith                  &	(âme en peine)\\
x & S\textunderscore xan                     &	(xan ou autre insecte extraordinaire)\\
X & S\textunderscore xorn                    &	(xorn)\\
Y & S\textunderscore yeti                    &	(créature simiesque)\\
Z & S\textunderscore zombie                  &	(zombie)\\
z & S\textunderscore zruty                   &	(zruty)\\
~ & S\textunderscore pet\textunderscore override     &	(tout familier si ACCESSIBILITY=1 est défini)\\
~ @ & S\textunderscore hero\textunderscore override    &	(le héros si ACCESSIBILITY=1 est défini)
\end{longtable}

\hyphenation{sysconf}	%no syllable breaks => don't hyphenate file name
%.lp
Remarques :

%.lp "*"
Plusieurs symboles de ce tableau semblent vides.
Il s'agit du caractère espace, sauf pour S\textunderscore pet\textunderscore override
et S\textunderscore hero\textunderscore override, qui n'ont pas de valeur par défaut
et ne peuvent être utilisés que s'ils sont autorisés dans le fichier ``sysconf''.

%.lp "*"
Le nom de S\textunderscore rock est trompeur : les cailloux et les pierres utilisent S\textunderscore gem.
Les ``rock'' dont il est question sont les statues et les rochers, mais depuis
la version 3.6.0, les statues sont affichées sous la forme du monstre qu'elles représentent.
S\textunderscore rock ne sert donc qu'aux rochers, et pas du tout s'il est
remplacé par le symbole plus spécifique S\textunderscore boulder.

%.lp
%.hn 2
\subsection*{Personnaliser la représentation des glyphes de la carte avec Unicode}

%.pg
Si votre plateforme ou votre terminal gère l'affichage des séquences de caractères
UTF-8, vous pouvez personnaliser l'affichage du jeu en attribuant aux représentations
des glyphes des points de code Unicode et des couleurs
rouge-vert-bleu. Ces personnalisations peuvent être indiquées pour un jeu de
symboles (symset) doté d'un gestionnaire UTF8 dans le fichier des symboles, comme le
jeu enhanced1, ou individuellement dans votre propre fichier nethack.rc.

Le format de définition d'une représentation de glyphe est :\\
\begin{verbatim}
OPTIONS=glyph:nom_du_glyphe/U+nnnn/R-V-B
\end{verbatim}

L'interface active doit savoir afficher les séquences de caractères UTF-8
et les couleurs rouge-vert-bleu explicites sur 24 bits pour que la représentation
du glyphe soit visible comme indiqué.

Par exemple, la ligne suivante dans votre fichier de configuration fera
utiliser à la représentation du glyphe G\textunderscore pool le point de code Unicode
U+224B et la couleur de valeur R-V-B 0-0-160 :\\
\begin{verbatim}
OPTIONS=glyph:G_pool/U+224B/0-0-160
\end{verbatim}

La liste des noms de glyphes acceptés peut être obtenue avec
\begin{verbatim}
    nethack --dumpglyphnames
\end{verbatim}
Chaque glyphe de NetHack peut être désigné avec le préfixe G\textunderscore ;
vous pouvez aussi utiliser un symbole S\textunderscore comme nom de glyphe et enregistrer la
représentation personnalisée pour tous les glyphes de NetHack qui correspondent
à ce symbole.

Vous devrez choisir un jeu de symboles doté d'un gestionnaire UTF8, comme le
jeu Enhanced, pour que les personnalisations soient affichées.

%.pg
%.hn 2
\subsection*{Configurer \textit{NetHack} pour les joueurs aveugles}

%.pg
\textit{NetHack} peut être configuré pour n'utiliser que des caractères ASCII standard pour
dessiner les cartes des donjons. Cela rend même les versions MS-DOS de \textit{NetHack}
(qui utilisent par défaut des caractères spéciaux de tracé de lignes) entièrement
accessibles aux aveugles qui utilisent la synthèse vocale et/ou le braille.
Les joueurs devront bien maîtriser les fonctions de revue de leur lecteur
d'écran et savoir naviguer horizontalement et
verticalement caractère par caractère. Les fonctions de recherche
de leur lecteur d'écran leur seront aussi très précieuses. Prenez soin
d'étudier ce manuel avant de jouer, afin d'avoir une idée de la disposition
de l'écran. Vous devrez aussi savoir localiser le curseur PC. Il se trouve
toujours là où est votre personnage. Chercher simplement un signe @
ne suffit pas toujours pour trouver votre personnage, car d'autres humanoïdes
sont représentés par le même signe. Votre lecteur d'écran devrait aussi avoir une
fonction indiquant la ligne et la colonne de votre curseur de revue et du curseur PC.
Ces coordonnées aident souvent les joueurs à mieux se représenter
la position générale des éléments à l'écran.
%.pg
\textit{NetHack} peut aussi être compilé avec la possibilité d'envoyer les messages
du jeu à un programme externe, par exemple un synthétiseur vocal. Si
la commande étendue ``\texttt{\#version}'' indique ``external program as a
message handler'', votre \textit{NetHack}
a été compilé avec cette possibilité. Si vous compilez \textit{NetHack}
à partir des sources
sous Linux ou un autre système POSIX, définissez \texttt{MSGHANDLER} pour l'activer.
Pour utiliser
cette possibilité, donnez à la variable d'environnement \texttt{NETHACK\textunderscore MSGHANDLER} le chemin
d'un exécutable, qui sera lancé avec le message du jeu comme
unique paramètre.
%.pg

Les réglages les plus importants pour rendre le jeu plus accessible sont :
%.pg
\begin{description}[font=\mdseries\itshape]
%.lp
\item[symset:plain]
Charge un jeu de symboles adapté aux joueurs aveugles.
%.lp
\item[menustyle:traditional]
Facilite l'interfaçage avec les synthétiseurs vocaux.
%.lp
\item[nomenu\textunderscore overlay]
Affiche les menus sur un écran effacé, alignés sur le bord gauche.
%.lp
\item[number\textunderscore pad]
Beaucoup de programmes d'accès vocal utilisent le pavé numérique pour parcourir l'écran.
Si c'est votre cas, désactivez l'option number\textunderscore pad et utilisez les
commandes traditionnelles à la Rogue.
%.lp
\item[paranoid\textunderscore confirmation:swim]
Empêche d'entrer dans l'eau ou la lave en marchant.
%.lp
\item[accessiblemsg]
Ajoute aux messages des informations de direction ou de position.
%.lp
\item[spot\textunderscore monsters]
Affiche un message lorsque le héros remarque un monstre ; à combiner avec accessiblemsg.
%.lp
\item[mon\textunderscore movement]
Affiche un message lorsque le héros remarque le déplacement d'un monstre ;
à combiner avec spot\textunderscore monsters et accessiblemsg.
%.lp
\item[autodescribe]
Décrit automatiquement le terrain sous le curseur lorsque vous visez.
%.lp
\item[mention\textunderscore map]
Affiche des messages lorsque des emplacements intéressants de la carte changent.
%.lp
\item[mention\textunderscore walls]
Affiche des messages lorsque vous marchez vers un mur ou lorsque la commande de
voyage a été interrompue.
%.lp
\item[whatis\textunderscore coord:compass]
Lorsque vous visez avec le curseur, décrit sa position par des coordonnées
relatives à votre personnage.
%.lp
\item[whatis\textunderscore filter:area]
Lorsque vous visez avec le curseur, filtre les emplacements possibles pour ne
retenir que ceux de la même zone (par ex. même pièce ou même couloir).
%.lp
\item[whatis\textunderscore moveskip]
Lorsque vous visez avec le curseur et utilisez le déplacement rapide, saute les
glyphes identiques au lieu d'avancer de 8 cases à la fois.
%.lp
\item[nostatus\textunderscore updates]
Empêche la mise à jour des lignes d'état en bas de l'écran, si
votre lecteur d'écran lit ces lignes. Les mêmes informations peuvent être
consultées avec la commande \texttt{\#attributes}.
%.lp
\item[showdamage]
Affiche un message indiquant les dégâts subis et les points de vie restants.
\end{description}

%.hn2
\subsection*{Configuration globale pour les administrateurs système}

%.pg
Si \textit{NetHack} est compilé avec l'option SYSCF, l'administrateur système
doit mettre en place une configuration globale : un fichier au
même format que le fichier de configuration traditionnel de chaque utilisateur (voir plus haut).

Ce fichier doit s'appeler sysconf et se trouver dans le même répertoire que
les autres fichiers de \textit{NetHack}.
Les options reconnues dans ce fichier sont listées ci-dessous. Toute option non
définie prend une valeur par défaut fixée à la compilation (qui n'est pas forcément
adaptée à votre système).

%.pg
\begin{description}[font=\mdseries\itshape]
%.lp
\item[WIZARDS]
Une liste, séparée par des espaces, des noms d'utilisateurs autorisés à
jouer en mode débogage (couramment appelé mode magicien, ou ``wizard mode'').
Une valeur constituée d'un seul
astérisque (*) permet à n'importe qui de lancer une partie en mode débogage.
%.lp
\item[SHELLERS]
Une liste des utilisateurs autorisés à utiliser la commande d'échappement vers le shell (`\texttt{!}').
La syntaxe est la même que pour WIZARDS.
%.lp
\item[EXPLORERS]
Une liste des utilisateurs autorisés à utiliser le mode exploration.
La syntaxe est la même que pour WIZARDS.
%.lp
\item[GENERICUSERS]
Une liste de noms d'utilisateurs pour lesquels le jeu demande ``Qui êtes-vous~?''
au lieu de les prendre comme nom du personnage.
%.lp
\item[MSGHANDLER]
Le chemin et le nom d'un exécutable.  Chaque fois qu'un message
est affiché dans la fenêtre des messages, \textit{NetHack} lance ce programme.  Le programme
reçoit le message comme unique paramètre.
%.lp
\item[MAXPLAYERS]
Limite le nombre maximum de parties pouvant se dérouler en même temps.
%.lp
\item[SUPPORT]
Une chaîne expliquant comment obtenir une assistance locale (pas de valeur par défaut).
%.lp
\item[RECOVER]
Une chaîne expliquant comment récupérer une partie sur ce système (pas de valeur par défaut).
%.lp
\item[SEDUCE]
0 ou 1 pour désactiver ou activer, respectivement, l'option SEDUCE.
Lorsqu'elle est désactivée, les incubes et les succubes se comportent comme des nymphes.
%.lp
\item[CHECK\textunderscore PLNAME]
Avec la valeur 1, EXPLORERS, WIZARDS et SHELLERS vérifient
le nom du joueur au lieu du nom de connexion de l'utilisateur.
%.lp
\item[CHECK\textunderscore SAVE\textunderscore UID]
0 ou 1 pour désactiver ou activer, respectivement, la vérification
de l'UID (numéro d'identification de l'utilisateur) des fichiers de sauvegarde (pour
vérifier que l'utilisateur qui restaure la partie est bien celui qui l'a sauvegardée).
\end{description}

%.pg
Les quatre options suivantes concernent le fichier des scores :
\begin{description}[font=\mdseries\itshape]
%.pg
%.lp
\item[PERSMAX]
Nombre maximum d'entrées pour une même personne.
%.lp
\item[ENTRYMAX]
Nombre maximum d'entrées dans le fichier des scores.
%.lp
\item[POINTSMIN]
Nombre minimum de points pour obtenir une entrée dans le fichier des scores.
%.lp
\item[PERS\textunderscore IS\textunderscore UID]
0 ou 1 pour utiliser, respectivement, les noms d'utilisateurs ou les identifiants
numériques afin de distinguer les personnes dans le fichier des scores.
%.lp
\item[HIDEUSAGE]
0 ou 1 pour afficher ou masquer l'entrée du menu d'aide consacrée
à l'utilisation en ligne de commande.
%.lp
\item[MAX\textunderscore STATUENAME\textunderscore RANK]
Nombre maximum d'entrées du fichier des scores utilisées pour
les noms aléatoires des statues (10 par défaut).
%.lp
\item[ACCESSIBILITY]
0 ou 1 pour désactiver ou activer, respectivement, la possibilité pour les joueurs
de définir les symboles S\textunderscore pet\textunderscore override et S\textunderscore hero\textunderscore override
dans leur fichier de configuration.
%.lp
\item[PORTABLE\textunderscore DEVICE\textunderscore PATHS]
0 ou 1, Windows uniquement : le jeu cherchera tous ses fichiers
externes, et écrira tous ses fichiers de sortie, en un seul endroit
plutôt qu'aux emplacements standard.
%.lp
\item[DUMPLOGFILE]
Le nom du fichier où est enregistré le journal de fin de partie (dumplog).
Si elle n'est pas définie, aucun dumplog n'est créé.
Disponible uniquement si votre jeu est compilé avec DUMPLOG.
Les marqueurs suivants sont acceptés :
% FIXME: this should be changed to a nested list or else be forcibly indented
\texttt{\%\%}  --- un `\texttt{\%}' littéral\\
\texttt{\%v}  --- version (par ex. ``\texttt{5.0.1-0}'')\\
\texttt{\%u}  --- UID de la partie\\
\texttt{\%t}  --- heure de début de la partie, au format horodatage UNIX\\
\texttt{\%T}  --- heure actuelle, au format horodatage UNIX\\
\texttt{\%d}  --- heure de début de la partie, au format AAAAMMJJhhmmss\\
\texttt{\%D}  --- heure actuelle, au format AAAAMMJJhhmmss\\
\texttt{\%n}  --- nom du joueur\\
\texttt{\%N}  --- premier caractère du nom du joueur
%.lp
\item[LIVELOG]
Un masque de bits des types d'événements à écrire dans
le fichier \textit{livelog} s'il existe.
Le fichier \textit{sysconf} d'exemple fourni avec le programme contient un
commentaire qui donne la signification des différents bits.
Destiné aux serveurs accueillant plusieurs joueurs simultanément
(chacun jouant, bien entendu, sa propre partie solo),
pour montrer leur progression à des observateurs.
Ne sert que si le programme a été compilé avec LIVELOG activé.
Lorsqu'elle est disponible, il vaut mieux la laisser en commentaire sur les
installations à un seul joueur, car avec le temps le fichier pourrait devenir
énorme s'il n'est pas activement entretenu.
%.lp
\item[CRASHREPORTURL]
Si elle vaut
\url{https://www.nethack.org/links/cr-5.0.0.html}
et que la prise en charge est compilée, ouvre une fenêtre de navigateur remplie avec
les informations nécessaires pour signaler un problème si le jeu panique ou se
retrouve dans un état interne incohérent, ou si la commande \#bugreport est
invoquée.
\end{description}

%.hn 1
\section{Score}

%.pg
\textit{NetHack} tient une liste des meilleurs scores ou des meilleurs joueurs de votre
machine, selon sa configuration.  Dans le second cas, chaque compte de
la machine ne peut inscrire qu'un seul score non victorieux dans cette liste.  Si
votre score dépasse celui de quelqu'un d'autre dans la liste, ou si vous battez votre
score précédent, vous serez inséré à la bonne place sous votre
nom actuel.  Le nombre de scores conservés peut aussi être réglé lors de
la compilation de \textit{NetHack}.

%.pg
Votre score dépend surtout de l'expérience que vous avez acquise, de la quantité
de butin que vous avez amassée, de la profondeur que vous avez explorée et de la
façon dont la partie s'est terminée.  Si vous abandonnez, vous repartez avec tout
votre or.
Si, en revanche, vous vous faites tuer dans les Mazes of Menace, la guilde
n'entendra parler que de 90\,\% de votre or lorsque votre cadavre sera découvert
(on a vu des aventuriers prélever une commission pour l'avoir trouvé).  Réfléchissez donc
bien : voulez-vous frapper une dernière fois ce monstre en espérant
survivre, ou abandonner et vous arrêter avec ce que vous avez~?  Si vous abandonnez,
vous gardez tout votre or ; mais si vous frappez et survivez, vous en trouverez peut-être davantage.

%.pg
Si vous voulez simplement voir la liste actuelle des meilleurs joueurs ou parties,
vous pouvez taper
\begin{verbatim}
    nethack -s all
\end{verbatim}
sur la plupart des versions.

%.hn 1
\section{Mode exploration}

%.pg
\textit{NetHack} est un jeu complexe et difficile.  Les novices pourraient
trembler de peur, conscients de leur ignorance des moyens de survivre.  Eh bien,
n'ayez crainte.  Votre donjon est équipé d'un mode ``exploration'' ou ``découverte''
qui vous permet de conserver d'anciennes sauvegardes et de tromper la mort, pour
le prix dérisoire de ne pas figurer dans la liste des meilleurs scores.

%.pg
Il y a deux façons d'activer le mode exploration.  La première consiste à lancer
le jeu avec l'option de ligne de commande \texttt{-X}
ou avec l'option
\textit{playmode:explore}.
La seconde consiste à utiliser la commande étendue `\texttt{\#exploremode}' en
cours de partie.  Commencer une nouvelle partie en mode exploration donne à votre
personnage une baguette de souhait dans son inventaire initial ; passer en ce
mode en cours de partie, non.  Les autres avantages du mode exploration sont laissés
à la découverte du lecteur intrépide.

%.pg
%.hn 2
\subsection*{Mode débogage}

%.pg
Le mode débogage, également appelé mode magicien (``wizard mode''), n'est pas
documenté en dehors de cette brève description et des diverses commandes
``mode débogage uniquement'' mentionnées parmi les descriptions des commandes.
Il est destiné à traquer les problèmes du
programme plutôt qu'à conférer des pouvoirs divins à votre personnage, et
les joueurs qui s'essaient au débogage sont censés découvrir par eux-mêmes
comment l'utiliser.
On l'active en lançant le jeu avec l'option de ligne de commande
\texttt{-D}
ou avec l'option
\textit{playmode:debug}.

%.pg
Sur certains systèmes, le joueur doit être connecté
sous un nom d'utilisateur particulier pour avoir le droit d'utiliser le mode débogage ;
sur d'autres, le héros doit porter un nom de personnage particulier (mais peut avoir
n'importe quel rôle ; il n'y a aucun lien entre le ``mode magicien'' et le rôle de \textit{Magicien}).
Tenter de lancer une partie en mode débogage sans en avoir le droit,
ou lorsqu'il n'est pas disponible, fait basculer en mode exploration à la place.

%.hn
\section{Crédits}
%.pg
Le jeu %
\textit{hack} original s'inspirait du jeu \textit{rogue} du système
%.ux
UNIX
de Berkeley.  De larges portions de ce document ont été effrontément
pompées sur %
\textit{A Guide to the Dungeons of Doom}, de Michael C. Toy
et Kenneth C. R. C. Arnold.  De petites portions ont été adaptées de
\textit{Further Exploration of the Dungeons of Doom}, de Ken Arromdee.

%.pg
\textit{NetHack} est le fruit du travail de littéralement des dizaines de personnes.
Les principales étapes du développement du jeu sont décrites ci-dessous :

%.pg
\bigskip
\noindent \textit{Jay Fenlason} a écrit le \textit{Hack} original, avec l'aide de {\it
Kenny Woodland}, \textit{Mike Thome} et \textit{Jon Payne}.

%.pg
\medskip
\noindent \textit{Andries Brouwer} en a fait une réécriture majeure lorsqu'il était au
Stichting Mathematisch Centrum (aujourd'hui Centrum Wiskunde \& Informatica),
transformant Hack en un jeu très différent.
Il a publié le code source de Hack pour les systèmes UNIX
en le diffusant sur le groupe de discussion Usenet
\textit{net.sources} (renommé plus tard \textit{comp.sources}),
avec la version 1.0 en décembre 1984, puis les versions 1.0.1, 1.0.2
et enfin 1.0.3 en juillet 1985.
Le groupe Usenet \textit{net.games.hack} (renommé
plus tard \textit{rec.games.hack}, et finalement remplacé
par \textit{rec.games.roguelike.nethack})
a été créé pour en discuter.

%.pg
\medskip
\noindent \textit{Don G. Kneller} a porté \textit{Hack} 1.0.3 vers Microsoft C et MS-DOS,
produisant \textit{PC Hack} 1.01e, a ajouté la prise en charge des graphismes DEC Rainbow dans
la version 1.03g, puis a produit au moins quatre autres versions (3.0, 3.2,
3.51 et 3.6 ;
notez qu'il s'agit d'anciens numéros de version de \textit{Hack}, et non de
\textit{NetHack}).

%.pg
\medskip
\noindent \textit{R. Black} a porté \textit{PC Hack} 3.51 vers Lattice C et l'Atari
520/1040ST, produisant \textit{ST Hack} 1.03.

%.pg
\medskip
\noindent \textit{Mike Stephenson} a fusionné ces différentes versions,
en incorporant nombre des fonctionnalités ajoutées, et a produit \textit{NetHack} version
1.4 en 1987.
Il a ensuite coordonné une distribution de milliers de personnes pour améliorer et
déboguer \textit{NetHack} 1.4, et a publié les versions 2.2 et 2.3 de \textit{NetHack}.
Comme Hack, elles ont été publiées en diffusant leur code source sur Usenet, où
elles sont restées disponibles dans diverses archives accessibles
par \textit{ftp} et \textit{uucp} après leur expiration du groupe de discussion.

%.pg
\medskip
\noindent Plus tard, Mike a coordonné une réécriture majeure du jeu, à la tête d'une équipe
comprenant \textit{Ken Arromdee}, \textit{Jean-Christophe Collet}, \textit{Steve Creps},
\textit{Eric Hendrickson}, \textit{Izchak Miller}, \textit{Eric S. Raymond}, \textit{John
Rupley}, \textit{Mike Threepoint} et \textit{Janet Walz}, pour produire
\textit{NetHack} 3.0c.

%.pg
\medskip
\noindent \textit{NetHack} 3.0 a été porté sur Atari par \textit{Eric R. Smith}, sur OS/2 par
\textit{Timo Hakulinen} et sur VMS par \textit{David Gentzel}.  Tous trois,
ainsi que \textit{Kevin Darcy}, ont ensuite rejoint l'équipe de développement principale de \textit{NetHack} (\textit{NetHack Development Team}) pour produire
les révisions suivantes de la 3.0.

%.pg
\medskip
\noindent \textit{Olaf Seibert} a porté \textit{NetHack} 2.3 et 3.0 sur Amiga.  {\it
Norm Meluch}, \textit{Stephen Spackman} et \textit{Pierre Martineau} ont conçu
le code de recouvrement (overlay) de \textit{NetHack PC} 3.0.  \textit{Johnny Lee} a porté
\textit{NetHack} 3.0 sur Macintosh.  Avec divers autres explorateurs de donjons, ils
ont continué à améliorer les versions PC, Macintosh et Amiga au fil des
révisions ultérieures de la 3.0.

%.pg
La version 3.0 a connu dix révisions de ``niveau de correctif'' (patch-level)
publiées à un rythme assez soutenu.
À l'époque, les versions s'appelaient 3.0 pour la version de base, puis
``3.0a'' à ``3.0j'',
``3.0~patchlevel~1'' à ``3.0~patchlevel~10'',
ou ``3.0pl1'' à ``3.0pl10''
plutôt que 3.0.0 et 3.0.1 à 3.0.10 ;
la numérotation à trois composantes a commencé à être utilisée avec la 3.1.0.

%.pg
\medskip
\noindent Dirigée par \textit{Mike Stephenson} et coordonnée par \textit{Izchak Miller}
et \textit{Janet Walz}, l'équipe de développement de \textit{NetHack}, qui comprenait désormais
\textit{Ken Arromdee},
\textit{David Cohrs}, \textit{Jean-Christophe Collet}, \textit{Kevin Darcy},
\textit{Matt Day}, \textit{Timo Hakulinen}, \textit{Steve Linhart}, \textit{Dean Luick},
\textit{Pat Rankin}, \textit{Eric Raymond} et \textit{Eric Smith}, a entrepris une
révision radicale de la 3.0.
Ils ont restructuré la conception du jeu et réécrit de grandes
parties du code.
Ils ont ajouté plusieurs donjons, un nouvel affichage, des quêtes
propres à chaque personnage, une nouvelle fin de partie et bien d'autres nouveautés, et
ont produit \textit{NetHack} 3.1.
La version 3.1.0 est sortie en janvier 1993.

%.pg
\medskip
\noindent \textit{Ken Lorber}, \textit{Gregg Wonderly} et \textit{Greg Olson}, avec l'aide
de \textit{Richard Addison}, \textit{Mike Passaretti} et \textit{Olaf Seibert},
ont développé \textit{NetHack} 3.1 pour l'Amiga.

%.pg
\medskip
\noindent \textit{Norm Meluch} et \textit{Kevin Smolkowski}, avec l'aide de
\textit{Carl Schelin}, \textit{Stephen Spackman}, \textit{Steve VanDevender}
et \textit{Paul Winner}, ont porté \textit{NetHack} 3.1 sur PC.

%.pg
\medskip
\noindent \textit{Jon W\textbraceleft tte} et \textit{Hao-yang Wang},
avec l'aide de \textit{Ross Brown}, \textit{Mike Engber}, \textit{David Hairston},
\textit{Michael Hamel}, \textit{Jonathan Handler}, \textit{Johnny Lee},
\textit{Tim Lennan}, \textit{Rob Menke} et \textit{Andy Swanson},
ont développé \textit{NetHack} 3.1 pour le Macintosh, en le portant pour MPW.
En s'appuyant sur leur travail, \textit{Bart House} a ajouté un portage pour Think C.

%.pg
\medskip
\noindent \textit{Timo Hakulinen} a porté \textit{NetHack} 3.1 sur OS/2.
\textit{Eric Smith} a porté \textit{NetHack} 3.1 sur Atari.
\textit{Pat Rankin}, avec l'aide de \textit{Joshua Delahunty},
a été responsable de la version VMS de \textit{NetHack} 3.1.
\textit{Michael Allison} a porté \textit{NetHack} 3.1 sur Windows NT.

%.pg
\medskip
\noindent \textit{Dean Luick}, avec l'aide de \textit{David Cohrs}, a développé
\textit{NetHack} 3.1 pour X11.
Cette version dessinait la carte en mode texte et non graphique, mais
incluait \texttt{nh10.bdf}, une police X11 personnalisée facultative contenant
de minuscules images à la place des lettres et de la ponctuation, précurseur des tuiles.
Ces images ne distinguaient pas chaque type de monstre et d'objet, et se contentaient
de remplacer les classes de monstres et d'objets (une image pour tous les
insectes ``\texttt{a}'', une autre pour toutes les armures ``\texttt{[}'' et ainsi
de suite, sans images distinctes pour les scarabées et les fourmis, ni pour les capes et les bottes).

%.pg
\medskip
\noindent \textit{Warwick Allison} a écrit pour l'Atari une version graphique
de \textit{NetHack}
dans laquelle les petites images étaient appelées ``icônes'' et
étaient propres à chaque type de monstre et d'objet, et non plus seulement
à leur classe.
Il les a offertes à l'équipe de développement de \textit{NetHack} (\textit{NetHack Development Team}), qui les a
rebaptisées ``tuiles'' (tiles), un usage original repris depuis
par divers autres jeux.
La prise en charge des tuiles de \textit{NetHack} a ensuite été mise en œuvre sur d'autres plateformes
(d'abord MS-DOS, puis Windows, Qt et X11).

%.pg
\medskip
\noindent L'équipe de développement de \textit{NetHack} 3.2, composée de \textit{Michael Allison}, \textit{Ken
Arromdee}, \textit{David Cohrs}, \textit{Jessie Collet}, \textit{Steve Creps}, {\it
Kevin Darcy}, \textit{Timo Hakulinen}, \textit{Steve Linhart}, \textit{Dean Luick},
\textit{Pat Rankin}, \textit{Eric Smith}, \textit{Mike Stephenson}, \textit{Janet Walz}
et \textit{Paul Winner}, a publié la version 3.2.0 en avril 1996.

%.pg
\medskip
\noindent La version 3.2 a marqué le dixième anniversaire de la formation de
l'équipe de développement.
Témoignage de leur dévouement au jeu, les treize membres
de l'équipe de développement de \textit{NetHack} d'origine faisaient toujours partie de l'équipe au
début des travaux sur cette version.
Entre la sortie de la 3.1.3 et celle de la 3.2.0,
l'un des membres fondateurs de l'équipe de développement de \textit{NetHack},
\textit{Dr. Izchak Miller}, a appris qu'il souffrait d'un cancer et nous a quittés.
Cette version du jeu lui a été
dédiée par les équipes de développement et de portage.

%.pg
La version 3.2 s'est révélée plus stable que les précédentes.
De nombreux bogues ont été corrigés, des abus éliminés, et des fonctionnalités ajustées
pour un meilleur équilibre de jeu.

%.pg
\medskip
Pendant la durée de vie de \textit{NetHack} 3.1 et 3.2, plusieurs passionnés
du jeu y ont ajouté
leurs propres modifications et ont rendu ces ``variantes'' publiquement
disponibles :

%.pg
\medskip
\textit{Tom Proudfoot} et \textit{Yuval Oren} ont créé \textit{NetHack++},
rapidement renommé \textit{NetHack$--$}
lorsque certains ont supposé à tort qu'il s'agissait d'une conversion du
code source \textit{C} en \textit{C++}.
De son côté, \textit{Stephen White} a écrit \textit{NetHack Plus}.
\textit{Tom Proudfoot} a plus tard fusionné \textit{NetHack Plus}
et son propre \textit{NetHack$--$} pour produire \textit{SLASH}.
\textit{Larry Stewart-Zerba} et \textit{Warwick Allison} ont amélioré le système
de lancement de sorts avec le Wizard Patch.
\textit{Warwick Allison} a aussi porté \textit{NetHack} vers l'interface Qt.

%.pg
\medskip
\textit{Warren Cheung} a combiné \textit{SLASH} avec le Wizard Patch
pour produire \textit{Slash'EM}, et,
avec l'aide de \textit{Kevin Hugo}, y a ajouté d'autres fonctionnalités.
Kevin a ensuite rejoint l'équipe de développement de \textit{NetHack} et a intégré
les meilleures de ces idées dans \textit{NetHack} 3.3.

%.pg
\medskip
La dernière mise à jour de la 3.2 a été la version corrective 3.2.3, publiée
en même temps que la 3.3.0 en décembre 1999, juste à temps pour l'an 2000.
En raison de la nouvelle version, la 3.2.3 n'a été publiée que sous forme de correctif
du code source, sans distribution prête à jouer pour les systèmes qui en avaient d'habitude.

%.pg
(Pour qui envisagerait de ressusciter une ancienne version : toutes les versions antérieures
à la 3.2.3 avaient un bogue de l'\textit{an 2000}.
Le fichier des meilleurs scores et le fichier journal contenaient
des dates dont l'année était écrite sur deux chiffres, et l'année 99 de 1999 était
suivie de l'année 100 de 2000.
Cela s'écrivait sans problème, mais
ajoutait involontairement une colonne supplémentaire dans le format du fichier, ce qui
empêchait de relire correctement les entrées de score, perturbant l'insertion
de nouveaux meilleurs scores et la récupération des noms d'anciens personnages utilisés
pour nommer au hasard les fantômes et les statues de la partie en cours.)

%.pg
\medskip
L'équipe de développement de \textit{NetHack} 3.3, composée de \textit{Michael Allison}, \textit{Ken Arromdee},
\textit{David Cohrs}, \textit{Jessie Collet}, \textit{Steve Creps}, \textit{Kevin Darcy},
\textit{Timo Hakulinen}, \textit{Kevin Hugo}, \textit{Steve Linhart}, \textit{Ken Lorber},
\textit{Dean Luick}, \textit{Pat Rankin}, \textit{Eric Smith}, \textit{Mike Stephenson},
\textit{Janet Walz} et \textit{Paul Winner}, a publié la 3.3.0 en
décembre 1999 et la 3.3.1 en août 2000.

%.pg
\medskip
La version 3.3 a apporté de nombreuses premières. Ce fut la première version à séparer
la race et la profession. La classe Elfe a été supprimée au profit d'une race elfe,
et les races naines, gnomes et orques ont fait leur première apparition dans
le jeu aux côtés de la familière race humaine.  Les rôles de Moine et de Rôdeur ont rejoint
les Archéologues, Barbares, Hommes des cavernes, Guérisseurs, Chevaliers, Prêtres, Voleurs, Samouraïs,
Touristes, Valkyries et, bien sûr, Magiciens.  Ce fut aussi la première version
à vous permettre de chevaucher une monture, et la première à disposer d'un site web
public recensant tous les bogues découverts.  Malgré
cette liste de bogues sans cesse croissante, la 3.3 s'est révélée assez stable pour durer
plus d'un an et demi.

%.pg
\medskip
L'équipe de développement de \textit{NetHack} 3.4 se composait initialement de
\textit{Michael Allison}, \textit{Ken Arromdee},
\textit{David Cohrs}, \textit{Jessie Collet}, \textit{Kevin Hugo}, \textit{Ken Lorber},
\textit{Dean Luick}, \textit{Pat Rankin}, \textit{Mike Stephenson},
\textit{Janet Walz} et \textit{Paul Winner}, rejoints par \textit{ Warwick Allison}
juste avant la sortie de \textit{NetHack} 3.4.0 en mars 2002.

%.pg
\medskip
Comme pour la version 3.3, diverses personnes ont contribué au jeu dans son ensemble,
ainsi qu'aux portages sur les différentes plateformes sur lesquelles \textit{NetHack}
fonctionne :

%.pg
\medskip
\noindent\textit{Pat Rankin} a maintenu la 3.4 pour VMS.

%.pg
\medskip
\noindent \textit{Michael Allison} a maintenu \textit{NetHack} 3.4 pour la plateforme
MS-DOS.
\textit{Paul Winner} et \textit{Yitzhak Sapir} l'ont encouragé.

%.pg
\medskip
\noindent \textit{Dean Luick}, \textit{Mark Modrall} et \textit{Kevin Hugo} ont maintenu et
amélioré le portage Macintosh de la 3.4.

%.pg
\medskip
\noindent \textit{Michael Allison}, \textit{David Cohrs}, \textit{Alex Kompel},
\textit{Dion Nicolaas} et
\textit{Yitzhak Sapir} ont maintenu et amélioré la 3.4 pour la plateforme
Microsoft Windows.
\textit{Alex Kompel} a fourni une nouvelle interface graphique pour le portage Windows.
\textit{Alex Kompel} a aussi fourni un portage Windows CE pour la 3.4.1.

%.pg
\medskip
\noindent \textit{Ron Van Iwaarden} a été le seul mainteneur de \textit{NetHack} pour
OS/2 pendant
plusieurs versions. Malheureusement, la dernière machine OS/2 de Ron a cessé de fonctionner
début 2006. Un grand merci à Ron pour avoir maintenu \textit{NetHack} en vie sur
OS/2 pendant toutes ces années.

%.pg
\medskip
\noindent \textit{Janne Salmij\"{a}rvi} et \textit{Teemu Suikki} ont maintenu
et amélioré le portage Amiga de la 3.4 après que \textit{Janne Salmij\"{a}rvi} l'eut ressuscité
pour la 3.3.1.

%.pg
\medskip
\noindent \textit{Christian ``Marvin'' Bressler} a maintenu la 3.4 pour l'Atari après l'avoir
ressuscitée pour la 3.3.1.

%.pg
\medskip
La sortie de \textit{NetHack} 3.4.3 en décembre 2003 a marqué le début d'une
longue pause dans les publications. La 3.4.3 s'est révélée remarquablement stable et
a continué à faire le bonheur de la communauté pendant plus de dix ans.
L'équipe de développement de \textit{NetHack} a continué à travailler lentement et discrètement sur le jeu en coulisses
pendant le règne de la 3.4.3. C'est pendant cette même période que plusieurs
nouvelles variantes sont apparues dans la communauté \textit{NetHack}. Notamment sporkhack de
Derek S. Ray, \textit{unnethack} de Patric Mueller, \textit{nitrohack} et ses
successeurs, d'abord de Daniel Thaler puis d'Alex Smith, et
\textit{Dynahack} de Tung Nguyen.
Certaines de ces variantes sont toujours
développées, maintenues et appréciées par la communauté aujourd'hui.

%.pg
\medskip
En septembre 2014, un instantané intermédiaire du code en cours de développement a été
rendu public par des tiers.
Comme ce code était un travail en cours
et n'avait pas été débogué au point de constituer une version publiable,
il a été décidé que les numéros de version présents dans cet instantané seraient
retirés et jamais utilisés pour une version officielle de \textit{NetHack}.
Une annonce a été publiée en ce sens sur le site web officiel de l'équipe de développement de \textit{NetHack},
\textit{nethack.org},
précisant qu'il n'y aurait jamais de version officielle 3.4.4, 3.5 ou
3.5.0.

%.pg
\medskip
En janvier 2015 a commencé la préparation de la sortie de \textit{NetHack} 3.6.

%.pg
\medskip
Au début du développement de ce qui allait devenir
la 3.6.0, l'équipe de développement de \textit{NetHack} se composait de \textit{Warwick Allison},
\textit{Michael Allison}, \textit{Ken Arromdee},
\textit{David Cohrs}, \textit{Jessie Collet},
\textit{Ken Lorber}, \textit{Dean Luick}, \textit{Pat Rankin},
\textit{Mike Stephenson}, \textit{Janet Walz} et \textit{Paul Winner}.
Début 2015, avant la sortie de la 3.6.0, de nouveaux membres,
\textit{Sean Hunt}, \textit{Pasi Kallinen} et \textit{Derek S. Ray},
ont rejoint l'équipe de développement de \textit{NetHack}.

%.pg
\medskip
Vers la fin du développement de la 3.6.0, l'une des principales sources d'inspiration
de nombre des aspects drôles et amusants du jeu,
l'auteur Terry Pratchett, nous a quittés. \textit{NetHack} 3.6.0 lui rend
hommage.

%.pg
\medskip
La 3.6.0 est sortie en décembre 2015 ; elle réunissait le travail accompli par l'équipe
de développement depuis la sortie de la 3.4.3 et certains des correctifs
de la communauté les plus appréciés.  De nombreux bogues ont été corrigés et une partie du code a été restructurée.

%.pg
\medskip
L'équipe de développement de \textit{NetHack}, ainsi que \textit{Steve VanDevender} et
\textit{Kevin Smolkowski}, ont veillé à ce que \textit{NetHack} 3.6 continue de
fonctionner sur diverses variantes du système UNIX et ont maintenu l'interface X11.

%.pg
\medskip
\textit{Ken Lorber}, \textit{Haoyang Wang}, \textit{Pat Rankin} et \textit{Dean Luick}
ont maintenu le portage de \textit{NetHack} 3.6 pour MacOS.

%.pg
\medskip
\textit{Michael Allison}, \textit{David Cohrs}, \textit{Bart House},
\textit{Pasi Kallinen}, \textit{Alex Kompel}, \textit{Dion Nicolaas},
\textit{Derek S. Ray} et  \textit{Yitzhak Sapir}
ont maintenu le portage de  \textit{NetHack} 3.6 pour Microsoft Windows.

%.pg
\medskip
\textit{Pat Rankin} a tenté de maintenir le portage VMS de \textit{NetHack} 3.6 en état de marche,
gêné par un accès limité.  \textit{Kevin Smolkowski} l'a mis à jour et testé
pour la version la plus récente d'OpenVMS (V8.4 au moment où nous écrivons) sur Alpha
et Integrity (alias Itanium, alias IA64), mais pas sur VAX.

%.pg
\medskip
\textit{Ray Chason}  a ressuscité le portage MS-DOS pour la 3.6 et a fourni les
mises à jour nécessaires à l'ensemble de la communauté.

%.pg
\medskip
Fin avril 2018, plusieurs centaines de corrections de bogues de la 3.6.0 et quelques nouvelles fonctionnalités
ont été rassemblées et publiées sous le nom de \textit{NetHack} 3.6.1.
L'équipe de développement de \textit{NetHack}, au
moment de la sortie de la 3.6.1, se composait de
\textit{Warwick Allison}, \textit{Michael Allison}, \textit{Ken Arromdee},
\textit{David Cohrs}, \textit{Jessie Collet},
\textit{Pasi Kallinen}, \textit{Ken Lorber}, \textit{Dean Luick},
\textit{Patric Mueller}, \textit{Pat Rankin}, \textit{Derek S. Ray},
\textit{Alex Smith}, \textit{Mike Stephenson}, \textit{Janet Walz} et
\textit{Paul Winner}.

%.pg
\medskip
Début mai 2019, 320 autres corrections de bogues, accompagnées de quelques améliorations et
de l'interface curses adoptée, ont été publiées sous le nom de 3.6.2.

%.pg
\medskip
\textit{Bart House}, qui contribuait au jeu au sein de l'équipe de portage
depuis des décennies, a rejoint l'équipe de développement de \textit{NetHack} fin mai 2019.

%.pg
\medskip
\textit{NetHack} 3.6.3 est sorti le 5 décembre 2019 avec plus de 190 corrections
de bogues de \textit{NetHack} 3.6.2.

%.pg
\medskip
\textit{NetHack} 3.6.4 est sorti le 18 décembre 2019 avec une correction de sécurité et
quelques corrections de bogues.

%.pg
\medskip
\textit{NetHack} 3.6.5 est sorti le 27 janvier 2020 avec quelques corrections de sécurité
et un petit nombre de corrections de bogues.

%.pg
\medskip
\textit{NetHack} 3.6.6 est sorti le 8 mars 2020 avec une correction de sécurité et
quelques corrections de bogues.

%.pg
\medskip
\textit{NetHack} 3.6.7 est sorti le 16 février 2023 avec une correction de sécurité et
quelques corrections de bogues.

%.pg
\medskip
Le développement de la version majeure destinée à succéder à \textit{NetHack} 3.6 a commencé en 2015,
à peu près au moment de la sortie de \textit{NetHack} 3.6.0. Ce développement
s'est poursuivi en parallèle de chacune des versions de \textit{NetHack} 3.6, de 2015
à 2023, et a continué jusqu'à fin avril 2026.  Pour la première fois,
ce développement a été partagé publiquement sur GitHub et SourceForge au fur et à mesure.
Il s'est fait sous l'étiquette NetHack-3.7 work-in-progress (WIP, travail en cours), bien que
le numéro de la prochaine version n'ait pas encore été arrêté.

%.pg
\medskip
L'ouverture du développement au public a apporté beaucoup de bonnes choses, et quelques
difficultés. Chacun a pu observer et critiquer les changements et les nouveautés
presque immédiatement, et ne s'en est souvent pas privé. Le système de pull requests
de GitHub a permis à chacun de contribuer facilement et directement au développement.
Les contributions ont résolu de très, très nombreux bogues du jeu, et nous remercions tous les
contributeurs.

%.pg
\medskip
Début 2026, le développement du jeu étant devenu assez stable pour envisager
une version officielle, l'équipe de développement a examiné la nature et le nombre
des changements apportés au jeu. Il était clair qu'ils avaient assez de profondeur et
d'ampleur pour justifier une version majeure, et le choix s'est porté sur la version 5.0.
Il s'agit d'une nouvelle version majeure succédant aux 3.x, sans l'ambiguïté
ni la confusion avec des variantes existantes qu'aurait pu entraîner
une sortie sous le numéro 4.0.

%.pg
\medskip
\textit{NetHack} 5.0.0 est sorti le 2 mai 2026.

%.pg
\medskip
Le code source de \textit{NetHack} 5.0.0 a été modifié et modernisé pour être
conforme à la norme C99. La version 5.0.0 contenait plus de 3100 corrections,
modifications et fonctionnalités mises à jour.

%.pg
\medskip
\textit{NetHack} 5.0 a été la première version à remplacer les compilateurs de niveaux
et de donjons basés sur lex et yacc des versions passées par une nouvelle approche
reposant sur un interpréteur Lua. Lua est aussi utilisé pour les textes des quêtes
dans \textit{NetHack} 5.0.0. Toute l'équipe de développement salue le travail
accompli par \textit{Pasi Kallinen} pour y parvenir.

%.pg
\medskip
Au moment de la sortie de \textit{NetHack} 5.0, l'équipe principale
de développement de \textit{NetHack},
membres actifs et anciens, comprenait \textit{Warwick Allison}, \textit{Michael Allison},
\textit{Ken Arromdee}, \textit{David Cohrs}, \textit{Jessie Collet}, \textit{Bart House},
\textit{Kevin Hugo}, \textit{Pasi Kallinen}, \textit{Ken Lorber}, \textit{Dean Luick},
\textit{Pat Rankin}, \textit{Derek S. Ray}, \textit{Alex Smith}, \textit{Patric Mueller},
\textit{Mike Stephenson}, \textit{Janet Walz}, \textit{Paul Winner}.

%.pg
\medskip
\textit{Ken Lorber}, \textit{Pat Rankin}, \textit{Patrick Mueller} et
\textit{Michael Allison} ont contribué à ce que \textit{NetHack} 5.0 fonctionne
sous macOS.

%.pg
\medskip
\textit{Ingo Paschke} a réussi, on ne sait trop comment, à faire revivre un portage de \textit{NetHack} 5.0
pour l'Amiga, en utilisant pour cela un compilateur croisé sur une plateforme moderne. Il a partagé
son travail, si bien que d'autres peuvent facilement produire \textit{NetHack} 5.0 pour l'Amiga grâce à
ses efforts.

%.pg
\medskip
\textit{Ray Chason} a assuré l'essentiel de la maintenance du
portage MS-DOS de \textit{NetHack} 5.0, y compris le portage de l'interface curses.
\textit{Michael Allison} a veillé à ce que les modifications du cœur de \textit{NetHack} 5.0
continuent de fonctionner avec le portage msdos, pour le maintenir en vie.  La compilation croisée
du portage MS-DOS a contribué à rendre cela possible, et presque indolore.

%.pg
\medskip
Parmi les personnes qui ont contribué au portage Windows de \textit{NetHack} 5.0 depuis le
début de son développement, il y a plus de onze ans, figurent
\textit{Michael Allison}, \textit{David Cohrs}, \textit{Bart House},
\textit{Pasi Kallinen}, \textit{Alex Kompel}, \textit{Dion Nicolaas},
\textit{Derek S. Ray} et \textit{Yitzhak Sapir}.

%.pg
\medskip
C'est avec tristesse que l'équipe de développement souhaite saluer et rappeler les
contributions passées de feu \textit{Ron Van Iwaarden}, qui fut le seul mainteneur
de \textit{NetHack} pour OS/2 pendant plusieurs versions de \textit{NetHack}. Ron nous manquera.

%.pg
\medskip
\noindent Le site web officiel de \textit{NetHack} est maintenu par \textit{Ken Lorber} à l'adresse
\url{https://www.nethack.org}.

%.pg
%.hn 2

\subsection*{Remerciements particuliers}
\noindent Au nom de la communauté \textit{NetHack}, merci encore une fois
à \textit{M. Drew Streib} et \textit{Pasi Kallinen} de fournir un
serveur \textit{NetHack} public à nethack.alt.org. Merci à \textit{Keith Simpson}
et \textit{Andy Thomson} pour hardfought.org. Merci à tous ces
explorateurs de donjons anonymes qui consacrent temps et efforts aux tournois
annuels de \textit{NetHack} tels que \textit{Junethack},
\textit{The November NetHack Tournament} et, autrefois,
\textit{devnull.net} (disparu pour l'instant, mais pas oublié).
\clearpage

%.hn 2
\subsection*{Explorateurs de donjons}
%.pg
\noindent De temps à autre, un individu dépravé quelque part sur le réseau envoie une
modification particulièrement intrigante pour améliorer le jeu.  L'équipe de
développement de \textit{NetHack} prend parfois note des noms des pires
de ces mécréants dans la présente liste des explorateurs de donjons :
\medskip
%.sd
\begin{center}
\begin{tabular}{llll}
%TABLE_START
Adam Aronow & Irina Rempt-Drijfhout & Mike Gallop\\
Alex Kompel & Izchak Miller & Mike Passaretti\\
Alex Smith & J. Ali Harlow & Mike Stephenson\\
Andreas Dorn & Janet Walz & Mikko Juola\\
Andy Church & Janne Salmij\"{a}rvi & Nathan Eady\\
Andy Swanson & Jean-Christophe Collet & Norm Meluch\\
Andy Thomson & Jeff Bailey & Olaf Seibert\\
Ari Huttunen & Jochen Erwied & Pasi Kallinen\\
Bart House & John Kallen & Pat Rankin\\
Benson I. Margulies & John Rupley & Patric Mueller\\
Bill Dyer & John S. Bien & Paul Winner\\
Boudewijn Waijers & Johnny Lee & Pierre Martineau\\
Bruce Cox & Jon W\{tte & Ralf Brown\\
Bruce Holloway & Jonathan Handler & Ray Chason\\
Bruce Mewborne & Joshua Delahunty & Richard Addison\\
Cameron Root & Karl Garrison & Richard Beigel\\
Carl Schelin & Keizo Yamamoto & Richard P. Hughey\\
Chris Russo & Keith Simpson & Rob Menke\\
David Cohrs & Ken Arnold & Robin Bandy\\
David Damerell & Ken Arromdee & Robin Johnson\\
David Gentzel & Ken Lorber & Roderick Schertler\\
David Hairston & Ken Washikita & Roland McGrath\\
Dean Luick & Kestrel Gregorich-Trevor & Ron Van Iwaarden\\
Del Lamb & Kevin Darcy & Ronnen Miller\\
Derek S. Ray & Kevin Hugo & Ross Brown\\
Deron Meranda & Kevin Sitze & Sascha Wostmann\\
Dion Nicolaas & Kevin Smolkowski & Scott Bigham\\
Dylan O'Donnell & Kevin Sweet & Scott R. Turner\\
Eric Backus & Lars Huttar & Sean Hunt\\
Eric Hendrickson & Leon Arnott & Stephen Spackman\\
Eric R. Smith & M. Drew Streib & Stefan Thielscher\\
Eric S. Raymond & Malcolm Ryan & Stephen White\\
Erik Andersen & Mark Gooderum & Steve Creps\\
Fredrik Ljungdahl & Mark Modrall & Steve Linhart\\
Frederick Roeber & Marvin Bressler & Steve VanDevender\\
G. Branden Robinson & Matthew Day & Teemu Suikki\\
Gil Neiger & Merlyn LeRoy & Tim Lennan\\
Greg Laskin & Michael Allison & Timo Hakulinen\\
Greg Olson & Michael Feir & Tom Almy\\
Gregg Wonderly & Michael Hamel & Tom West\\
Hao-yang Wang & Michael Meyer & Warren Cheung\\
Helge Hafting & Michael Sokolov & Warwick Allison\\
Ingo Paschke & Mike Engber & Yitzhak Sapir
%TABLE_END  Do not delete this line.
\end{tabular}
\end{center}
%.ed
\clearpage
%\vfill
%\begin{flushleft}
%\small
%Microsoft and MS-DOS are registered trademarks of Microsoft Corporation.\\
%%%Don't need next line if a UNIX macro automatically inserts footnotes.
%UNIX is a registered trademark of AT\&T.\\
%Lattice is a trademark of Lattice, Inc.\\
%Atari and 1040ST are trademarks of Atari, Inc.\\
%Amiga is a trademark of Commodore-Amiga, Inc.\\
%%.sm
%Brand and product names are trademarks or registered trademarks
%of their respective holders.
%\end{flushleft}

\end{document}
